% concatenated from moduli-v1.tex
%%%%%%%%%%%%%%%%%%%%%%%%%%%%%%%%%%%%%%%%%%%%%%%%%%%%%%%%%%%%%
\documentclass[12pt,reqno]{amsart}
\setlength{\textheight}{23cm}
\setlength{\textwidth}{16cm}
\setlength{\topmargin}{-0.8cm}
%\setlength{\parskip}{0.5\baselineskip}
\hoffset=-1.4cm 

\usepackage[skip=5pt plus1pt, indent=20pt]{parskip}

%%%%%%%%%%%%%%%%%%%%%%%%%%%%%%%%%%%%%%%%%%%%%%%%%%%%%%%%%%%%%
%\usepackage{stix}
\usepackage{palatino}
\usepackage{mathpazo}
\usepackage{stmaryrd}
\usepackage{tikz-cd}
\usepackage{textcomp}
%\usepackage[T1]{fontenc}
%\usepackage{lmodern}
%\usepackage{newpxmath}
\usepackage{amsmath,amssymb,amsxtra,color,calligra,mathrsfs,tcolorbox,stmaryrd,amscd}

%%%%%%%%%%%%%%%%%%%%%%%%%%%%%%%%%%%%%%
%\usepackage{tocloft}




%%%%%%%%%%%%%%%%%%%%%%%%%%%%%%%%%%%%%%%%%%%

\usepackage{xcolor} % Required for specifying custom colours
\colorlet{mdtRed}{red!50!black}
\colorlet{dblue}{blue!50!black}
%\definecolor{dblue}{rgb}{0,0,.6}
\usepackage[colorlinks,pagebackref=true]{hyperref}
\hypersetup{linkcolor=dblue,citecolor=dblue,filecolor=dullmagenta,urlcolor=mdtRed}
\renewcommand*{\backref}[1]{}
\renewcommand*{\backrefalt}[4]{[{%
		\ifcase #1 Not cited.%
		\or $\uparrow$~#2.%
		\else $\uparrow$~#2.%
		\fi%
	}]}
\usepackage[all]{xy}
\usepackage{tikz,tikz-cd,tkz-graph,enumerate}
\usetikzlibrary{matrix,arrows,decorations.pathmorphing}

\usepackage[us,12hr]{datetime}
\usepackage{comment}

%%%%%%%%%%%%%%%%%%%%%%%%%%%%%%%%%%%%%%%%%%%%%%%%%%%%%%%%%%%%%
\DeclareMathOperator{\Pic}{\textnormal{Pic}}
\DeclareMathOperator{\codim}{\textnormal{codim}}
\DeclareMathOperator{\Spec}{{\rm Spec}}
\DeclareMathOperator{\Id}{{\rm Id}}
\DeclareMathOperator{\Hom}{{\rm Hom}}
\DeclareMathOperator{\Aut}{{\rm Aut}}
\DeclareMathOperator{\End}{{\rm End}}
\DeclareMathOperator{\Sch}{{\rm Sch}}
\DeclareMathOperator{\op}{{\rm op}}
\DeclareMathOperator{\Vect}{{\bf Vect}}
\DeclareMathOperator{\PVect}{{\bf PVect}}
\DeclareMathOperator{\RPVect}{{\bf RealParVect}}
\DeclareMathOperator{\Con}{{\rm Con}}
\DeclareMathOperator{\Lie}{{\rm Lie}}
\DeclareMathOperator{\At}{{\rm At}}
\DeclareMathOperator{\Ad}{{\rm Ad}}
\DeclareMathOperator{\ad}{{\rm ad}}

\DeclareMathOperator{\Der}{{\mathcal{D}\textit{er}}}
\DeclareMathOperator{\sHom}{{\mathcal{H}\textit{om}}}

\DeclareMathOperator{\GL}{{\rm GL}}
\DeclareMathOperator{\SL}{{\rm SL}}
\DeclareMathOperator{\Sp}{{\rm Sp}}

%\newcommand{\op}{\operatorname}

\newcommand{\mf}[1]{\mathfrak{#1}}
\newcommand{\mc}[1]{\mathcal{#1}}
\newcommand{\ms}[1]{\mathscr{#1}}
\newcommand{\bb}[1]{\mathbb{#1}}
\newcommand{\dv}{\vee\vee}
\newcommand{\ul}[1]{\underline{#1}}
\newcommand{\dual}{^\vee}
\newcommand{\norm}[1]{|\!|{#1}|\!|}
\newcommand{\diagram}[1]{\begin{align}
		\xymatrix{#1}
\end{align}}

\newcommand{\inv}{\textnormal{inv}}

\renewcommand{\ker}{\mathrm{Ker}}
\newcommand{\coker}{\mathrm{Coker}}
\renewcommand{\Im}{\mathrm{Im}}
\renewcommand{\Re}{\mathrm{Re}}


\newcommand{\doi}[1]{\href{https://doi.org/#1}{doi:#1}}
%%%%%%%%%%%%%%%%%%%%%%%%%%%%%%%%%%%%%%%%%%%%%%%%%%%%%%%%%%%%%
%\numberwithin{equation}{section}
\numberwithin{equation}{subsection}

\newtheorem{thm}[equation]{Theorem}
\newtheorem{corollary}[equation]{Corollary}
\newtheorem{lem}[equation]{Lemma}
\newtheorem{proposition}[equation]{Proposition}
\newtheorem{notation}[equation]{Notation}
\newtheorem{conjecture}[equation]{Conjecture}

\newtheorem{fact}{Fact}

\theoremstyle{definition}
\newtheorem{defn}[equation]{Definition}
\newtheorem{rem}[equation]{Remark}
\newtheorem{example}[equation]{Example}
\newtheorem{question}[equation]{Question}
\newtheorem{exercise}[equation]{Exercise}

%CustomTheorem
\theoremstyle{theorem}
\newtheorem{innercustomthm}{Theorem}
\newenvironment{customthm}[1]
{\renewcommand\theinnercustomthm{#1}\innercustomthm}
{\endinnercustomthm}

%%%%%%%% Symbols/letters/numbers for author section footnotes. 
\makeatletter
\newcommand\fnsymb[1]{\textsuperscript{\@fnsymbol{#1}}}
\newcommand\fnletter[1]{\lowercase{\textsuperscript{\@alph{#1}}}}
\newcommand\fnnum[1]{\textsuperscript{#1}}
\makeatother

%%%%%%%% Address, Email, Corresponding author: 
\usepackage{marvosym} % For Email/Letter icons. 
\renewcommand{\email}[2][1]{\thanks{\textit{Email address}#1: \href{mailto:#2}{#2}}}
\newcommand{\emailbox}[2][1]{\thanks{\Letter#1: \href{mailto:#2}{#2}}}
\renewcommand{\address}[2][1]{\thanks{\textit{Address}#1: #2}} 
%\newcommand{\corauth}[2][1]{\thanks{#2Corresponding author: #1}} 
\newcommand{\corauth}[2][]{\thanks{#2Corresponding author}} 



\begin{document}
	
	\baselineskip=15.5pt 
	
	\title[On Semiprojective Moduli of $\mathcal{L}$-Twisted Principal Objects and Motives]{Lie algebroid Connections, Moduli of $\mathcal{L}$--twisted Principal Objects and motives}
	
	\author[S. Ghosh]{Samit Ghosh\fnsymb{1}}
	
	\author[A. Paul]{Arjun Paul\fnsymb{2}}
	
	\email[\fnsymb{1}]{sg23rs005@iiserkol.ac.in}
	
	\email[\fnsymb{2}]{arjun.paul@iiserkol.ac.in}
	
	\address[]{Department of Mathematics and Statistics, 
		Indian Institute of Science Education and Research Kolkata, 
		Mohanpur - 741 246, Nadia, West Bengal, India.
	} 
	
	
	\thanks{Corresponding author: Arjun Paul}
	
	\subjclass[2020]{14H60, 53B15, 70G45, 17B45}
	
	\keywords{Lie algebroids, Tannakian categories, Moduli spaces, Higgs bundles, Integrable connections, Semiprojective varieties, Motives, $E$--polynomial.} 
	
	\begin{abstract}
    Let \(X\) be an irreducible smooth complex projective variety, and let \(G\) be a connected reductive linear algebraic group over \(\mathbb{C}\). In this paper, we first classify integrable transitive algebraic Lie algebroids on $X$. We then introduce Higgs bundles associated to a Lie algebroid and study their moduli spaces. In particular, we show that the category of vector bundles equipped with integrable \(\mathcal{L}\)-connections and the category of \(\mathcal{L}\)-twisted Higgs bundles of semiharmonic type on \(X\) are neutral Tannakian categories, provided that \(\mathcal{L}\) is a transitive Lie algebroid. Using this Tannakian framework, we obtain a characterization of principal \(G\)-bundles with integrable \(\mathcal{L}\)-connections and \(\mathcal{L}\)-twisted principal \(G\)-Higgs bundles of semiharmonic type on \(X\), and construct their moduli spaces via Mumford's geometric invariant theory. We further introduce the notion of the \(\mathcal{L}\)-Hodge moduli space for principal \(G\)-bundles and prove that the moduli spaces of principal \(G\)-bundles with integrable \(\mathcal{L}\)-connections, \(\mathcal{L}\)-twisted principal \(G\)-Higgs bundles of harmonic type, and the associated \(\mathcal{L}\)-Hodge moduli spaces are semiprojective varieties. Finally, using the semiprojectivity of the \(\mathcal{L}\)-Hodge moduli spaces for principal \(G\)-bundles, we obtain a description of smooth locus of these moduli spaces in the Grothendieck ring of varieties and establish a motivic non-abelian Hodge correspondence type theorem.
\end{abstract}
	
	\baselineskip=15.5pt 
	
	\date{Last updated on \today\,at \currenttime\,(IST)}
	
	\maketitle 
	
    \begin{center}
		\small
		\tableofcontents
	\end{center}

	
	
	

	
	
%%%%%%%%%%%%%%%%%%%%%%%%%%%%%%%%%%%%%%%%%%%%%%%%%%%%%%%%%%%%%%%%%%%%%%%%%%%%%%%%%%%%%%%%%%%%%%%%%%%%%%%%%%%%%%%%%%%%%%%%%%%%%%%%%%%%%%%%%%%%%%%%%%%%

\section{Introduction}\label{sec:introduction}

In \cite{CT.Simpson-1994}, C.~T.~Simpson introduced the notion of $\Lambda$--modules and constructed their moduli spaces using geometric invariant theory. 
The category of $\Lambda$--modules provides a unifying framework that simultaneously generalizes several important geometric structures appearing in differential geometry, algebraic geometry, and mathematical physics (see, for example, \cite{Cortes-Martinez-2004}, \cite{Lazzarini-Masson-2012}). 
Notable examples include Higgs bundles \cite{N.J.Hitchin-1987,CT. Simpson-1992,CT.Simpson-1995}, twisted Higgs bundles \cite{N.Nitsure-1991}, integrable connections \cite{CT. Simpson-1992}, and logarithmic or meromorphic connections \cite{P.Deligne-1970,N.Nitsure-1993,P.Boalch-2002,Bremer-Sage-2012}.

Beilinson and Bernstein \cite{Beilinson-Bernstein-1993} introduced the notion of a \textit{\(D\)-algebra} over a smooth projective variety \(X\). A \(D\)-algebra is an \(\mathcal{O}_{X}\)-algebra over \(X\) satisfying properties analogous to those of \(\mathcal{D}_{X}\), the sheaf of differential operators on \(X\). In \cite{P.Tortella-2011,P.Tortella-2011T}, P.~Tortella established a correspondence between holomorphic (or algebraic) Lie algebroids on $X$ and \(D\)-algebras \(\Lambda\) over \(X\). This correspondence provides a bridge between these two classes of objects. Using this framework, Tortella constructed moduli spaces of holomorphic integrable Lie algebroid connections on vector bundles via Simpson’s formalism of \(\Lambda\)-modules.
In \cite{K.Libor-2009,K.Libor-2010}, moduli spaces of Lie algebroid connections and flat Lie algebroid connections associated to a fixed Lie algebroid were constructed using analytic methods.


 
In \cite{Ghosh-Paul-2025}, Ghosh and Paul generalized the notion of holomorphic (or algebraic ) Lie algebroid connections to the setting of holomorphic (or algebraic) principal $G$--bundles over smooth projective varieties. Earlier, in \cite{CT.Simpson-1995}, Simpson constructed moduli spaces of principal \(G\)-bundles equipped with integrable \(\mathcal{T}_{X}\)-connections, as well as moduli spaces of principal \(G\)-Higgs bundles of harmonic type, by interpreting principal objects through the Tannakian formalism in terms of associated vector bundles and then applying Mumford's geometric invariant theory. More recently, in \cite{Alfaya-Oliveire-2024}, the semiprojectivity of the moduli space of algebraic Lie algebroid connections on vector bundles, as well as of the corresponding $\mathcal{L}$--Hodge moduli space, was established. 
The semiprojectivity of moduli spaces of principal $G$--bundles with $\lambda$--connections was studied in \cite{Roy-Singh-2024}.

A holomorphic (or algebraic) Lie algebroid connection generalizes the classical notion of a holomorphic (or algebraic) connection \cite{Atiyah-1957}, a structure that naturally arises on complex manifolds, Poisson varieties, and spaces endowed with singular foliations. 
Moduli spaces of vector bundles with integrable $\mathcal{L}$--connections, $\mathcal{L}$--twisted Higgs bundles, and $\mathcal{L}$--Hodge moduli spaces for vector bundles—are typically quasi-projective and non-proper. 
Semiprojectivity \cite{T.Hausel-F. Villegas} provides a natural framework in which this non-properness is controlled via a $\mathbb{C}^{\ast}=\mathbf{G}_{m}$--action with projective fixed-point locus and well-defined limits. 
This notion plays a central role in the study of topology, Hodge theory, and degeneration phenomena for moduli spaces. 
For a smooth semiprojective variety, the Grothendieck motivic class admits an explicit description in terms of the projective fixed-point locus via the Bia{\l}ynicki--Birula decomposition.


The aim of this paper is to develop a systematic moduli-theoretic framework for principal $G$--bundles with integrable $\mathcal{L}$--connections and $\mathcal{L}$--twisted principal $G$--Higgs bundles of semiharmonic type over $X$, as well as the associated $\mathcal{L}$--Hodge moduli spaces for principal $G$--bundles. We study the global geometry of these moduli spaces, with particular emphasis on establishing their semiprojectivity, and motivic classes of smooth loci of resulting semiprojective moduli spaces.



%%%%%%%%%%%%%%%%%%%%%%%%%%%%%%%%%%%%%%%%%%%%%%%%%%%%%%%%%%%%%%%%%%%%%%%%%%


\subsection{Main Contribution}

Let \(X\) be an irreducible smooth complex projective variety, let \(\mathcal{L}\) be a transitive algebraic Lie algebroid on \(X\), and let \(G\) be a connected complex reductive algebraic group. The ideas developed in this paper are motivated by and build upon the works of \cite{Alfaya-Oliveire-2024}, \cite{CT.Simpson-1994}, \cite{CT.Simpson-1995}, \cite{Beilinson-Bernstein-1993}, and the PhD thesis of Pietro Tortella \cite{P.Tortella-2011T}.

\medskip

\noindent
\textbf{(1) Characterization of integrable transitive algebraic Lie algebroid over $X$.} The theory of classical Higgs bundles is by now very well developed. Twisted Higgs bundles, particularly those twisted by line bundles and low rank vector bundles (for instance rank one and rank two), have also been extensively studied (cf. \cite{N.Nitsure-1991},\cite{Gallego-Oscar-Narasimhan-2024}). However, Higgs bundles associated to an arbitrary Lie algebroid appear to be largely unexplored and provide a genuinely new class of geometric objects. This highlights the importance of understanding and classifying Lie algebroids over smooth projective varieties, especially over curves.

For smooth projective curves of genus \(g \geqslant 2\), Lie algebroids of rank one are completely classified (see \cite[Proposition 2.1]{GLP-2018} and \cite[Theorem 1.2]{Alfaya-Oliveire-2024}). In contrast, no general classification is known for higher-rank Lie algebroids, and even the case of rank two remains poorly understood. 

In this paper, we characterize an important class of Lie algebroids, namely integrable transitive algebraic Lie algebroids. Although the corresponding characterization is known for smooth transitive Lie algebroids over connected smooth manifolds (cf. \cite{Moerdijk-2003}).


\begin{thm}\eqref{cor:classifcation of integrable transitive-Lie algebroid}.
Let $X$ be a smooth irreducible projective algebraic variety over $\mathbb{C}$ and let $\mathcal{A}$ be a transitive algebraic Lie algebroid on $X$. Then $\mathcal{A}$ is algebraically integrable if and only if $\mathcal{A}$ is isomorphic to the Atiyah algebroid of a principal
$G$--bundle on $X$, for some algebraic group $G$.
\end{thm}


\medskip

\noindent
\textbf{(2) Tannakian description of $\mathcal{L}$--twisted principal $G$--objects and construction of moduli spaces.}
In \cite[\S\,9]{CT. Simpson-1992}, C.~T.~Simpson proved that the category of vector bundles equipped with integrable connections and the category of Higgs bundles of semiharmonic type are neutral Tannakian categories. In this paper, we extend this framework to the setting of transitive Lie algebroids, establishing analogous Tannakian descriptions for vector bundles with integrable \(\mathcal{L}\)-connections and \(\mathcal{L}\)-twisted Higgs bundles.


\begin{proposition}\eqref{prop: tannakian category Vect_L(X)},\eqref{prop:Tannakian-L-Higgs}
Let $X$ be a smooth irreducible projective algebraic variety over $\mathbb{C}$
	\begin{enumerate}
		\item The category $\texttt{Vect}^{\emph{int}}_{\mathcal{L}}(X)$ of vector bundles equipped with integrable $\mathcal{L}$--connections is a neutral Tannakian category.
		\item The category $\texttt{Higg}^{\emph{sh}}_{\mathcal{L}}(X)$ of $\mathcal{L}$--twisted Higgs bundles of semiharmonic type on $X$ is also a neutral Tannakian category.
	\end{enumerate}
\end{proposition}

In \cite[\S\,9]{CT.Simpson-1995}, C.~T.~Simpson developed a Tannakian description of principal \(G\)-bundles endowed with integrable connections and of principal \(G\)-Higgs bundles of semiharmonic type on \(X\). In this paper, we extend this framework to the setting of algebraic transitive Lie algebroids \(\mathcal{L}\) on \(X\).

We first generalize the notion of \(\mathcal{L}\)-twisted Higgs bundles  to principal \(G\)-bundles. Our first main result provides a Tannakian description of principal \(G\)-bundles equipped with integrable \(\mathcal{L}\)-connections and \(\mathcal{L}\)-twisted principal \(G\)-Higgs bundles.

\begin{thm}\eqref{lem:PB-with-L-conection-Tannakian equivalence},\eqref{lem: L-twisted PB-Higgs bundle-tannakian}.
	\label{thm:tannakian-principal-objects}
	Let $\mathcal{L}$ be an algebraic transitive Lie algebroid. The assignment $E_{G}\longmapsto \rho_{E_{G}}$ defined in
	\eqref{eq:nori-functor} induces the following equivalences of categories.
	\begin{enumerate}
		\item The category of principal $G$--bundles on $X$ equipped with integrable $\mathcal{L}$--connections is equivalent to the category of $G$ torsor in $\texttt{Vect}^{\emph{int}}_{\mathcal{L}}(X)$,
		such that for every closed point $x\in X$, the functor $
		V\longmapsto \rho(V)_{x}$ defines a fiber functor on $\texttt{Rep}(G)$.
		
		\item The category of $\mathcal{L}$--twisted principal $G$--Higgs bundles of semiharmonic type on $X$ is equivalent to the category of $G$--torsor in $\texttt{Higg}^{\emph{sh}}_{\mathcal{L}}(X)$,
		such that for every closed point $x\in X$, the functor $
		V\longmapsto \rho(V)_{x}$  defines a fiber functor on $\texttt{Rep}(G)$.
	\end{enumerate}
\end{thm}


Using the above Tannakian description together with Mumford’s geometric invariant theory, we construct moduli spaces of $\mathcal{L}$--twisted principal $G$--bundles with integrable $\mathcal{L}$--connection and $\mathcal{L}$--twisted principal $G$--Higgs bundles of harmonic type on $X$.

\begin{thm}\eqref{cor:moduli of int-L-connection PB}
	There exists a quasi projective variety $\mathcal{M}^{\text{DR}}_{\mathcal{L}}(X,G)$ which is universally co-represents the moduli functor $\mathcal{M}^{\text{DR}^{\natural}}_{\mathcal{L}}(X,G)$ defined in \eqref{moduli of principal G- bundles with integrable connection}. 
	In particular, $\mathcal{M}^{\text{DR}}_{\mathcal{L}}(X,G)$ is a quasi-projective variety.
\end{thm}

\begin{thm}\eqref{thm:moduli-L-twisted-principal-G-Higgs}
	There exists a quasi projective variety $\mathcal{M}^{\text{Dol}}_{\mathcal{L}}(X,G)$ which is universally co-represents the moduli functor $\mathcal{M}^{\text{Dol}^{\natural}}_{\mathcal{L}}(X,G)$ defined in \eqref{moduli-of-L-twisted princpal G--Higgs bundles}. 
	In particular, $\mathcal{M}^{\text{Dol}}_{\mathcal{L}}(X,G)$ is a quasi-projective variety.
\end{thm}

\medskip

\noindent
\textbf{(3) $\mathcal{L}$--Hodge moduli spaces for principal $G$ bundles and semiprojectivity.}
We further generalize the notion of \(\mathcal{L}\)-Hodge moduli spaces introduced in \cite{Alfaya-Oliveire-2024} to the setting of principal \(G\)-bundles (see \eqref{The-L-Hodge-moduli-spaces-for-principal-bundles}). Another main result of this paper concerns the global geometry of the resulting moduli spaces.

\begin{thm} \eqref{semiprojectivity of MLDOL(X,G)}, \eqref{semiprojectivity of MLHOD(X,G)}
Let $X$ be a smooth irreducible projective variety and $\mathcal{L}$ be an algebraic transitive Lie algebroid.
	\begin{enumerate}
		\item The moduli space $\mathcal{M}^{\mathrm{Dol}}_{\mathcal{L}}(X,G)$ of $\mathcal{L}$--twisted principal $G$--Higgs bundles of harmonic type is a semiprojective variety.
		\item The $\mathcal{L}$--Hodge moduli space $\mathcal{M}^{\mathrm{Hod}}_{\mathcal{L}}(X,G)$ of principal $G$--bundles is a semiprojective variety.
	\end{enumerate}
\end{thm}


\noindent
\textbf{(4) Motivic classes of the resulting moduli spaces in the Grothendieck ring of varieties.}

Let \(X\) be a smooth projective curve of genus \(g \geqslant 2\) and $\gcd(r,d) = 1$. For Higgs bundles, let \(\mathcal{M}^{\mathrm{Dol}}(r,d)\) denote the moduli space of Higgs bundles of rank \(r\) and degree \(d\). For connections, fix a point \(x \in X\), and let \(\mathcal{M}^{\mathrm{DR}}(r,d)\) denote the moduli space of logarithmic connections with allowed poles at \(x\), and with residue fixed by the degree \(d\). Using the semiprojectivity and smoothness of \(\mathcal{M}^{\mathrm{Hod}}(r,d)\), Hoskins and Lehalleur \cite{Hoskins-Lehalleur-2021} established what they called the \textquotedblleft motivic non-abelian Hodge correspondence\textquotedblright\ by proving an equality between the Voevodsky motives of \(\mathcal{M}^{\mathrm{Dol}}(r,d)\) and \(\mathcal{M}^{\mathrm{DR}}(r,d)\). Later, David Alfaya and Andr\'e Oliveira generalized these results to the setting of rank one Lie algebroids whose degree is less than that of \(\mathcal{T}_{X}\) (cf.~\cite{Alfaya-Oliveire-2024}).

In general, the \(\mathcal{L}\)-Hodge moduli space \(\mathcal{M}^{\mathrm{Hod}}_{\mathcal{L}}(X,G)\) is not smooth. Nevertheless, by restricting to the smooth loci of the corresponding moduli spaces, one can still obtain analogous motivic results.

Let \(Y\) be a semiprojective variety over \(\mathbb{C}\), and denote by \(\widehat{Y}\) its smooth locus. Observe that the smooth locus of a semiprojective variety is again semiprojective.
\begin{thm}\eqref{thm:motivic equality of dr-dol-Hod principal moduli space}.
	Let $X$ be a smooth irreducible projective variety over $\mathbb{C}$, and $\mathcal{L} = (V,[·,·],\delta)$ be a transitive Lie algebroid on $X$. Then the following equalities hold $\hat{\mathcal{K}}(Var_{\mathbb{C}})$,
	
	\begin{enumerate}
		\item $[\widehat{\mathcal{M}^{\text{Dol}}_{\mathcal{L}}(X,G)}]=\sum_{\alpha\in X^{\mathbf{G}_{m}}}\mathbb{L}^{N_{\alpha}^{+}}[F_{\alpha}]$
		
		\item $[\widehat{\mathcal{M}^{\text{DR}}_{\mathcal{L}}(X,G)}] = [\widehat{\mathcal{M}_{\mathcal{L}}^{\text{Dol}}(X,G)}, \ \  [\widehat{\mathcal{M}_{\mathcal{L}}^{\text{Hod}}(X,G)]} = \mathbb{L}[\widehat{\mathcal{M}_{\mathcal{L}}^{\text{Dol}}(X,G)}]$,
		
		\item $E(\widehat{\mathcal{M}_{\mathcal{L}}^{\text{DR}}(X,G))} = E(\widehat{\mathcal{M}_{\mathcal{L}}^{\text{Dol}}(X,G))}$, \ \  $E(\widehat{\mathcal{M}_{\mathcal{L}}^{\text{Hod}}(X,G))} = xyE(\widehat{\mathcal{M}_{\mathcal{L}}^{\text{Dol}}(X,G))}$,
		
		\item we have an isomorphism of Hodge structures,
		\[H^{\bullet}(\widehat{\mathcal{M}_{\mathcal{L}}^{\text{DR}}(X,G)}) \cong H^{\bullet}(\widehat{\mathcal{M}_{\mathcal{L}}^{\text{Dol}}(X,G)})\]and both $\widehat{\mathcal{M}_{\mathcal{L}}^{\text{DR}}(X,G)}$ and $\widehat{\mathcal{M}_{\mathcal{L}}^{\text{Hod}}(X,G)}$ have pure mixed Hodge structure.
		
		
		\item A smooth fixed point 
		\[
		\alpha \in \mathcal{M}^{\mathrm{DR}}_{\mathcal{L}}(X,G)
		\quad 
		\left(\text{respectively} \ 
		\mathcal{M}^{\mathrm{Dol}}_{\mathcal{L}}(X,G),\ 
		\mathcal{M}^{\mathrm{Hod}}_{\mathcal{L}}(X,G)
		\right)
		\]
		is very stable if and only if the corresponding upward flow \(F^{+}_{\alpha}\) is closed.
	\end{enumerate}
	
\end{thm}

\subsection{Open question and future work}. 
We are presently able to construct moduli spaces of \(\mathcal{L}\)-twisted principal \(G\)-objects only in the case where \(\mathcal{L}\) is a transitive Lie algebroid. It is natural to expect that one can construct moduli spaces of \(\mathcal{L}\)-twisted principal \(G\)-objects for arbitrary Lie algebroids \(\mathcal{L}\), together with the smooth loci of these moduli spaces, and establish a motivic invariance theorem analogous to \cite[Theorem~7.1]{Alfaya-Oliveire-2024}, at least in the case of rank one and rank two Lie algebroids.

To achieve this, it would be necessary to investigate the fixed-point loci of the corresponding moduli spaces. It would also be interesting to compute the homotopy groups of the smooth loci of the resulting moduli spaces.

%%%%%%%%%%%%%%%%%%%%%%%%%%%%%%%%%%%%%%%%%%%%%%%%%%%%%%%%%%%%%%%%%%%%%%%%%%%%%%%%%%%%%%%%%%%%%%%%%%%%%%%%%%%%%%%%%%%%%%%%%%%%%%%%%%%%%%%%%%%%%%%%%



\section{Lie algebroid connections, and moduli spaces}

Throughout this paper, let $X$ be an irreducible smooth projective variety over $\mathbb{C}$, and let $\mathcal{O}_X(1)$ denote a fixed very ample line bundle on $X$. For computational purposes, we shall mainly work with irreducible smooth projective curves over $\mathbb{C}$ of genus $g \geq 2$. Let $G$ be a linear algebraic group over $\mathbb{C}$. Unless explicitly stated otherwise, all objects considered in this paper are algebraic. In particular, we work with algebraic vector bundles, algebraic principal $G$-bundles (locally trivial in the \'{e}tale topology; see, for example, \cite{Sch-2008, C.Sorger-1999}), algebraic Lie algebroids defined over irreducible smooth projective varieties over $\mathbb{C}$. This convention will be assumed throughout, even when the adjective ``algebraic'' is omitted.

We shall also use the standard identification between algebraic vector bundles and locally free sheaves on $X$.

\subsection{Lie algebroids and $\mathcal{L}$-connection on $G$-bundles}
Let $\mc O_X$ be the sheaf of functions on $X$, and let 
$\mathcal{T}_{X}$ be the tangent bundle of $X$. 

\begin{defn}\cite[\S\,3.1]{Alfaya-Oliveire-2024}
	An {\it algebraic Lie algebroid} on $X$ is a triple $\mc L := (V, [\cdot, \cdot], \delta)$, where 
	\begin{enumerate}[(i)]
		\item $V$ is a vector bundle on $X$, 
		
		\item $[\cdot, \cdot] : V\times V \to V$ is a $\bb C$--bilinear skew-symmetric 
		morphism of sheaves such that for all locally defined sections 
		$u, v, w$ of $V$, the following Jacobi identity holds: 
		$$[u, [v, w]]+[v, [w, u]]+[w, [u, v]] = 0;$$
		
		\item $\delta : V \to \mathcal{T}_{X}$ is a vector bundle homomorphism 
		satisfying the following properties: for all locally defined sections $s, t$ of $V$ and 
		locally defined section $f$ of $\mc O_X$, we have 
		\begin{enumerate}[(a)]
			\item {\it Compatibility of Lie algebra structures}: 
			$\delta([s, t]) = [\delta(s), \delta(t)]$, and 
			\item {\it Leibniz rule}: $[fs, t] = f[s, t] - \delta(t)(f)s$.
		\end{enumerate}
	\end{enumerate}
	The homomorphism $\delta$ is called the {\it anchor map} of the Lie algebroid $\mc L$. 
	The {\it degree} and the {\it rank} of $\mathcal{L}$ is defined to be the degree and the rank, respectively, of the underlying vector bundle $V$ of $\mc L$. A Lie algebroid $\mathcal{L} = (V,[\cdot,\cdot],\delta)$ on $X$ is called
	\emph{transitive} if its anchor map
	\[
	\delta : V \longrightarrow \mathcal{T}_X
	\]
	is surjective.
	A Lie algebroid $\mathcal{L}=(V,[\cdot,\cdot],\delta)$ on $X$ will be called \textit{split} if there is an $\mathcal{O}_{X}$--linear homomorphism,
	\[\gamma:\mathcal{T}_{X}\rightarrow V\]
	such that $\delta\circ \gamma = \text{Id}_{\mathcal{T}_{X}}$. A Lie algebroid $\mathcal{L}$ on $X$ will be called \textit{nonsplit} if it is not split.
\end{defn}

A \textit{Lie algebroid map} $f:\mathcal{L}\rightarrow \mathcal{L}'$ between $\mathcal{L}=(V,[\cdot,\cdot]_{V},\delta_{V})$ and $\mathcal{L}'=(V',[\cdot,\cdot]_{V'},\delta_{V'})$ is an algebraic $\mathbb{C}$--Lie algebra bundle map  $f:V\rightarrow V'$ such that $\delta_{V'}\circ f = \delta_{V}$. A \textit{Lie algebroid isomorphism} is a Lie algebroid map which is an isomorphism of the underlying bundles; in that case the Lie algebroids are said to be \textit{isomorphic}.
\begin{example}[Standard Examples of Lie Algebroids]
Let $X$ be a smooth projective variety over $\mathbb{C}$ and $D \subset X$ is a simple normal crossing divisor.
	
\begin{enumerate}
\item \textit{Tangent Lie algebroid.}  
The tangent bundle $\mathcal{T}_X$ with anchor
		\[
		\rho = \mathrm{id}_{\mathcal{T}_X} : \mathcal{T}_X \to \mathcal{T}_X
		\]
		and the usual Lie bracket of vector fields.  
		This Lie algebroid is \emph{transitive}.
		
		\item \textit{Logarithmic tangent Lie algebroid.}  
		Define $\mathcal{T}_X(-\log D)$ to be the subbundle of $\mathcal{T}_{X}$ whose section are the vector fields $\mathcal{V}$ such that $\mathcal{V}(I_{D})\subseteq I_{D}$, for $I_{D}$ the ideal sheaf of $D$ (in other words, $\mathcal{T}_{X}(-\log D)$ consisting of vector fields tangent to $D$). The Lie bracket of sections in $\mathcal{T}_{X}(-\log D)$ is still in  $\mathcal{T}_{X}(-\log D)$, so that it inherits a Lie algebroid structure from $\mathcal{T}_{X}$
		with the natural inclusion anchor
		\[
		i : \mathcal{T}_X(-\log D) \hookrightarrow \mathcal{T}_X.
		\]
		This Lie algebroid is generally \emph{not transitive}.
		
		\item \textit{Atiyah Lie algebroid of a principal bundle.} \label{example:Atiyah Lie algebroid of a principal bundle}
		Let $E\xrightarrow{\pi} X$ be a principal $G$-bundle on $X$. Denote by $\mathfrak{g}$ the Lie algebra of $G$. Consider the exact sequence obtained from the differential of $\pi$,
		\[0\longrightarrow E\times \mathfrak{g}\longrightarrow \mathcal{T}_{E}\longrightarrow \pi^{*}\mathcal{T}_{X}\longrightarrow 0\]The group $G$ acts naturally on each element of this sequence, and after taking the quotient one obtains the exact sequence of vector bundles over $X$
		\begin{equation}\label{Atiyah sequence for PB}
		0 \longrightarrow \operatorname{ad}(E)
		\longrightarrow \frac{\mathcal{T}_{E}}{G}
		\xrightarrow{\;\rho\;}
		\mathcal{T}_X
		\longrightarrow 0.
		\end{equation}
		The \textit{Atiyah algebroid} of $E$ is $\text{At}(E):=\mathcal{T}_{E}/G$, with anchor equal to the quotient map of the sequence, and bracket induced by the bracket of the vector fields in $\mathcal{T}_{E}$. This Lie algebroid is \emph{transitive}. An algebraic $\mathcal{T}_{X}$--connection (or, simply connection) on $E$ is a map of algebraic vector bundle $\nabla:\mathcal{T}_{X}\rightarrow \text{At}(E)$ such that $\rho\circ \nabla = \text{Id}_{\mathcal{T}_{X}}$. We say that the connection $\nabla$ is \textit{integrable} if moreover $\nabla$ is a morphism of Lie algebroid.
		
		Assume that $E$ does not admit any algebraic connection. For example, set $G = GL(r, \mathbb{C})$ and take $E$ to be the  principal $GL(r, \mathbb{C})$–bundle over $X$ associated to a vector bundle of rank $r$ and nonzero degree over $X$. Then the Lie algebroid $(\text{At}(E), [\cdot,\cdot],\rho)$ is nonsplit. On the other hand, if $E$ admits a connection, then the Lie algebroid $(\text{At}(E), [\cdot,\cdot],\rho)$ in is split.
		
		\item \textit{Log Picard Lie algebroid.} A log Picard algebroid $(\mathcal{A}, [\cdot,\cdot], \sigma, e)$ on $(X, D)$ is a locally free $\mathcal{O}_{X}$-module $\mathcal{A}$ equipped with a Lie bracket $[\cdot,\cdot]$, a bracket-preserving morphism of $\mathcal{O}_{X}$--modules $\sigma : \mathcal{A} \rightarrow
		\mathcal{T}_{X}(- \log D)$, and a central section $e$, such that the Leibniz rule
		\[[a_{1},fa_{2}]=f[a_{1},a_{2}]+\sigma(a_{1})(f)a_{2}\]
         holds for all $f \in \mathcal{O}_{X}$, $a_{1}, a_{2} \in \mathcal{A}$, and the sequence
		\[0\longrightarrow \mathcal{O}_{X}\xrightarrow{e}\mathcal{A}\xrightarrow{\sigma} \mathcal{T}_{X}(-\log D)\longrightarrow 0\]is exact. These objects were considered by Marco Gualtieri and Kevin Luk in \cite{GL-2021}.
		
		
		\item \textit{Trivial Lie algebroid.}  
		Any algebraic vector bundle $V$ on $X$ can be viewed as a Lie algebroid by equipping it with zero bracket $[s_1,s_2]=0$
		and zero anchor $\rho = 0 : V \to \mathcal{T}_X.$
		This Lie algebroid is \emph{not transitive}.
		
	\end{enumerate}
\end{example}

Let $\mathcal{L} = (V ,[\cdot,\cdot],\delta)$ be a Lie algebroid. We now define a differential on the complex of exterior powers, $$\Omega^{\bullet}_{\mathcal{L}} := \bigwedge^{\bullet}V^{*}, \ d_{\mathcal{L}}:\Omega^{k}_{\mathcal{L}}\rightarrow \Omega^{k+1}_{\mathcal{L}},$$
generalising the classical de Rham complex $d :\Omega^{k}_{X}\rightarrow \Omega^{k+1}_{X}$. In degree $0$, define $d_{\mathcal{L}}: \mathcal{O}_{X} \rightarrow V^{\ast}$ as the
composition of the canonical differential $d : \mathcal{O}_{X}\rightarrow \Omega_{X}^{1}(=\mathcal{T}^{*}_{X})$ with the dual of the anchor, $\delta^{*}: \Omega^{1}_{X} \rightarrow V^{\ast}$. Thus, given $v \in V$ and $f$ a local algebraic function on $X$,
\[d_{\mathcal{L}}(f)(v):=df(\delta(v))=\delta(v)(f)\]
The map $d_{\mathcal{L}} : O_{X} \rightarrow V^{*}(=\Omega^{1}_{\mathcal{L}})$ is clearly a $V^{*}$-valued derivation, we can extend it to higher order exterior powers through the usual recursive equation, for details see \cite[\S\,3.1]{Alfaya-Oliveire-2024}. This differential satisfies $d_{\mathcal{L}}^{2}=0$, so $(\Omega^{\bullet}_{\mathcal{L}},d_{\mathcal{L}})$ is a complex, called the \textit{Chevalley-Elienberg-de Rham complex} of $\mathcal{L}$. Note that $d_{\mathcal{L}}=0$ for a trivial algebroid.

We denote by $\tau^{\geqslant r}\Omega^{\bullet}_{\mathcal{L}}$
L the \textit{b$\hat{e}$te filtration }of the complex $\Omega^{\bullet}_{\mathcal{L}}$, which is the complex,
\[\tau^{\geqslant r}\Omega^{\bullet}_{\mathcal{L}}=
\begin{cases}
0 \ \ \ \ \ \text{if} \ k<r,\\
\Omega^{k}_{\mathcal{L}} \ \ \text{if} \ k\geqslant r
\end{cases}\]
Let $\mathfrak{U} = \{U_{i}\}$ be a sufficiently fine open covering of $X$ , such that we have an isomorphism between the sheaf and $\check{C}$ech cohomology over it. Consider the double complex
\[
K^{p,q}_{\mathcal{L}}:=\check{C}^{q}(\mathfrak{U},\Omega^{p}_{\mathcal{L}})
\] with differentials given by $d_{\mathcal{L}}$ and the $\check{C}$ech coboundary, and recall that its associated total complex $T^{\bullet}_{\mathcal{L}}$ computes the hypercohomology of $\Omega^{\bullet}_{\mathcal{L}}$. Remark that
the hypercohomology of the complex $\tau^{\geqslant r}\Omega^{\bullet}_{\mathcal{L}}$ is isomorphic to the cohomology of the
complex of vector spaces
\[T^{k}_{\tau^{\geqslant r}\Omega^{\bullet}_{\mathcal{L}}}=\bigoplus_{p+q=k,p\geqslant r}K^{p,q}_{\mathcal{L}}\]

\begin{defn}\label{def:Lie-alg-conn-on-VB}
	Let $\mathcal{E}$ be a vector bundle on $X$, and let $\mathcal{L}$ be an algebraic Lie algebroid over $X$. An {\it $\mc L$-- connection} on  
	$\mc E$ on $X$ is a $\bb C$--linear homomorphism of sheaves 
	$$D_{\mathcal{L}} : \mc E \longrightarrow \mc E\otimes \Omega^{1}_{\mathcal{L}}$$ 
	satisfying the {\it $\delta^*$--twisted Leibniz rule}: 
	\begin{equation}\label{eqn:phi^*-twisted-Leibniz-rule}
		D_{\mathcal{L}}(f\cdot s) = fD_{\mathcal{L}}(s)+s\otimes\delta^*(df), 
	\end{equation}
	for all locally defined section $s$ of $\mc E$ and for all locally defined 
	section $f$ of $\mc O_X$. 
\end{defn}
\begin{defn}
Let $(\mathcal{E},D_{\mathcal{L}})$ be an $\mathcal{L}$-connection on $X$. The composition map $D^{2}_{\mathcal{L}}:\mathcal{E}\rightarrow \mathcal{E}\otimes \Omega^{2}_{\mathcal{L}}$, called the curvature of $(\mathcal{E},D_{\mathcal{L}})$. An $\mathcal{L}$-connection is \textit{integrable} if its curvature vanishes.
\end{defn}
Let $X$ be an irreducible smooth projective curve over $\mathbb{C}$. It is a well-known fact that not every algebraic vector bundle admits an algebraic $\mathcal{T}_{X}$--connection. A theorem of Atiyah \cite{Atiyah-1957} states that an algebraic vector bundle $\mathcal{E}$ admits an algebraic $\mathcal{T}_{X}$--connection if and only if each indecomposable component of $\mathcal{E}$ has degree zero. It can be shown that $\mathcal{E}$ admits a tautological Lie algebroid connection with respect to the Lie algebroid $(\operatorname{At}(\mathcal{E}), [\cdot,\cdot], \rho)$ (see, for example, \cite{ABKA-2025}); moreover, this tautological Lie algebroid connection on $\mathcal{E}$ is integrable. Furthermore, \cite{ABKA-2025} established an Atiyah--Weil type criterion for the existence of algebraic $\mathcal{L}$--connections, which we state below.

\begin{thm}\cite[\S\,1, Theorem 1.1]{ABKA-2025}
Let $X$ be an irreducible smooth projective curve over $\mathbb{C}$, and $\mathcal{L}=(V, [\cdot,\cdot],\delta)$ be a algebraic Lie algebroid over $X$.
\begin{itemize}
\item If $\mathcal{L}$ be a nonsplit Lie algebroid. Then every vector bundle $\mathcal{E}$ on $X$ admits a Lie algebroid connection for $\mathcal{L}$.
\item If $\mathcal{L}$ be a split Lie algebroid. Then the following two statements are equivalent:
\begin{enumerate}
\item $\mathcal{E}$ admits a Lie algebroid connection for $\mathcal{L}$.
\item Each indecomposable component of $\mathcal{E}$ is of degree zero.
\end{enumerate}
\end{itemize}
\end{thm}


Let $\mathcal{F}$ be a vector bundle on $X$. The above theorem and example 2.1.2 \eqref{example:Atiyah Lie algebroid of a principal bundle} gives the following two corollary:
\begin{corollary}\cite[Corollary 5.1]{ABKA-2025}
Assume that $\mathcal{F}$ admits an $\mathcal{T}_{X}$--connection, and consider the Atiyah Lie algebroid $(\text{At}(\mathcal{F}), [\cdot,\cdot],\rho)$ associated to $\mathcal{F}$. Let $\mathcal{E}$ be an algebraic vector bundle on $X$. Then the following three statements are equivalent:
\begin{enumerate}
\item Each indecomposable component of $\mathcal{E}$ is of degree zero.
\item $E$ admits a Lie algebroid connection for $(\text{At}(\mathcal{F}),[\cdot,\cdot],\rho)$.
\item $E$ admits an integrable Lie algebroid connection for $(\text{At}(\mathcal{F}),[\cdot,\cdot],\rho)$.
\end{enumerate}
\end{corollary}

\begin{corollary}\cite[Corollary 5.2]{ABKA-2025}
Assume that $\mathcal{F}$ does not admit any $\mathcal{T}_{X}$--connection, and
consider the Atiyah Lie algebroid $(\text{At}(\mathcal{F}), [\cdot,\cdot],\rho)$ associated to $\mathcal{F}$. Then any holomorphic vector bundle $\mathcal{E}$ on $X$ admits a Lie algebroid connection for $\text{At}(\mathcal{F})$.
\end{corollary}




The notion of Lie algebroid connection to the case of principal bundles 
introduced in (\cite{Ghosh-Paul-2025}) and for reader's convenience recall the definitions.
Let $G$ be a linear algebraic group over $\mathbb{C}$ with the Lie algebra $\mf{g} := \Lie(G)$. 
Let $p : E_G \to X$ be a algebraic principal $G$--bundle on $X$. The adjoint representation 
$${\rm ad} : G \longrightarrow \GL(\mf{g})$$ 
of $G$ on its Lie algebra $\mf{g}$ gives rise to a vector bundle 
$$\ad(E_G) := E_G\times^{\rm ad}\mf{g}$$ 
on $X$, called the {\it adjoint vector bundle} of $E_G$. If $E=\textbf{Fr}(\mathcal{E})$ is the frame bundle of a vector bundle $\mc E$ of rank $n$ on $X$, 
then we have $\ad(E) \cong \End(\mc E)$, the endomorphism bundle of $\mc E$. The surjective submersion $p : E_G \to X$ gives rise to an exact sequence of  vector bundles,
\begin{equation}
	\xymatrix{
		0 \ar[r] & \ad(E_G) \ar[r] & \At(E_G) \ar[r]^-{d'p} & TX \ar[r] & 0 
	}
\end{equation}
called the \textit{Atiyah exact sequence} of $E_G$.

Fix a Lie algebroid $\mc L = (V, [\cdot\,,\,\cdot], \delta)$ on $X$, and consider the map 
$$\rho : \At(E_G)\oplus V \longrightarrow TX$$ 
defined by 
\begin{equation}
	\rho(\xi, v) = d'p(\xi)-\delta(v), 
\end{equation}
for all locally defined section $\xi$ of $\At(E_G)$ and locally defined section $v$ of $V$. Note that $\rho$ is a vector bundle homomorphism and
\begin{equation}
	\At_{\delta}(E_G) := \rho^{-1}(0) 
\end{equation}
is a vector bundle on $X$ (as $\rho$ is surjective). 
The restriction of the second projection map gives rise 
to a vector bundle homomorphism 
\begin{equation}\label{eqn:rho-tilde-map}
	\widetilde{\rho} : \At_{\delta}(E_G) \longrightarrow V
\end{equation} 
with kernel 
$$\ker(\widetilde{\rho}) = \ad(E_G).$$ 

Thus we have the following short exact sequence 
\begin{equation}\label{eqn:Atiyah-exact-seq-for-(V,phi)-valued-connection}
	0 \longrightarrow \ad(E_G) \longrightarrow \At_{\delta}(E_G) 
	\stackrel{\widetilde{\rho}}{\longrightarrow} V \longrightarrow 0 
\end{equation}
of vector bundles on $X$, which fits into the following 
commutative diagram 


\begin{equation}
\begin{CD}
	0@>>> \text{ad}(E_{G}) @>>> \At_{\delta}(E_G) @>\tilde{\rho}>> V @>>> 0\\
	@.    @| @VVV @VV\delta V @.\\
	0@>>> \text{ad}(E_{G}) @>>> \At(E_G) @>d'p>> TX @>>> 0\\
\end{CD}
\end{equation}
of vector bundle homomorphisms with all rows exact. This exact sequence is called the $\mathcal{L}$-\textit{Atiyah exact sequence for principal $G$-bundles}.

\begin{defn}\label{def:Lie-alg-conn-on-G-Bundles}\cite[\S\,2,Definition 2.2.7]{Ghosh-Paul-2025}
	An algebraic {\it $\mc L$-- connection on principal $G$-bundle} $E_G$, is a 
	vector bundle homomorphism 
	$$\nabla_{\mathcal{L}} : V \longrightarrow \At_{\delta}(E_G)$$
	such that $\widetilde{\rho}\circ\nabla_{\mathcal{L}} = \Id_V$, 
	where $\widetilde{\rho}$ is defined in \eqref{eqn:rho-tilde-map}. 
\end{defn}

The short exact sequence \eqref{eqn:Atiyah-exact-seq-for-(V,phi)-valued-connection} 
defines a cohomology class 
\begin{equation}\label{defn:V-valued-Atiyah-class-of-E_G}
	\Phi_{\mc L}(E_G) \in H^1(X, \ad(E_G)\otimes \Omega^{1}_{\mathcal{L}}), 
\end{equation}
such that the exact sequence \eqref{eqn:Atiyah-exact-seq-for-(V,phi)-valued-connection} 
splits algebraically if and only if $\Phi_{\mc V}(E_G) = 0$. 

\begin{proposition}\label{prop:E_G-admits-Lie-alg-conn-iff-Atiyah-cls-vanishes}\cite[\S\,2, Prop 2.2.9]{Ghosh-Paul-2025}
	A principal $G$--bundle $E_G$ on $X$ admits a $\mc L$--connection if and only if $\Phi_{\mc L}(E_G) = 0$. 
	We call $\Phi_{\mc L}(E_G)$ the {\it $\mc L$-- Atiyah class of $E_G$}. 
\end{proposition}

Let $\nabla_{\mathcal{L}} : V \to \At_{\delta}(E_G)$ be a $\mc L$-- connection 
on $E_G$ over $X$. For all locally defined sections $s$ and $t$ of $V$, let 
$$\kappa_{\nabla_{\mathcal{L}}}(s, t) := [\nabla_{\mathcal{L}}(s), \nabla_{\mathcal{L}}(t)] - \nabla_{\mathcal{L}}([s, t]).$$ 
Since the homomorphism $\widetilde{\rho} : \At_{\delta}(E_G) \to V$ respects the Lie algebra 
structures on the sheaves of sections, $\kappa_{\nabla_{\mathcal{L}}}(s, t)$ defines a holomorphic local 
section of $\ad(E_G)$. Thus we obtain a section 
$$\kappa_{\nabla_{\mathcal{L}}} \in H^0(X, \ad(E_G)\otimes \Omega^{2}_{\mathcal{L}}),$$
called the {\it curvature} of the $\mc L$--connection 
$\nabla_{\mathcal{L}}$ on $E_G$. The section $\kappa_{\nabla_{\mathcal{L}}}$ can be considered as an obstruction 
for $\nabla_{\mathcal{L}}$ to be a Lie algebra homomorphism. 

\begin{defn}
	An $\mc L$-- connection $\nabla_{\mathcal{L}}$ on a 
	principal $G$--bundle $E_G$ on $X$ is said to be {\it integrable} if $\kappa_{\nabla_{\mathcal{L}}} = 0$.
\end{defn}



Let $X$ be an irreducible smooth projective algebraic curve over $\mathbb{C}$, and $G$ be a complex reductive affine algebraic group. Let $E_{G}$ be a principal $G$--bundle over $X$. Let $\chi : L(P) \longrightarrow \mathbb{G}_{m} = \mathbb{C}^{*}$
be a character of a Levi subgroup $L(P)$ of a parabolic subgroup $P \subset G$. Let $ E_{L(P)} \subset E_{G}$ be a reduction of structure group of $E_{G}$ to $L(P)$. By extending the structure group of $E_{L(P)}$ via the character $\chi$, we obtain a principal $\mathbb{C}^{*}$--bundle $E_{L(P)} \times^{L(P)} \mathbb{C}^{*}$ over $X$. Using the standard multiplication action of $\mathbb{C}^{*}$ on $\mathbb{C}$, this principal $\mathbb{C}^{*}$--bundle defines a line bundle $\mathscr{L}(E_{L(P)}, \chi) \longrightarrow X$.

The principal $G$--bundle $E_{G}$ admits an algebraic $\mathcal{T}_{X}$--connection if and only if, for every triple $(P, L(P), \chi)$ as above, and every reduction of structure group $E_{L(P)}$ of $E_{G}$ to $L(P)$, one has
\[
\deg \bigl(\mathscr{L}(E_{L(P)}, \chi)\bigr)=0
\]
(see \cite{Azad-Biswas-2002}).

Let $F_{G}$ be a principal $G$--bundles on $X$. Recently, \cite[Theorem 6.1]{Biswas-2026} generalized this criterion to the setting of $\mathcal{L}$--connections. A direct consequence of this theorem and example 2.1.2 \eqref{example:Atiyah Lie algebroid of a principal bundle} is stated below.

\begin{proposition}
	 Assume that $F_{G}$ admits an $\mathcal{T}_{X}$--connection, and consider the Atiyah Lie algebroid $(\text{At}(F_{G}), [\cdot,\cdot],\rho)$. Let $E_{G}$ be a principal $G$-- bundle on $X$. Then the following three statements are equivalent:
	\begin{enumerate}
		\item For every triple $(P, L(P), \chi)$, and every algebraic reduction of structure group $E_{L(P)}$ of
		$E_{G}$ to $L(P)$, $\deg(\mathscr{L}(E_{L(P)}, \chi)) = 0$.
		\item $E_{G}$ admits a Lie algebroid $G$--connection for $(\text{At}(F_{G}),[\cdot,\cdot],\rho)$.
		\item $E_{G}$ admits an integrable Lie algebroid $G$--connection for $(\text{At}(F_{G}),[\cdot,\cdot],\rho)$. 
	\end{enumerate}
\end{proposition}

\begin{proposition}
	Assume that $F_{G}$ does not admit any $\mathcal{T}_{X}$--connection, and
	consider the Atiyah Lie algebroid $(\text{At}(F_{G}), [\cdot,\cdot],\rho)$ associated to $F_{G}$. Then any principal $G$--bundles bundle $E_{G}$ on $X$ admits a Lie algebroid connection for $(\text{At}(F_{G}), [\cdot,\cdot],\rho)$.
\end{proposition}






%%%%%%%%%%%%%%%%%%%%%%%%%%%%%%%%%%%%%%%%%%%%%%%%%%%%%%%%%%%%%%%%%%%%%%%%%%%%%%%%%%%%%%%%%%%%%%%%%%%%%%%%%%%%%%%%%%%%%%%%%%%%%%%%%%%%%%%%%%%%%%%%%%%%%%%%%%%%%%%%

\subsection{Characterization of integrable transitive algebraic Lie algebroids}.
Let $X$ be an irreducible smooth projective algebraic variety over $\mathbb{C}$ and let $(\mathcal{A},[\ ,\ ],\rho)$ be an algebraic Lie algebroid on $X$.
Denote by $X^{an}$ the associated complex analytic space.

Recall that analytification defines a functor from algebraic coherent
sheaves on $X$ to coherent analytic sheaves on $X^{an}$. In particular,
the vector bundle $\mathcal{A}$ gives rise to a holomorphic vector bundle
$\mathcal{A}^{an}$ on $X^{an}$. The anchor morphism
\[
\rho : \mathcal{A} \longrightarrow T_X
\]
analytifies to a holomorphic morphism
\[
\rho^{an} : \mathcal{A}^{an} \longrightarrow T_{X^{an}} .
\]
Similarly, the Lie bracket
\[
[\ ,\ ] : \mathcal{A} \otimes_{\mathcal{O}_X} \mathcal{A}
\longrightarrow \mathcal{A}
\]
induces a holomorphic bracket
\[
[\ ,\ ]^{an} : \mathcal{A}^{an} \otimes_{\mathcal{O}_{X^{an}}}
\mathcal{A}^{an} \longrightarrow \mathcal{A}^{an}.
\]
These structures satisfy the Leibniz rule and Jacobi identity after
analytification. Hence we obtain a holomorphic Lie algebroid
\[
(\mathcal{A}^{an},[\ ,\ ]^{an},\rho^{an})
\]
on $X^{an}$, called the \emph{analytic Lie algebroid associated to
	$\mathcal{A}$}.

There is a standard construction that associates with a Lie groupoid $\mathcal{G}$ a Lie algebroid $\mathcal{L}(\mathcal{G})$. A Lie algebroid $\mathcal{A}$ on $X$ (in the smooth, holomorphic, or algebraic category) is said to be \textit{integrable} if there exists a Lie groupoid $\mathcal{G} \rightrightarrows X$ in the same category whose associated Lie algebroid is isomorphic to $\mathcal{A}$. Lie algebroids are not always integrable, but for $\mathcal{L}$ an integrable Lie algebroid, there exists a unique $s$-connected $s$--simply connected Lie groupoid $\tilde{\mathcal{G}}$ integrating $\mathcal{L}$.

\begin{example}
Let $\pi:E_{G}\rightarrow X$ be a princpal $G$ bundle. Consider $\mathfrak{G}(E_{G}):=\frac{E_{G}\times E_{G}}{G}$, with $G$ acting diagonally on the product. The product $[a,b]\cdot[b',c] = [a,cg]$, where $g\in G$ is such that $b'=bg^{-1}$, defines a Lie groupoid structure on $\mathfrak{G}(E_{G})$, with source map $s([a,b])=\pi(b)$ and target map $t([a,b]) = \pi(a)$, called the \textit{gauge groupoid} of $E_{G}$. The Lie algebroid associated with $\mathfrak{G}(E_{G})$ is $\text{At}(E_{G})$. Therefore the Atyah algebroid $\text{At}(E_{G})$ is integrable.
\end{example}

\begin{lem}\label{lemma:-analytic-integrability-variety}
	Let $X$ be a smooth projective algebraic variety over $\mathbb{C}$ and let
	$(\mathcal{A},[\ ,\ ],\rho)$ be an algebraic Lie algebroid on $X$.
	Denote by $(\mathcal{A}^{an},[\ ,\ ]^{an},\rho^{an})$ the associated
	holomorphic Lie algebroid on the complex analytic space $X^{an}$.
	Then $\mathcal{A}$ is algebraically integrable if and only if
	$\mathcal{A}^{an}$ is analytically integrable.
\end{lem}

\begin{proof}
	Assume first that $\mathcal{A}$ is algebraically integrable. Then there
	exists an algebraic Lie groupoid $\mathcal{G} \rightrightarrows X$
	whose Lie algebroid is isomorphic to $\mathcal{A}$. Passing to the
	associated analytic spaces, we obtain a holomorphic Lie groupoid $
	\mathcal{G}^{an} \rightrightarrows X^{an}$.
	The Lie algebroid of $\mathcal{G}^{an}$ identifies naturally with
	$\mathcal{A}^{an}$, hence $\mathcal{A}^{an}$ is analytically integrable.
	
	Conversely, assume that $\mathcal{A}^{an}$ is analytically integrable.
	Then there exists a holomorphic Lie groupoid $\mathcal{G}^{an} \rightrightarrows X^{an}$ integrating $\mathcal{A}^{an}$. Replacing $\mathcal{G}^{an}$ by the source simply connected integration if necessary, we may assume that $\mathcal{G}^{an}$ is source simply connected.
	
	Since $X$ is a smooth projective algebraic variety, $X^{an}$ is compact.
	The structure maps of the Lie groupoid
	\[
	s,t : \mathcal{G}^{an} \longrightarrow X^{an}, \qquad
	m : \mathcal{G}^{an}\,{}_{s}\!\times_{t}\mathcal{G}^{an}
	\longrightarrow \mathcal{G}^{an},
	\]
	together with the inversion and unit maps, are holomorphic. The graph of
	these morphisms is closed in the corresponding fiber products, and hence
	defines coherent analytic subspaces. By the GAGA principle for projective
	varieties, coherent analytic sheaves and closed analytic subspaces over
	$X^{an}$ correspond to algebraic ones over $X$.
	
	Therefore there exists an algebraic space $\mathcal{G}$ whose
	analytification is $\mathcal{G}^{an}$, and the groupoid structure maps
	are induced by algebraic morphisms. Consequently, $\mathcal{G} \rightrightarrows X$ is an algebraic Lie groupoid.Finally, the Lie algebroid of $\mathcal{G}$ coincides with $\mathcal{A}$,
	since analytification preserves the Lie algebroid construction. Hence
	$\mathcal{A}$ is algebraically integrable.This completes the proof.
\end{proof}

\begin{proposition}\label{prop:transitive-integrable-atiyah}
	Let $X$ be a smooth connected manifold and let $\mathcal{A}$ be a transitive smooth Lie algebroid on $X$. Then $\mathcal{A}$ is integrable if and only if $\mathcal{A}$ is isomorphic to the Atiyah algebroid of a smooth principal $G$--bundle over $X$ for some Lie group $G$.
\end{proposition}

\begin{proof}
	This is a classical result; see, for example, \cite[Corollary 6.4]{Moerdijk-2003}.
\end{proof}

It is known that the integrability of a holomorphic Lie algebroid is
equivalent to the integrability of its underlying smooth real Lie
algebroid (see \cite[Theorem 3.17]{GSX-2009}). Hence by proposition \eqref{prop:transitive-integrable-atiyah} we obtain the
following corollary.
\begin{corollary}\label{cor:classifcation of integrable transitive-Lie algebroid}
	Let $X$ be a smooth irreducible projective algebraic variety over $\mathbb{C}$ and let $\mathcal{A}$ be a transitive algebraic Lie algebroid on $X$. Then $\mathcal{A}$ is algebraically integrable if and only if
	$\mathcal{A}$ is isomorphic to the Atiyah algebroid of a principal
	$G$--bundle on $X$, for some algebraic group $G$.
\end{corollary}





%%%%%%%%%%%%%%%%%%%%%%%%%%%%%%%%%%%%%%%%%%%%%%%%%%%%%%%%%%%%%%%%%%%%%%%%%%%%%%%%%%%%%%%%%%%%%%%%%%%%%%%%%%%%%%%%%%%%%%%%%%%%%%%%%%%%%%%%%%%%%%%%%%%%%%%%%%%%%%%%





\subsection{$\Lambda$-modules and Moduli spaces of $\mathcal{L}$-connections}
Let $X$ be a smooth projective variety over $\mathbb{C}$. Let us recall some definitions from (\cite[\S\,1]{Beilinson-Bernstein-1993},\cite[\S\,2.1]{FT-2017},\cite[\S\,2]{CT.Simpson-1994}). A differential $\mathcal{O}_{X}$-\textit{bimodule} is a quasicoherent sheaf on $X \times X$ supported on the diagonal $\Delta(X)\subset X \times X$. We can regard a differential $\mathcal{O}_{X}$--bimodule as a sheaf of $\mathcal{O}_{X}$--bimodules over $X$. An $\mathcal{O}_{X}$--\textit{differential algebra} or simply $D$-\textit{algebra} on $X$ is a sheaf of associative algebras (not necessarily commutative) $\Lambda$ on the Zariski topology of $X$ equipped with a morphism
of algebras $i : \mathcal{O}_{X} \rightarrow \Lambda$ such that $\Lambda$ is differential $\mathcal{O}_{X}$--bimodule. This implies that $\Lambda$ comes with an increasing filtration
\[0=\Lambda_{-1}\subset \Lambda_{0}\subset\Lambda_{1}\subset \Lambda_{2}\dots\]
such that $\Lambda =\cup_{i}\Lambda_{i}$ and for any $f$ in $\mathcal{O}_{X}$ and $\lambda \in \Lambda_{i}$ one has $f\cdot \lambda-\lambda\cdot f \in \Lambda_{i-1}$. We denote
$\text{Gr}_{i}\Lambda=\Lambda_{i}/\Lambda_{i-1}$ and $Gr_{\bullet}\Lambda:= \bigoplus\text{Gr}_{i}\Lambda$. Recalling that $\Lambda$ is a differential $\mathcal{O}_{X}$--bimodule,
we denote by $\mathscr{I}(\Lambda)$ the associated quasicoherent sheaf on $X \times X$ supported on the diagonal.

We will focus on $D$-algebras that are almost polynomial (cf. \cite[\S\,2, page 81]{CT.Simpson-1994}), namely those $D$-algebras $\Lambda$ such that $\Lambda_{0}=\mathcal{O}_{X}$, $\text{Gr}_{1}(\Lambda)$ is a locally free $\mathcal{O}_{X}$-module and whose associated graded algebra is isomorphic to the symmetric product over the first graded piece, $\text{Gr}_{\bullet}\Lambda = \text{Sym}_{\mathcal{O}_{X}}^{\bullet}(\text{Gr}_{1}(\Lambda))$.

Almost polynomial $D$-algebras may be described in terms of Lie algebroids. The relation between Lie algebroids and D-algebras is stated in the following lemma:

\begin{lem}\cite[Theorem 34]{P.Tortella-2011T}
Let $X$ be a smooth projective variety over $\mathbb{C}$, and let $V$ be a locally free $\mathcal{O}_{X}$-module of finite rank. There is a bijective correspondence between isomorphism classes of:
\begin{enumerate}
\item pairs $(\Lambda,\Xi)$, with $\Lambda$ an almost polynomial $D$-algebra and $\Xi$ an isomorphism of the associated graded algebra $\text{Gr}_{\bullet}\Lambda$ with the symmetric algebra $\text{Sym}_{\mathcal{O}_{X}}^{\bullet}(\text{Gr}_{1}(\Lambda))$.
\item pairs $(\mathcal{L},\tilde{\mathcal{L}})$, with $\mathcal{L}$ a Lie algebroid structure on $V$ and $\tilde{\mathcal{L}}$ a central extension of $\mathcal{L}$ by $\mathcal{O}_{X}$.
\item pairs $(\mathcal{L}, \Sigma)$, with $\mathcal{L}$ a Lie algebroid structure on $V$ and $\Sigma \in \mathbb{H}^{2}(X, \tau^{\geqslant 1}\Omega^{\bullet}_{\mathcal{L}})$.
\end{enumerate}
\end{lem}

A \emph{split almost polynomial $D$--algebra} is an almost polynomial $D$--algebra $\Lambda$ together with a morphism
\[
\zeta:\operatorname{Gr}_{1}(\Lambda)\longrightarrow \Lambda_{1}
\]
of left $\mathcal{O}_{X}$--modules splitting the natural projection $\Lambda_{1}\to \operatorname{Gr}_{1}(\Lambda)$ (cf. \cite[\S\,2, page 81]{CT.Simpson-1994}). One can check that universal enveloping $D$-algebra (cf. \cite[\S\,1]{Beilinson-Bernstein-1993}) associated to Lie algebroid $\mathcal{L}$ is a split almost polynomial $D$--algebra.

\begin{lem}\label{Lie algebroid is-equivalent-to-Lambda_L}\cite[\S\,3, Theorem  3.12]{Alfaya-Oliveire-2024}
	The following correspondences gives inverse equivalence of following categories:
	
	$$
	\left\{
	\begin{array}{c}
		isomorphism \ classes \ of \\ Lie \ algebroids \ on \ X
	\end{array}
	\right\}
	\Longleftrightarrow 
	\left\{
	\begin{array}{c}
		isomorphism \ classes \ of \\ split \ almost \ polynomial \  D-algebra \ on \ X
	\end{array}
	\right\}
	$$
	$$\mathcal{L}\mapsto \Lambda_{\mathcal{L}}(\text{universal envoloping $D$--algebra of $\mathcal{L})$}.$$
\end{lem}

An important class of algebras are constructed from Lie algebroids supported on the tangent bundle. In the untwisted case, we have of course the \textit{algebra of differential operators}, or \textit{De Rham D-algebra} in Simpson’s notation (\cite[\S\,6]{CT.Simpson-1995},\cite[\S\,2]{CT.Simpson-1994}) $\mathscr{D}_{X}=\Lambda^{DR}$ which arises as the universal enveloping algebra of the canonical Lie algebroid $(\mathcal{T}_{X}, [\cdot, \cdot],\text{Id}_{\mathcal{T}_{X}})$, i.e. the Lie algebroid supported on $\mathcal{T}_{X}$ obtained after setting the anchor to be the identity morphism. The abelianization of $\Lambda^{\text{DR}}$ is the \textit{Dolbeault $D$--algebra}, $\Lambda^{\text{Dol}}:= \text{Gr}(\Lambda^{\text{DR}})=\text{Sym}^{\bullet}(\mathcal{T}_{X})$,
which can be obtained as the universal enveloping algebra of the trivial Lie algebroid supported on the tangent bundle, $(\mathcal{T}_{X}, [\cdot, \cdot],0)$, where the anchor is the $0$ map. We can construct also a family of $D$--algebras which is a deformation from $\Lambda^{\text{DR}}$ to $\Lambda^{\text{Dol}}$. Set, for each $t \in \mathbb{C}$, the Lie algebroid $(\mathcal{T}_{X},[\cdot,\cdot],t\cdot \Id_{\mathcal{T}_{X}})$, where the anchor consists on scaling by $t$, and define $\Lambda^{t}$ to be the universal enveloping $D$ algebra of it. Similerly, for a given Lie algebroid $\mathcal{L}=(V,[\cdot,\cdot],\delta)$, we can construct \textit{$\mathcal{L}$--twisted De Rham $D$--algebra} $\Lambda^{DR}_{\mathcal{L}}$ (universal enveloping algebra of the Lie algebroid $\mathcal{L}$) and \textit{$\mathcal{L}$--twisted Dolbeault $D$--algebra}, $\Lambda_{\mathcal{L}}^{\text{Dol}}:= \text{Gr}(\Lambda_{\mathcal{L}}^{\text{DR}})=\text{Sym}^{\bullet}(V)$, which can be obtained as the universal enveloping algebra of the trivial Lie algebroid $(V,[\cdot,\cdot],0)$.  We can construct also a family of $D$--algebras which is a deformation from $\Lambda_{\mathcal{L}}^{\text{DR}}$ to $\Lambda_{\mathcal{L}}^{\text{Dol}}$.


\begin{defn}
Let $\Lambda$ be a split almost polynomial $D$ algebra. Let $\mathcal{E}$ be a coherent sheaf on $X$. A $\Lambda$--\textit{module structure on $\mathcal{E}$} is a $\mathcal{O}_{X}$-morphism $\varphi : \Lambda \otimes \mathcal{E} \rightarrow \mathcal{E}$ satisfying the usual module axioms and such that the $\mathcal{O}_{X}$-module structure
on $\mathcal{E}$ induced by $\mathcal{O}_{X} \rightarrow \Lambda$ coincides with the original one.
\end{defn}

Therefore by lemma \eqref{Lie algebroid is-equivalent-to-Lambda_L} we have the following:
\begin{lem}\label{Lambda_L module is equvalent to $L$-connection }\cite[\S\,2, lemma 2.13]{CT.Simpson-1994}
	Let $\mathcal{L}$ be a Lie algebroid on $X$ and let $\mathcal{E}$ be a vector bundle on $X$.
	Then giving a $\Lambda^{\text{DR}}_{\mathcal{L}}$--module structure on $\mathcal{E}$ is equivalent to the choice of an integrable $\mathcal{L}$--connection
	\[
	D_{\mathcal{L}}:\mathcal{E}\longrightarrow \mathcal{E}\otimes \Omega^{1}_{\mathcal{L}}
	\]
	on $\mathcal{E}$.
\end{lem}

A $\Lambda$--module $\mathcal{E}$ is said to be of \emph{pure dimension $d$} if the underlying $\mathcal{O}_{X}$--coherent sheaf is of pure dimension $d$. 
The \emph{Hilbert polynomial}, the \emph{rank}, and the \emph{slope} of $\mathcal{E}$ are defined to be those of the underlying coherent sheaf (cf.~\cite[\S\,1]{CT.Simpson-1994}, \cite[Ch.~1]{Huybrechts-Lehn-2010}).

A $\Lambda$--module $\mathcal{E}$ is called \emph{$\mathcal{P}$--semistable} (resp.~\emph{$\mathcal{P}$--stable}) if it is of pure dimension and for every nonzero proper $\Lambda$--submodule
$\mathscr{F}\subset \mathcal{E}$ with
\[
0 < r(\mathscr{F}) < r(\mathcal{E}),
\]
there exists an integer $N$ such that
\[
\frac{\mathcal{P}(\mathscr{F},n)}{r(\mathscr{F})}
\leqslant
\frac{\mathcal{P}(\mathcal{E},n)}{r(\mathcal{E})}
\quad
(\text{resp. } <)
\]
for all $n \geqslant N$.

A $\Lambda$--module $\mathcal{E}$ is called \emph{$\mu$--semistable} (resp.~\emph{$\mu$--stable}) if it is of pure dimension and for every nonzero proper $\Lambda$--submodule
$\mathscr{F}\subset \mathcal{E}$ with
\[
0 < r(\mathscr{F}) < r(\mathcal{E}),
\]
we have
\[
\mu(\mathscr{F}) \leqslant \mu(\mathcal{E})
\quad
(\text{resp. } <).
\]
Moreover, $\mathcal{P}$--semistability implies $\mu$--semistability, while $\mu$--stability implies $\mathcal{P}$--stability (cf.~\cite[\S\,1]{CT.Simpson-1994}).

\medskip
\noindent
\textbf{Moduli functor.}
Let $S$ be a scheme over $\mathbb{C}$, and let $p_{S}, p_{X} : S \times X \longrightarrow S, X$
denote the natural projections. Set $\Lambda^{S} := p_{X}^{*}\Lambda$.
We define a contravariant functor
\[
\mathcal{F}_{\Lambda}(\mathcal{P}) : \texttt{Sch}_{/\mathbb{C}} \longrightarrow \texttt{Set}
\]
as follows.

\begin{itemize}
	\item To a scheme $S$, the functor associates the set of isomorphism classes of $\Lambda^{S}$--modules $(\mathcal{E},\mu)$ such that:
	\begin{itemize}
		\item $\mathcal{E}$ is flat over $S$, and
		\item for every closed point $s \in S$, the fiber $(\mathcal{E}_{s},\mu_{s})$ is a semistable $\Lambda$--module on $X$ with Hilbert polynomial $\mathcal{P}$.
	\end{itemize}
	Two such families $(\mathcal{E},\mu)$ and $(\mathcal{E}',\mu')$ are identified if there exists a line bundle $L$ on $S$ such that
	\[
	\mathcal{E}' \cong \mathcal{E} \otimes p_{S}^{*}L,
	\qquad
	\mu' = \mu \otimes \mathrm{id}_{p_{S}^{*}L}.
	\]
	\item To a morphism of schemes $\psi : T \to S$, the functor associates the pullback family: if $(\mathcal{E},\mu)$ is a $\Lambda^{S}$--module flat over $S$, then
	$\psi^{*}(\mathcal{E},\mu)$ is a $\Lambda^{T}$--module flat over $T$.
\end{itemize}

\begin{thm}\label{moduli-of-Lambda-module}\cite[\S\,4, Th 4.7.]{CT.Simpson-1994}
Let $X$ be a smooth projective variety over $\mathbb{C}$, let $\Lambda$ be a sheaf of rings of differential operators on $X$, and let $\mathcal{P}$ be a numerical polynomial. 
Then there exists a quasi-projective variety $\mathcal{M}(\Lambda,\mathcal{P})$ which universally corepresents the moduli functor $\mathcal{F}_{\Lambda}(\mathcal{P})$. 
The closed points of $\mathcal{M}(\Lambda,\mathcal{P})$ are in one-to-one correspondence with Jordan equivalence classes of semistable $\Lambda$--modules on $X$ with Hilbert polynomial equal to $\mathcal{P}$.

\end{thm}


One may similarly define an $\mathcal{L}$--connection on a coherent sheaf (see Definition~\eqref{def:Lie-alg-conn-on-VB}). However, it is well known that a coherent sheaf equipped with a $\mathcal{T}_{X}$--connection is necessarily locally free, and moreover its $\mathcal{T}_{X}$--Chern classes vanish
(see \cite[\S\,VI, Proposition~1.7]{A.Borel-1987} and \cite[\S\,3, Theorem~4]{Atiyah-1957}). Consequently, by the \textit{Hirzebruch--Riemann--Roch Theorem} \cite{A. Moroianu-2007}, the normalized Hilbert polynomial of $\mathcal{E}$ coincides with that of $\mathcal{O}_{X}$.

In \cite[\S\,4]{R.L.Fernandez-2002}, R.L.Fernandez introduce Chern classes associated to Lie algebroids. In contrast, for a general Lie algebroid $\mathcal{L}$, a coherent sheaf may admit an $\mathcal{L}$--connection without being locally free (see \cite[\S\,2.2, Example~2.12]{GLP-2018}).

Another example is the following: let $\mathcal{L}$ be a holomorphic Lie algebroid, let $\mathcal{G}$ denote the associated holomorphic foliation, let $\mathcal{E}$ be a coherent $\mathcal{O}_{X}$--module, and let $D_{\mathcal{L}}$ be a holomorphic $\mathcal{L}$--connection on $\mathcal{E}$. Then one can only conclude that the restriction $\mathcal{E}|_{\mathfrak{G}}$ is locally free for every leaf $\mathfrak{G}$ of $\mathcal{G}$ (see \cite{R.L.Fernandez-2002}).



\begin{lem}\label{lem:transitive-Lie-algebroid-locally-free}
	Let $\mathcal{L}$ be a transitive Lie algebroid. Then any coherent sheaf $\mathcal{E}$ equipped with an $\mathcal{L}$--connection (not necessarily integrable) is locally free, and its $\mathcal{L}$--Chern classes vanish. Moreover, the normalized Hilbert polynomial of $\mathcal{E}$ coincides with that of $\mathcal{O}_{X}$.
\end{lem}

\begin{proof}
	Since $\mathcal{L}$ is a transitive Lie algebroid, it follows from \cite[\S\,2.2, Proposition~2.15]{GLP-2018} that $\mathcal{E}$ is locally free. Furthermore, the $\mathcal{L}$--Chern classes of $\mathcal{E}$ vanish by \cite[\S\,2, Theorem~18]{P.Tortella-2011T}.
	
	Again, since $\mathcal{L}$ is transitive, the vanishing of the $\mathcal{L}$--Chern classes implies the vanishing of the usual Chern classes (cf. \cite{Crainic-2009}). Hence, by the \textit{Hirzebruch--Riemann--Roch Theorem}, the normalized Hilbert polynomial of $\mathcal{E}$ coincides with that of $\mathcal{O}_{X}$.
\end{proof}

In view of the above discussion, we may assume that the pure dimension $d$ coincides with the dimension of $X$, and that the normalized Hilbert polynomial $\mathcal{P}_{0}$ is equal to that of $\mathcal{O}_{X}$; otherwise, the corresponding moduli spaces are empty.

Furthermore, any subsheaf of $\mathcal{E}$ preserved by the $\mathcal{L}$--connection is again a vector bundle equipped with an $\mathcal{L}$--connection, having the same normalized Hilbert polynomial $\mathcal{P}_{0}$. Therefore, for a fixed transitive Lie algebroid $\mathcal{L}$, every $\Lambda^{\mathrm{DR}}_{\mathcal{L}}$--module of rank $n$ is automatically $n\mathcal{P}_{0}$--semistable, and is $n\mathcal{P}_{0}$--stable.


The following theorem is an analogue of \cite[\S\,6, Theorem~6.13]{CT.Simpson-1995} in the setting of Lie algebroids.
\begin{thm}\label{eqn: moduli-of-L-connection on VB}
Let $\mathcal{L}$ be a transitive Lie algebroid and let $X$ be a smooth projective variety over $\mathbb{C}$. There exists a
quasi-projective scheme $\mathcal{M}^{\text{DR}}_{\mathcal{L}}(X,n)$ over $\mathbb{C}$ which universally corepresents the functor 
$\underline{\mathcal{M}^{\text{DR}}_{\mathcal{L}}}(X,n),$
assigning to a $\mathbb{C}$--scheme $S$ the set of isomorphism classes of
vector bundles $\mathcal{E}$ of rank $n$ on $X' := X \times S$
equipped with integrable $\mathcal{L}$--connections, and having Hilbert
polynomial $n\mathcal{P}_{0}$.
	
Fix a base point $\mathfrak{p} \in X$.  
Then there exists a quasi-projective scheme $\mathcal{R}^{\text{DR}}_{\mathcal{L}}(X,\mathfrak{p},n)$ over $\mathbb{C}$ representing the functor which assigns to a $\mathbb{C}$--scheme $S$ the set of isomorphism classes of triples $(\mathcal{E},D_{\mathcal{L}},\beta)$, where
\begin{itemize}
		\item $(\mathcal{E},D_{\mathcal{L}})$ is a vector bundle of rank $n$ on $X' = X \times S$ equipped with an integrable $\mathcal{L}$--connection, and
		\item $\beta : \mathcal{E}|_{\mathfrak{p}} \xrightarrow{\sim} \mathbb{C}^{n}$ is a framing at the point $\mathfrak{p}$.
\end{itemize}
	Moreover, with respect to a suitable linearization, all points of $\mathcal{R}^{\text{DR}}_{\mathcal{L}}(X,\mathfrak{p},n)$ are semistable for the natural action of $GL_{n}(\mathbb{C})$, and the resulting good quotient is naturally identified with $\mathcal{M}^{\text{DR}}_{\mathcal{L}}(X,n)$.
\end{thm}
\begin{proof}
Let $\Lambda^{\text{DR}}_{\mathcal{L}}$ be the $\mathcal{L}$--twisted de Rham $D$--algebra associated to the Lie algebroid $\mathcal{L}$. In Simpson’s description of $\Lambda^{\text{DR}}_{\mathcal{L}}$ as triplets $(H,\delta,\gamma)$ \cite[\S\,2, p.~82]{CT.Simpson-1994}, the sheaf $H$ is identified with
$\Omega^{1}_{\mathcal{L}}$, and its dual $H^{*}$ is the Lie algebroid
$\mathcal{L}$. The derivation is given by the anchor map
$\delta \colon \mathcal{L} \to \mathcal{T}_{X}$, and the bracket
$\{\,,\,\}_{\gamma}$ coincides with the Lie algebroid bracket on
$\mathcal{L}$.

The description of $\Lambda^{\text{DR}}_{\mathcal{L}}$--modules given in
\cite[\S\,2, Lemma~2.13]{CT.Simpson-1994} therefore coincides with the notion
of a coherent sheaf (or vector bundle) equipped with an integrable
$\mathcal{L}$--connection. If $\mathcal{E}$ is a vector bundle on $X$ endowed with an integrable $\mathcal{L}$--connection, then any subsheaf of $E$ preserved by the $\mathcal{L}$--connection has the same normalized Hilbert polynomial $\mathcal{P}_{0}$ (see, lemma \eqref{lem:transitive-Lie-algebroid-locally-free}). Consequently, $\mathcal{E}$ is $p$--semistable as a $\Lambda^{\text{DR}}_{\mathcal{L}}$--module.

Hence, the first assertion follows from
Theorem~\eqref{moduli-of-Lambda-module}, while the second assertion follows
from \cite[\S\,4, Theorem~4.10]{CT.Simpson-1994}.

\end{proof}
We refer to $\mathcal{M}^{\text{DR}}_{\mathcal{L}}(X,n)$ as the \emph{$\mathcal{L}$--twisted de Rham moduli space}, and to $\mathcal{R}^{\text{DR}}_{\mathcal{L}}(X,\mathfrak{p},n)$ as the \emph{$\mathcal{L}$--twisted de Rham representation space}.

%%%%%%%%%%%%%%%%%%%%%%%%%%%%%%%%%%%%%%%%%%%%%%%%%%%%%%%%%%%%%%%%%%%%%%%%%%%%%%%%%%%%%%%%%%%%%%%%%%%%%%%%%%%%%%%%%%%%%%%%%%%%%%%%%%%%%%%%%%%%%%%%%%%%%%%%%%%%%%%%%%%%%%%%%%%%%%%%%%%%%%%%%%%%%%%%%%%%%%%%%%%%%%%%%%%%%%%%%%%%%%%%%%%%%%%%%%%%%%%%

\subsection{The moduli of $\mathcal{L}$-twisted Higgs bundles and $\mathcal{L}$-Hodge moduli spaces}\label{subsec: L-Hodge modul space}
The usual notion of a Higgs bundle on $X$ admits a natural extension to the setting of Lie algebroids. Let $\mathcal{L} = (V,[\cdot,\cdot],\delta)$ be a Lie algebroid on $X$.

\begin{defn}\label{def:L-twisted-Higgs bundle}
An \textit{$\mathcal{L}$--twisted Higgs bundle} on $X$ is a pair $(\mathcal{E},\theta_{\mathcal{L}})$ consisting of an algebraic vector bundle $\mathcal{E}$ on $X$ together with an $\mathcal{O}_{X}$--linear morphism
\[
\theta_{\mathcal{L}} : \mathcal{E} \longrightarrow \mathcal{E} \otimes \Omega^{1}_{\mathcal{L}},
\]
called the \emph{$\mathcal{L}$--twisted Higgs field}, such that, 
$
\theta_{\mathcal{L}} \wedge \theta_{\mathcal{L}} = 0 \in H^{0}\!\left(X, \operatorname{End}(\mathcal{E})\otimes \Omega^{2}_{\mathcal{L}}\right).
$\end{defn}
Therefore, by \cite[lemma 2.13]{CT.Simpson-1994}, we have the following lemma,
\begin{lem}\label{lem: Lambda module structure--equivalent to L twisted Higgs bundle} Let $\mathcal{L}$ be a Lie algebroid on $X$ and let $\mathcal{E}$ be a coherent sheaf on $X$. Then giving a $\Lambda^{\text{Dol}}_{\mathcal{L}}$--module structure on $\mathcal{E}$ is equivalent to a $\mathcal{L}$--twisted Higgs bundle structure on $\mathcal{E}$. Also note that a $\Lambda^{\text{Dol}}_{\mathcal{L}}$--structure on $\mathcal{E}$ is equivalent to an integrable $\mathcal{L}_{0}=(V,[\cdot,\cdot],0)$--connection on $\mathcal{E}$.
\end{lem}

%The above definition \eqref{def:L-twisted-Higgs bundle} is equivalent to the notion of a $V^{*}$--twisted Higgs bundle introduced in \cite[\S\,2, Definition~2.1]{Alfaya-Oliveire-2024}.

Let $\mathcal{M}^{\mathrm{Higgs}^{\natural}}_{\mathcal{L}}$ denote the
functor which associates to a $\mathbb{C}$--scheme $S$ the set of
isomorphism classes of $p$--semistable $\mathcal{L}$--twisted Higgs
bundles (or, more generally, $\mathcal{L}$--twisted Higgs sheaves) on
$X' := X \times S$ with Hilbert polynomial $\mathcal{P}$. This functor is
universally corepresented by the moduli space
\[
\mathcal{M}^{\mathrm{Higgs}}_{\mathcal{L}}(X,\mathcal{P})
:= \mathcal{M}\!\left(\Lambda^{\mathrm{Dol}}_{\mathcal{L}}, \mathcal{P}\right),
\]
constructed in Theorem~\eqref{moduli-of-Lambda-module}, which is a
quasi-projective scheme over $\mathbb{C}$. The closed points of
$\mathcal{M}^{\mathrm{Higgs}}_{\mathcal{L}}(X,\mathcal{P})$ parametrize Jordan
equivalence classes of $p$--semistable $\mathcal{L}$--twisted Higgs
bundles with Hilbert polynomial $\mathcal{P}$.

Let $\mathcal{P}_{0}$ denote the Hilbert polynomial of $\mathcal{O}_{X}$.
Let $\mathcal{M}^{\mathrm{Dol}^{\natural}}_{\mathcal{L}}(X,n)$ denote the
functor which assigns to a $\mathbb{C}$--scheme $S$ the set of
isomorphism classes of $p$--semistable $\mathcal{L}$--twisted Higgs sheaves $\mathcal{E}$ on $X' := X \times S$ with Hilbert polynomial
$n\mathcal{P}_{0}$ and rank $n$, such that the $\mathcal{L}$--Chern classes vanishes for all closed point $s\in S$. 

This functor is universally corepresented by a scheme
$\mathcal{M}^{\mathrm{Dol}}_{\mathcal{L}}(X,n)$, which is a disjoint union
of certain connected components of
$\mathcal{M}^{\mathrm{Higgs}}_{\mathcal{L}}(X,n\mathcal{P}_{0})$ (the fact that $\mathcal{M}^{\text{Dol}}_{\mathcal{L}}(X,n)$ may be proper subset of $\mathcal{M}^{\text{Higgs}}_{\mathcal{L}}(X,n\mathcal{P}_{0})$. The points of $\mathcal{M}^{\mathrm{Dol}}_{\mathcal{L}}(X,n)$ correspond to Jordan equivalence classes of $p$--semistable torsion free $\mathcal{L}$-- twisted torsion free Higgs sheves of rank $n$ on $X$ with $\mathcal{L}$--Chern classes vanishes. There is an open set $\mathcal{M}^{s\text{Dol}}_{\mathcal{L}}(X,n)$ parameterizing $p$--stable $\mathcal{L}$--twisted Higgs sheaves.



\begin{proposition}\label{prop:locally-free-L-Higgs}
	Let $X$ be a smooth projective variety and let $
	\mathcal L 
	\xrightarrow{\ \delta\ } \mathcal{T}_X \longrightarrow 0$,
	be a transitive Lie algebroid. Let $(\mathcal{E},\theta_{\mathcal L})$ be a $\mu$--semistable
	torsion--free $\mathcal L$--twisted Higgs sheaf with vanishing
	$\mathcal L$--Chern classes. Then $\mathcal{E}$ is locally free and is in fact an extension of $\mu$--stable Higgs bundles whose $\mathcal{L}$--Chern classes vanish. Any sub $\mathcal{L}$--twisted Higgs sheaf  of degree zero is a strict sub bundle with vanishing $\mathcal{L}$--Chern classes.
\end{proposition}

\begin{proof}
	Since $\mathcal L$ is transitive, the anchor map
	$\delta : \mathcal L \to \mathcal{T}_X$ is surjective. As the question is local on
	$X$, we may choose an $\mathcal O_X$--linear splitting
	$b : \mathcal{T}_X \to \mathcal L$ such that $\delta \circ b = \mathrm{id}_{\mathcal{T}_X}$.
	Dualizing, we obtain an morphism
	\[
	b^\vee : \Omega^1_{\mathcal L} \longrightarrow \Omega^1_X .
	\]
	
	Composing the $\mathcal L$--Higgs field
	$\theta_{\mathcal L} : \mathcal{E} \to \mathcal{E} \otimes \Omega^1_{\mathcal L}$
	with $\mathrm{id}_\mathcal{E} \otimes b^\vee$, we define
	\[
	\theta_X := (\mathrm{id}_\mathcal{E} \otimes b^\vee)\circ \theta_{\mathcal L}
	: \mathcal{E} \longrightarrow \mathcal{E} \otimes \Omega^1_X .
	\]
	
	The dual splitting $b^\vee:\Omega^1_{\mathcal L}\to\Omega^1_X$ induces a
	morphism of graded algebras
	$\wedge^\bullet\Omega^1_{\mathcal L}\to\wedge^\bullet\Omega^1_X$.
	Hence the induced Higgs field
	$\theta_X=(\mathrm{id}_\mathcal{E}\otimes b^\vee)\circ\theta_{\mathcal L}$
	satisfies
	\[
	\theta_X\wedge\theta_X
	=(\mathrm{id}_\mathcal{E}\otimes\wedge^2 b^\vee)
	(\theta_{\mathcal L}\wedge\theta_{\mathcal L})=0,
	\]
	Hence $\theta_X$ is integrable, and $(\mathcal{E},\theta_X)$ is a Higgs sheaf in the usual sense.
	
	
	Let $\mathcal{F} \subset \mathcal{E}$ be a $\theta_X$--invariant subsheaf.
	Since $\rho$ is surjective, $\mathcal{F}$ is also invariant under
	$\theta_{\mathcal L}$. Therefore the $\mu$--semistability of
	$(\mathcal{E},\theta_{\mathcal L})$ implies that $(\mathcal{E},\theta_X)$ is
	$\mu$--semistable as a Higgs sheaf.
	
	For a transitive Lie algebroid, vanishing of the $\mathcal L$--Chern
	classes implies vanishing of the usual Chern classes, that is
	$c_i(\mathcal{E})=0$ for all $i$.
	Thus $(\mathcal{E},\theta_X)$ is a $\mu$--semistable torsion--free Higgs sheaf
	with trivial Chern classes.
	
	By \cite[\S\,6, Proposition~6.6]{CT.Simpson-1995}, any $\mu$--semistable
	torsion--free Higgs sheaf on a smooth projective variety with vanishing
	Chern classes is locally free. Applying this result to the Higgs sheaf
	$(\mathcal{E},\theta_X)$, we conclude that $\mathcal{E}$ is locally free.
	
	
\end{proof}

\begin{corollary}
If $X$ is smooth and projective over $\mathbb{C}$, If $S$ is an $\mathbb{C}$--scheme, and if $\mathcal{E}$ is an element of $\mathcal{M}^{\text{Dol}^{\sharp}}_{\mathcal{L}}(X,n)(S)$, then $\mathcal{E}$ is locally free over $X'=X\times S$. The point of $\mathcal{M}^{\text{Dol}}_{\mathcal{L}}(X,n)$ correspond to sums of $\mathcal{L}$--twisted $\mu$--stable Higgs bundles with vanishing rational $\mathcal{L}$--Chern classes on $X$.
\end{corollary}

\begin{rem}
For $\mathcal{L}$--twisted Higgs sheaves with vanishing $\mathcal{L}$--Chern classes, $p$--semi-stablity (resp. $p$--stability) is equivalent to $\mu$--semistability (resp. $\mu$--stability). This is follows from, proposition \eqref{prop:locally-free-L-Higgs}.
\end{rem}

Suppose $X$ is smooth and projective over $\mathbb{C}$. 
An $\mathcal{L}$--twisted Higgs sheaves $\mathcal{E}$ on $X' := X \times S$, flat over $S$, is said to be of \emph{semiharmonic type} if, for every closed point $s \in S$, the restriction $\mathcal{E}_{s}$ is a $p$--semistable
$\mathcal{L}$--twisted Higgs bundle with vanishing rational
$\mathcal{L}$--Chern classes. It is said to be of \emph{harmonic type} if each $E_{s}$ is a direct sum of
stable $\mathcal{L}$--twisted Higgs sheaves with vanishing rational
$\mathcal{L}$--Chern classes (cf.~\cite[\S\,6, p.~17]{CT.Simpson-1995}).


Hence, we obtain the following theorem, which extends the construction of
\cite[\S\,6, p.~17]{CT.Simpson-1995} to the setting of
$\mathcal{L}$--twisted Higgs bundles.


\begin{thm}\label{eqn: moduli-of-L- twisted Higgs bundles}
	There exists a quasi-projective scheme $\mathcal{M}^{\mathrm{Dol}}_{\mathcal{L}}(X,n)$ over $\mathbb{C}$ which universally corepresents the functor
	$\mathcal{M}^{\mathrm{Dol}^{\natural}}_{\mathcal{L}}(X,n)$ assigning to a $\mathbb{C}$--scheme $S$ the set of isomorphism classes of $\mathcal{L}$--twisted Higgs bundles of semiharmonic type on $X' := X \times S$, with Hilbert polynomial $n\mathcal{P}_{0}$ and rank $n$.
	
	Fix a base point $\mathfrak{p} \in X$. 
	Then there exists a quasi-projective scheme
	$\mathcal{R}^{\mathrm{Dol}}_{\mathcal{L}}(X,\mathfrak{p},n)$ over $\mathbb{C}$ representing the functor which assigns to a $\mathbb{C}$--scheme $S$ the set of isomorphism classes of triples $(\mathcal{E},\theta_{\mathcal{L}},\beta)$, where
	\begin{itemize}
		\item $(\mathcal{E},\theta_{\mathcal{L}})$ is an $\mathcal{L}$--twisted Higgs bundle of semiharmonic type on $X' := X \times S$, and
		\item $\beta : \mathcal{E}|_{\mathfrak{p}} \xrightarrow{\sim} \mathbb{C}^{n}$ is a framing at $\mathfrak{p}$.
	\end{itemize}
	Moreover, with respect to a suitable linearization, all points of
	$\mathcal{R}^{\mathrm{Dol}}_{\mathcal{L}}(X,\mathfrak{p},n)$ are semistable for the natural action of $GL_{n}(\mathbb{C})$, and the resulting good quotient is naturally identified with
	$\mathcal{M}^{\mathrm{Dol}}_{\mathcal{L}}(X,n)$.
\end{thm}

\begin{proof}
	The first assertion follows from
	\eqref{lem: Lambda module structure--equivalent to L twisted Higgs bundle} together with Theorem~\eqref{moduli-of-Lambda-module} and the preceding discussion.
	
	Fix a base point $\mathfrak{p} \in X$, and let
	$\mathcal{R}(\Lambda^{\text{Dol}}_{\mathcal{L}},\mathfrak{p},n)$ denote the representation space
	of framed $\Lambda^{\mathrm{Dol}}_{\mathcal{L}}$--modules constructed in
	\cite[\S\,4, Theorem~4.10]{CT.Simpson-1994}. By proposition \eqref{prop:locally-free-L-Higgs} all $p$--semistable $\mathcal{L}$--twisted sheaves with vanishing rational $\mathcal{L}$--Chern classes satisfy condition $\text{LF}(\xi_{\mathfrak{p}})$, where $\xi_{\mathfrak{p}}:\text{Spec}(\mathbb{C})\rightarrow X, \mathfrak{m}\mapsto \mathfrak{p}$. Let
	$\mathcal{R}^{\mathrm{Dol}}_{\mathcal{L}}(X,\mathfrak{p},n)$ denote the disjoint union of those connected components of
	$\mathcal{R}(\Lambda^{\text{Dol}}_{\mathcal{L}},\mathfrak{p},n)$ corresponding to
	$\mathcal{L}$--twisted Higgs sheaves with vanishing rational
	$\mathcal{L}$--Chern classes.
	
	Then $\mathcal{R}^{\mathrm{Dol}}_{\mathcal{L}}(X,\mathfrak{p},n)$ represents the functor which assigns to a $\mathbb{C}$--scheme $S$ the set of isomorphism classes of triples $(\mathcal{E},\theta_{\mathcal{L}},\beta)$, where
	$(\mathcal{E},\theta_{\mathcal{L}})$ is an $\mathcal{L}$--twisted Higgs bundle of semiharmonic type on $X \times S$ and
	$\beta \colon \mathcal{E}|_{\mathfrak{p}} \xrightarrow{\sim} \mathbb{C}^{n}$ is a framing.
	
	
\end{proof}

We refer to $\mathcal{M}^{\mathrm{Dol}}_{\mathcal{L}}(X,n)$ as the
\emph{$\mathcal{L}$--twisted Dolbeault moduli space}, and to
$\mathcal{R}^{\mathrm{Dol}}_{\mathcal{L}}(X,\mathfrak{p},n)$ as the
\emph{$\mathcal{L}$--twisted Dolbeault representation space}.


Again, from remark after lemma \eqref{lem: Lambda module structure--equivalent to L twisted Higgs bundle} one can associate to $\Lambda^{\text{DR}}_{\mathcal{L}}$ a $D$--algebra
$\Lambda_{\mathcal{L}}^{\mathrm{red}}$ on $X \times \mathbf{A}^{1}$,
flat over $\mathbf{A}^{1}$, whose fiber over $1$ is $\Lambda^{\text{DR}}_{\mathcal{L}}$
and whose fiber over $0$ is isomorphic to its associated graded algebra
\[
\Lambda^{\mathrm{Dol}}_{\mathcal{L}}
= \operatorname{Gr}_{\bullet}(\Lambda^{\text{DR}}_{\mathcal{L}})
\cong \operatorname{Sym}^{\bullet}(V).
\]


On the other hand, given a Lie algebroid
$\mathcal{L} = (V,[\cdot,\cdot],\delta)$ on $X$ and a parameter
$\lambda \in \mathbf{A}^{1}$, we define a new Lie algebroid
\[
\mathcal{L}_{\lambda} := (V,[\cdot,\cdot],\lambda\delta).
\]
A straightforward computation shows that
\[
\Lambda^{\mathrm{red}}_{\mathcal{L}}|_{X\times\{\lambda\}}
\cong \Lambda^{\text{DR}}_{\mathcal{L}_{\lambda}}.
\]

Observe that for $\lambda \neq 0$, multiplication by $\lambda$ induces
an isomorphism of Lie algebroids
\[
\mathcal{L}_{\lambda} \cong \mathcal{L},
\qquad v \longmapsto \lambda v,
\]
whereas for $\lambda = 0$ we have
$\mathcal{L}_{0} = (V,[\cdot,\cdot],0)$, the trivial Lie algebroid on $V$.
Consequently, the fibers of $\Lambda^{\mathrm{red}}_{\mathcal{L}}$
over $\lambda \in \mathbf{A}^{1}$ satisfy the following properties,
which are well-known consequences of the \emph{Rees construction}:
\begin{itemize}
	\item for every $\lambda \neq 0$,
	\[
	\Lambda^{\mathrm{red}}_{\mathcal{L}}|_{X\times\{\lambda\}}
	\cong \Lambda^{\mathrm{red}}_{\mathcal{L}}|_{X\times\{1\}}
	\cong \Lambda^{\text{DR}}_{\mathcal{L}},
	\]
	\item for $\lambda = 0$,
	\[
	\Lambda^{\mathrm{red}}_{\mathcal{L}}|_{X\times\{0\}}
	\cong \Lambda^{\mathrm{Dol}}_{\mathcal{L}}
	\cong \operatorname{Gr}_{\bullet}(\Lambda^{\text{DR}}_{\mathcal{L}})
	\cong \operatorname{Sym}^{\bullet}(V).
	\]
\end{itemize}

By lemma \eqref{Lie algebroid is-equivalent-to-Lambda_L}, let
$\mathcal{L}^{\mathrm{red}}$ denote the Lie algebroid on
$X \times \mathbf{A}^{1}$ corresponding to
$\Lambda^{\mathrm{red}}_{\mathcal{L}}$.
We consider the moduli space
$\mathcal{M}^{\mathrm{Hod}}_{\mathcal{L}}(X,n)$
of vector bundles of rank $n$ equipped with integrable
$\mathcal{L}^{\mathrm{red}}$--connections on $X \times \mathbf{A}^{1}$.
By Theorem~\eqref{eqn: moduli-of-L-connection on VB}, this is a
quasi-projective variety endowed with a natural morphism
\[
\pi : \mathcal{M}^{\mathrm{Hod}}_{\mathcal{L}}(X,n)
\longrightarrow \mathbf{A}^{1}, \ \
(\mathcal{E},D_{\mathcal{L}},\lambda)\mapsto \lambda.
\]

It follows immediately that for all $\lambda \in \mathbf{A}^{1}\setminus\{0\}$,
\[
\pi^{-1}(\lambda)
= \mathcal{M}^{\text{DR}}_{\mathcal{L}_{\lambda}}(X,n)
\cong \mathcal{M}^{\text{DR}}_{\mathcal{L}_{1}}(X,n)
= \pi^{-1}(1),
\qquad
\pi^{-1}(0)
= \mathcal{M}^{\mathrm{Dol}}_{\mathcal{L}}(X,n).
\]
Thus, $\mathcal{M}^{\mathrm{Hod}}_{\mathcal{L}}(X,n)$
provides an interpolation between the moduli space of
$\mathcal{L}$--connections on vector bundles and the moduli space
of $\mathcal{L}$--twisted Higgs bundles.
We call $\mathcal{M}^{\mathrm{Hod}}_{\mathcal{L}}(X,n)$ the
\emph{$\mathcal{L}$--Hodge moduli space for vector bundles}.

Hence, we obtain the following result.
\begin{thm}\label{def:L-Hodge-moduli-space-for vector bundles}
	The $\mathcal{L}$--Hodge moduli space of rank $n$ vector bundles,
	$\mathcal{M}^{\mathrm{Hod}}_{\mathcal{L}}(X,n)$,
	is a quasi-projective scheme equipped with a morphism
	\[
	\pi : \mathcal{M}^{\mathrm{Hod}}_{\mathcal{L}}(X,n)
	\longrightarrow \mathbf{A}^{1}.
	\]
\end{thm}


We now introduce the notion of $\mathcal{L}$--twisted principal $G$--Higgs bundles,
which generalizes the classical notion of principal $G$--Higgs bundles.
Let $\mathfrak{g}$ denote the Lie algebra of $G$, equipped with the adjoint action of $G$.

\begin{defn}\label{def: L-twisted-principal-G-Higgs-bundles}
	An \emph{$\mathcal{L}$--twisted principal $G$--Higgs bundle on $X$} consists of a
	principal $G$--bundle $E_{G} \to X$ together with a section
	\[
	\theta_{\mathcal{L}} \in H^{0}\!\left(X,
	\text{ad}(E_{G}) \otimes \Omega^{1}_{\mathcal{L}}\right),
	\]
	called the \emph{$\mathcal{L}$--twisted $G$--Higgs field}, satisfying the integrability condition
	\[
	[\theta_{\mathcal{L}}, \theta_{\mathcal{L}}] = 0
	\quad \text{in} \quad
	H^{0}\!\left(X, \text{ad}(E_{G}) \otimes \Omega^{2}_{\mathcal{L}}\right).
	\]
\end{defn}

\begin{example}
An $\mathcal{L}$--twisted $SL(n, \mathbb{C})$--Higgs bundle is a pair $(\mathcal{E}, \theta_{\mathcal{L}})$, where $\mathcal{E} \rightarrow X$ is a rank
$n$ vector bundle with $\det(\mathcal{E}) = \mathcal{O}_{X}$ and $\theta_{\mathcal{L}} \in H^{0}(X, \End(\mathcal{E}) \otimes \Omega^{1}_{\mathcal{L}})$ with $tr(\theta_{\mathcal{L}}) = 0$.
\end{example}

\begin{example}
An $\mathcal{L}$--twisted $SO(n, \mathbb{C})$--Higgs bundle is a pair $(E, \theta_{\mathcal{L}})$ where $E$ is a $SO(n, \mathbb{C})$--bundle
and $\theta_{\mathcal{L}}\in H^{0}(E(\mathfrak{so}(n, C)) \otimes \Omega^{1}_{\mathcal{L}})$. Using the standard representations of $SO(n, \mathbb{C})$ in $\mathbb{C}^{n}$ we can associate to $E$ a vector bundle $W$ of rank $n$ with trivial determinant, $W = E\times^{SO(n,\mathbb{C})}\mathbb{C}^{n}$, together with a non-degenerate symmetric quadratic form $Q\in H^{0}(S^{2}W^{*})$; we can think
of $Q$ as a symmetric isomorphism $Q : W \rightarrow W^{*}$. The Higgs field in terms of the vector bundle $W$ is a section $\theta_{\mathcal{L}}\in H^{0}(End(W)\otimes \Omega^{1}_{\mathcal{L}})$ satisfying $Q(u, \theta_{\mathcal{L}}(v)) =
-Q(\theta_{\mathcal{L}}(u), v)$ and $\text{tr}(\theta_{\mathcal{L}}) = 0$.
\end{example}




%%%%%%%%%%%%%%%%%%%%%%%%%%%%%%%%%%%%%%%%%%%%%%%%%%%%%%%%%%%%%%%%%%%%%%%%%
%%%%%%%%%%%%%%%%%%%%%%%%%%%%%%%%%%%%%%%%%%%%%%%%%%%%%%%%%%%%%%%%%%%%%%%%%



\section{Tannakian description of $\mathcal{L}$--twisted principal objects}\label{sec:Tannakian Considerations}

The theory of Higgs bundles and integrable connections admits a natural extension in the framework of Lie algebroids, unifying several geometric structures such as ordinary Higgs bundles, logarithmic Higgs bundles, and generalized connections. From the Tannakian perspective, principal bundles can be described by tensor functors from the category of representations of a linear algebraic group, providing a conceptual approach to moduli problems in the classical setting \cite{CT.Simpson-1995}.

In this paper, we develop a Tannakian description of principal $\mathcal{L}$-twisted $G$-Higgs bundles and principal $G$-bundles equipped with integrable Lie algebroid connections, where $\mathcal{L}$ is a Lie algebroid over a smooth projective variety $X$. This framework provides a natural foundation for the construction and study of their moduli spaces, extending the classical theory of principal Higgs bundles and integrable principal bundles to the Lie algebroid setting.


For completeness, we briefly recall the necessary background on tensor categories and neutral Tannakian categories, for details, see \cite[Chapters~1--2]{J.S.Milne-2025}, \cite{Deligne-Milne-Ogus-Shih-1982}, and \cite[\S\,6]{CT. Simpson-1992}.


A \emph{tensor category} is a category $\mathcal{C}$ equipped with a bifunctor
\[
\otimes : \mathcal{C} \times \mathcal{C} \longrightarrow \mathcal{C},
\]
called the \emph{tensor product}, together with natural isomorphisms expressing
associativity and commutativity of $\otimes$, and a distinguished object
$\mathbb{1} \in \mathcal{C}$ called the \emph{unit object}.  
These natural isomorphisms are referred to as \emph{constraints}; for example,
the commutativity constraint is a natural isomorphism
\[
U \otimes V \cong V \otimes U.
\]
The constraints are required to satisfy a collection of coherence conditions,
usually called the \emph{canonical axioms}, which assert that any natural
automorphism obtained by composing the constraints is the identity.

The unit object $\mathbb{1}$ is equipped with natural isomorphisms
\[
\mathbb{1} \otimes V \cong V,
\]
satisfying the canonical axioms.

A \emph{tensor functor} between tensor categories is a functor
$\mathscr{F} : \mathcal{C} \to \mathcal{D}$ endowed with natural isomorphisms
\[
\mathscr{F}(U \otimes V) \cong \mathscr{F}(U) \otimes \mathscr{F}(V),
\]
compatible with the constraints and satisfying the canonical axioms.

A \emph{neutral Tannakian category} is an abelian tensor category $\mathcal{T}$
which is rigid (i.e.\ every object admits a dual), satisfies
$\mathrm{End}(\mathbb{1}) = \mathbb{C}$, and is equipped with an exact faithful
$\mathbb{C}$--linear tensor functor
\[
\omega : \mathcal{T} \longrightarrow \texttt{Vect}_{\mathbb{C}},
\]
called a \emph{fiber functor}.

If $G$ is an affine algebraic group over $\mathbb{C}$, then the category
$\texttt{Rep}(G)$ of finite-dimensional complex representations of $G$ is a
neutral Tannakian category. The forgetful functor
\[
\omega_{G} : \texttt{Rep}(G) \longrightarrow \texttt{Vect}_{\mathbb{C}},
\]
which sends a representation to its underlying vector space, is a fiber
functor. The group $G$ is recovered as the group
\[
G \cong \mathrm{Aut}^{\otimes}(\omega_{G}, \texttt{Rep}(G))
\]
of tensor automorphisms of the fiber functor.

Conversely, given a neutral Tannakian category $(\mathcal{T}, \omega)$, setting
\[
G := \mathrm{Aut}^{\otimes}(\omega, \mathcal{T}),
\]
one recovers an equivalence of tensor categories
\[
\mathcal{T} \simeq \texttt{Rep}(G).
\]


Let $\texttt{Vect}(X)$ denote the category of vector bundles (i.e.\ locally free
$\mathcal{O}_{X}$--modules of finite rank) on $X$. The category
$\texttt{Vect}(X)$ is naturally equipped with a tensor product
\[
\otimes : \texttt{Vect}(X) \times \texttt{Vect}(X) \longrightarrow
\texttt{Vect}(X),
\qquad (E,F) \longmapsto E \otimes F,
\]
and with a direct sum operation $\oplus$. With these structures,
$\texttt{Vect}(X)$ is a rigid, additive, $\mathbb{C}$--linear tensor category
whose unit object is $\mathbb{1}=\mathcal{O}_{X}$ and which satisfies
$\mathrm{End}(\mathbb{1})=\mathbb{C}$.

A morphism $u:E\to F$ in $\texttt{Vect}(X)$ is said to be \emph{strict} if
$\mathrm{coker}(u)$ is a locally free sheaf. In this case, both
$\ker(u)$ and $\mathrm{im}(u)$ are locally free. 
An alternative description of principal $G$--bundles, due to M.\,V.~Nori, may be
formulated in this framework. Let $E_{G}$ be a principal $G$--bundle on $X$.
Then $E_{G}$ determines a tensor functor
\begin{equation}\label{eq:nori-functor}
	\rho_{E_{G}} : \texttt{Rep}(G) \longrightarrow \texttt{Vect}(X),
\end{equation}
defined by
\[
\rho_{E_{G}}(V) := E_{G} \times^{G} V .
\]
This functor has the following properties:
\begin{itemize}
	\item $\rho_{E_{G}}$ is a strict tensor functor; that is, for any morphism
	$u:V\to W$ in $\texttt{Rep}(G)$, the induced morphism
	$\rho_{E_{G}}(u)$ is strict in $\texttt{Vect}(X)$;
	
	\item $\rho_{E_{G}}$ is exact and faithful;
	
	\item for every closed point $x\in X$, the functor
	\[
	\texttt{Rep}(G)\longrightarrow \texttt{Vec}_{\mathbb{C}},
	\qquad V \longmapsto \rho_{E_{G}}(V)_{x},
	\]
	is a fiber functor.
\end{itemize}
(see \cite{Balaji-2006}).

Nori proved the following converse statement.

\begin{proposition}\label{def:Nori-descrpton of PB}\cite{M.V.Nori-1976}
	Suppose $\rho:\texttt{Rep}(G)\longrightarrow \texttt{Vect}(X)$
	is a strict, exact, and faithful tensor functor such that for every closed point
	$x\in X$ the functor $V\mapsto \rho(V)_{x}$ is a fiber functor on
	$\texttt{Rep}(G)$. Then there exists a principal right $G$--bundle $E_{G}$ over
	$X$ and an isomorphism of tensor functors $\rho \cong \rho_{E_{G}}$. Moreover,
	$E_{G}$ is unique up to isomorphism.
\end{proposition}


\medskip

\noindent\textbf{Tannakian description of principal $G$--bundle equipped with integrable $\mathcal{L}$--connection}. Let $\mathcal{L}$ be an algebraic Lie algebroid on $X$. We denote by
$\texttt{Vect}^{\text{int}}_{\mathcal{L}}(X)$ the category whose objects are pairs
$(E,\nabla_{\mathcal{L}})$, where $E$ is an algebraic vector bundle on $X$
and $\nabla_{\mathcal{L}}$ is an integrable $\mathcal{L}$--connection on $E$.
A morphism
\[
h : (E,\nabla_{\mathcal{L}}) \longrightarrow (E',\nabla'_{\mathcal{L}})
\]
is a morphism of vector bundles $h:E\to E'$ such that
\[
\nabla'_{\mathcal{L}} \circ h
= (h \otimes \mathrm{id}_{\Omega^{1}_{\mathcal{L}}}) \circ \nabla_{\mathcal{L}} .
\]

Let $(E,\nabla^{E}_{\mathcal{L}})$ and $(F,\nabla^{F}_{\mathcal{L}})$ be vector
bundles equipped with integrable $\mathcal{L}$--connections. Define a
connection on $E\otimes F$ by
\[
\nabla^{E\otimes F}_{\mathcal{L}}(e\otimes f)
:= \nabla^{E}_{\mathcal{L}}(e)\otimes f
+ e\otimes \nabla^{F}_{\mathcal{L}}(f),
\]
for all locally defined sections $e$ of $E$ and $f$ of $F$. One checks
immediately that $\nabla^{E\otimes F}_{\mathcal{L}}$ is an integrable
$\mathcal{L}$--connection.

The usual associativity and symmetry isomorphisms of vector bundles,
\[
E\otimes(F\otimes G)\cong (E\otimes F)\otimes G,
\qquad
E\otimes F \cong F\otimes E,
\]
are compatible with the induced $\mathcal{L}$--connections. The unit object
is given by
\[
\mathbb{1} := (\mathcal{O}_{X}, d_{\mathcal{L}}),
\]
where $d_{\mathcal{L}}(f)=\delta^{*}(df)\in \Omega^{1}_{\mathcal{L}}$.

For a vector bundle $E$, let $E^{*}:=\mathcal{H}om_{\mathcal{O}_{X}}(E,\mathcal{O}_{X})$.
If $\nabla_{\mathcal{L}}$ is an integrable $\mathcal{L}$--connection on $E$,
then the dual bundle $E^{*}$ carries an induced integrable
$\mathcal{L}$--connection $\nabla^{*}_{\mathcal{L}}$, defined by
\[
d_{\mathcal{L}}\langle e,\xi\rangle
= \langle\nabla^{*}_{\mathcal{L}}(e),\xi\rangle
+ \langle e,\nabla_{\mathcal{L}}(\xi)\rangle,
\]
for all local sections $e\in E^{*}$ and $\xi\in E$, where
$\langle\cdot,\cdot\rangle : E^{*}\times E\to \mathcal{O}_{X}$ denotes the
natural pairing.

It follows that $\texttt{Vect}^{\text{int}}_{\mathcal{L}}(X)$ is a rigid, additive,
$\mathbb{C}$--linear tensor category with
$\mathrm{End}(\mathbb{1})=\mathbb{C}$.


It is well known that a coherent sheaf equipped with an integrable
$\mathcal{T}_{X}$--connection is necessarily locally free (see
\cite[\S\,VI, Proposition~1.7]{A.Borel-1987}). However, for a general Lie
algebroid $\mathcal{L}$, a coherent sheaf equipped with an integrable
$\mathcal{L}$--connection need not be locally free (see
\cite[\S\,2.2, Example~2.12]{GLP-2018}).

 However, if  $\mathcal{L}$ is a \emph{transitive Lie algebroid}. then this pathology above does not occur.

\begin{proposition}\label{def:locally free L-modulues}
	\cite[\S\,2.2, Proposition~2.15]{GLP-2018}
	Let $\mathcal{L}$ be a transitive Lie algebroid on $X$ and let
	$(\mathcal{E},\nabla_{\mathcal{L}})$ be a coherent sheaf equipped with an
	$\mathcal{L}$--connection (not necessarily integrable). Then $\mathcal{E}$ is
	locally free, i.e.\ a vector bundle.
\end{proposition}


Now we prove that $\texttt{Vect}^{\text{int}}_{\mathcal{L}}(X)$ is a neutral Tannakian
category when $\mathcal{L}$ is transitive.

\begin{proposition}\label{prop: tannakian category Vect_L(X)}
	Let $\mathcal{L}$ be a transitive Lie algebroid on $X$. Then the category
	$\texttt{Vect}^{\emph{int}}_{\mathcal{L}}(X)$ of vector bundles equipped with integrable
	$\mathcal{L}$--connections is a neutral Tannakian category over $\mathbb{C}$.
\end{proposition}

\begin{proof}
	By the previous discussion, $\texttt{Vect}^{\text{int}}_{\mathcal{L}}(X)$ is a rigid,
	additive, $\mathbb{C}$--linear tensor category with
	$\mathrm{End}(\mathbb{1})=\mathbb{C}$. It remains to show that it is abelian and
	admits a fiber functor.
	
	\medskip
	\noindent
	\emph{Abelian property.}
	Let
	\[
	h:(E,\nabla_{\mathcal{L}})\longrightarrow (E',\nabla'_{\mathcal{L}})
	\]
	be a morphism in $\texttt{Vect}^{\text{int}}_{\mathcal{L}}(X)$. Then we have a commutative
	diagram of $\mathcal{O}_{X}$--modules
	\[
	\begin{CD}
		0 @>>> \ker(h) @>>> E @>h>> E' @>>> \operatorname{coker}(h) @>>> 0 \\
		@. @VVV @V\nabla_{\mathcal{L}}VV @V\nabla'_{\mathcal{L}}VV @VVV @. \\
		0 @>>> \ker(h)\!\otimes\!\Omega^{1}_{\mathcal{L}}
		@>>> E\!\otimes\!\Omega^{1}_{\mathcal{L}}
		@>h\otimes\mathrm{id}>>
		E'\!\otimes\!\Omega^{1}_{\mathcal{L}}
		@>>> \operatorname{coker}(h)\!\otimes\!\Omega^{1}_{\mathcal{L}}
		@>>> 0 .
	\end{CD}
	\]
	
	The commutativity of the diagram shows that $\ker(h)$ and
	$\operatorname{coker}(h)$ inherit $\mathcal{L}$--connections. Since $h$
	intertwines integrable connections, these induced connections are also
	integrable.
	
	Because $\mathcal{L}$ is transitive, Proposition
	\eqref{def:locally free L-modulues} implies that both $\ker(h)$ and
	$\operatorname{coker}(h)$ are locally free. Hence kernels and cokernels exist
	in $\texttt{Vect}^{\text{int}}_{\mathcal{L}}(X)$, and the category is abelian.
	
	\medskip
	\noindent
	\emph{Fiber functor.}
	Fix a closed point $x\in X$. Define a functor
	\[
	\omega_{x}:\texttt{Vect}^{\text{int}}_{\mathcal{L}}(X)\longrightarrow
	\texttt{Vec}_{\mathbb{C}}, \qquad
	(E,\nabla_{\mathcal{L}})\longmapsto E_{x}.
	\]
	This functor is $\mathbb{C}$--linear, exact (since taking fibers of vector
	bundles is exact), faithful, and compatible with tensor products and duals.
	Therefore, $\omega_{x}$ is a fiber functor.
	
	\medskip
	Combining these facts, we conclude that $\texttt{Vect}^{\text{int}}_{\mathcal{L}}(X)$ is a
	neutral Tannakian category over $\mathbb{C}$.
\end{proof}

\begin{defn}
$\mathfrak{T}$ be a neutral Tannakian category. A \textit{$G$--torsor in $\mathfrak{T}$} is a strict exact faithful tensor functor $p:\text{Rep}(G)\rightarrow \mathfrak{T}$.
\end{defn}

The following lemma is well known for principal $G$--bundles equipped with
integrable $\mathcal{T}_{X}$--connections
(cf.~\cite[\S\,9, p.~55]{CT.Simpson-1995}). We extend this result to the setting of algebraic transitive Lie algebroid.


\begin{lem}\label{lem:PB-with-L-conection-Tannakian equivalence}
	Let $\mathcal{L}$ be a transitive Lie algebroid. The construction $E_{G}\longmapsto \rho_{E_{G}}$
	in \eqref{eq:nori-functor} induces an equivalence between:
	\begin{itemize}
		\item the category of principal $G$--bundles on $X$ equipped with integrable
		$\mathcal{L}$--connections, and
		\item the category of $G$ torsor in $\texttt{Vect}^{\text{int}}_{\mathcal{L}}(X)$,
		such that for every closed point $x\in X$, the functor
		$V\mapsto \rho(V)_{x}$ is a fiber functor on $\texttt{Rep}(G)$.
	\end{itemize}
\end{lem}
\begin{proof}
	Let $(E_{G},\nabla_{\mathcal{L}})$ be a principal $G$--bundle on $X$ equipped with an integrable $\mathcal{L}$--connection.  
	For any representation $\varphi:G\to GL(W)$ in $\texttt{Rep}(G)$, the induced Lie algebra homomorphism
	\[
	d\varphi:\mathfrak g\longrightarrow \mathfrak{gl}(W)
	\]
	gives rise, by extension of structure group, to a morphism of associated vector bundles
	\[
	\ad(E_G)=E_G\times^{\mathrm{ad}}\mathfrak g
	\;\longrightarrow\;
	\End(E_G\times^{\varphi}W)=E_G\times^{\mathrm{ad}}\mathfrak{gl}(W).
	\]
	
	This morphism is compatible with the $\mathcal{L}$--twisted Atiyah sequences, yielding a morphism of $\mathcal{L}$--Lie algebroids
	\[
	\At_{\delta}(E_G)\longrightarrow \At_{\delta}(E_G\times^{\varphi}W).
	\]
	Consequently, we obtain the commutative diagram
	\[
	\begin{CD}
		0@>>>\ad(E_G)@>>>\At_{\delta}(E_G)@>>> \mathcal{L}@>>>0\\
		@. @VVV @VVV @| @.\\
		0@>>>\End(E_G\times^{\varphi}W)@>>>\At_{\delta}(E_G\times^{\varphi}W)@>>> \mathcal{L}@>>>0 .
	\end{CD}
	\]
	The integrable $\mathcal{L}$--connection $\nabla_{\mathcal{L}}$ corresponds to a splitting of the top exact sequence, and by functoriality this induces a splitting of the bottom sequence. Hence the associated vector bundle
	$E_G\times^{\varphi}W$ carries a natural integrable $\mathcal{L}$--connection.
	
	Thus the assignment
	\[
	\rho_{E_G}:\texttt{Rep}(G)\longrightarrow \texttt{Vect}^{\text{int}}_{\mathcal{L}}(X),\qquad
	W\longmapsto (E_G\times^{G}W,\nabla^{\varphi}_{\mathcal{L}})
	\]
	defines a strict, exact, faithful tensor functor.
	
	\medskip
	
	Conversely, let
	\[
	\rho:\texttt{Rep}(G)\longrightarrow \texttt{Vect}^{\text{int}}_{\mathcal{L}}(X)
	\]
	be a strict, exact, faithful tensor functor such that for each closed point $x\in X$ the functor
	$V\mapsto \rho(V)_x$ is a fiber functor on $\texttt{Rep}(G)$.  
	Since $\texttt{Vect}_{\mathcal{L}}(X)$ is a neutral Tannakian category by Proposition~\eqref{prop: tannakian category Vect_L(X)}, Tannakian reconstruction (cf.\ \cite[\S\,6, remark after lemma 6.13]{CT. Simpson-1992} and Nori's description \eqref{def:Nori-descrpton of PB}) yields a principal $G$--bundle $E_G$ on $X$ together with a flat family of fiber functors
	\[
	\mathscr{F}_x:\texttt{Rep}(G)\longrightarrow \Vect_{\mathbb C},\qquad
	\mathscr{F}_x(V)=\rho(V)_x,
	\]
	varying algebraically with $x\in X$, and satisfying $\rho(V)=E_G\times^{G}V$.
	
	\medskip
	
	\noindent\textit{Induced integrable $\mathcal{L}$--connection.}
	The $\mathcal{L}$--connection on each object $\rho(V)$ is functorial in $V$ and compatible with tensor products, duals, and morphisms in $\texttt{Rep}(G)$. Equivalently, the family of fiber functors $\{\mathscr{F}_x\}_{x\in X}$ is flat along $\mathcal{L}$ in the sense that for every $V\in\texttt{Rep}(G)$, the vector bundle $E_G\times^{G}V$ is equipped with an integrable $\mathcal{L}$--connection.
	
	This functorial flatness uniquely determines a splitting of the $\mathcal{L}$--twisted Atiyah sequence
	\[
	0\longrightarrow \ad(E_G)\longrightarrow \At_{\delta}(E_G)\longrightarrow \mathcal{L}\longrightarrow 0,
	\]
	hence an integrable $\mathcal{L}$--connection on the principal bundle $E_G$.  
	This completes the proof.
\end{proof}

                          %%%%%%%%%%%%%%%%%%%%%
                          
                          
\medskip

\noindent\textbf{Tannakian description of $\mathcal{L}$--twisted principal $G$--Higgs bundle}. Let $E_G$ be a $\mathcal{L}$--twisted principal $G$--Higgs bundle on $X$. For any invariant polynomial
$P \in \mathbb{C}[\mathfrak g]^G$, the associated $\mathcal L$--Chern class
is defined by the Chern--Weil construction as
\[
c^{\mathcal L}_P(E_G) := P(F_{\nabla_{\mathcal L}}) \in H^{2\deg P}_{\mathcal L}(X),
\]
where $F_{\nabla_{\mathcal L}}$ is the curvature of the $\mathcal L$--connection.
These classes are independent of the choice of $\nabla_{\mathcal L}$ and
satisfy $c^{\mathcal L}_P(E_G)=\delta^*(c_P(E_G))$ via the anchor map (cf.\cite{PM-2020},\cite{Crainic-2009}).

We say $E_{G}$ is of
\textit{semiharmonic type} if the $\mathcal{L}$--chern classes of $E_{G}$ is zero in rational cohomology, and if there exist a faithful representation $V$ such that $E_{G}\times^{G}V$ is of semiharmonic type. Similarly, we say that $E_{G}$ is of \emph{harmonic type} if its
$\mathcal{L}$--Chern classes vanish in rational cohomology and if there
exists a faithful representation $V$ of $G$ such that the associated
vector bundle $E_{G}\times^{G} V$ is of harmonic type. In this case, the same is true for any other representation (cf. \cite[\S\,6, remarks after lemma 6.13]{CT. Simpson-1992}).



We now prove that the category of $\mathcal{L}$--twisted Higgs bundles of
semiharmonic type on $X$ is also a neutral Tannakian category.  
The proof is entirely analogous to the case of vector bundles equipped with integrable $\mathcal{L}$--connections.
\begin{proposition}\label{prop:Tannakian-L-Higgs}
	Let $X$ be a smooth projective variety and let $\mathcal L$ be a transitive Lie algebroid on $X$.  
	Then the category $\texttt{Higg}^{\emph{sh}}_{\mathcal L}(X)$ of
	$\mathcal L$--twisted Higgs bundles of semiharmonic type on $X$ is a
	neutral Tannakian category.
\end{proposition}

\begin{proof}
	By definition, objects of $\texttt{Higg}^{\text{sh}}_{\mathcal L}(X)$ are
	pairs $(E,\theta_{\mathcal L})$ where $E$ is a vector bundle and
	$\theta_{\mathcal L}$ is an integrable $\mathcal L$--twisted Higgs field
	satisfying the semiharmonic condition. Morphisms are Higgs--compatible
	bundle morphisms.
	
	The tensor product of two objects is defined using the usual tensor
	product of vector bundles together with the induced $\mathcal L$--Higgs
	field, and duals exist with the natural induced Higgs fields. Hence the
	category is rigid, additive, and $\mathbb C$--linear, with unit object
	$(\mathcal O_X,0)$ and $\mathrm{End}(\mathbb 1)=\mathbb C$.
	
	Since $\mathcal L$ is transitive, the kernel and cokernel of any morphism of
	$\mathcal L$--twisted Higgs bundles are locally free by
	Proposition~\eqref{prop:locally-free-L-Higgs}, and they naturally inherit
	$\mathcal L$--twisted Higgs fields of semiharmonic type.  
	Therefore, the category $\texttt{Higg}^{\text{sh}}_{\mathcal L}(X)$ is an abelian tensor category.
	
	Finally, choosing a closed point $x\in X$, the fiber functor
	\[
	(E,\theta_{\mathcal L}) \longmapsto E_x
	\]
	is exact, faithful, and tensor compatible. Hence
	$\texttt{Higg}^{\text{sh}}_{\mathcal L}(X)$ is a neutral Tannakian
	category.
\end{proof}
Hence, we obtain the following lemma, which is the $\mathcal{L}$--twisted
Higgs analogue of Lemma~\eqref{lem:PB-with-L-conection-Tannakian equivalence}.
It generalizes \cite[\S\,9, Lemma~9.2]{CT.Simpson-1995} to the setting of
Lie algebroids.

\begin{lem}\label{lem: L-twisted PB-Higgs bundle-tannakian}
	The construction $E_{G}\mapsto \rho_{E_{G}}$ in
	\eqref{eq:nori-functor} provides an equivalence between the category of $\mathcal{L}$--twisted principal $G$--Higgs bundles of semiharmonic type and the category of $G$--torsor in $\texttt{Higg}^{\emph{sh}}_{\mathcal L}(X)$,
	such that for every closed point $x\in X$ the functor
	$V\mapsto \rho(V)_x$ is a fiber functor on $\texttt{Rep}(G)$.
\end{lem}

\begin{proof}
	The proof is a direct adaptation of
	Lemma~\eqref{lem:PB-with-L-conection-Tannakian equivalence}, following
	Simpson’s argument in \cite[\S\,9, Lemma~9.2]{CT.Simpson-1995}.
	
	Let $(E_G,\theta_{\mathcal L})$ be an $\mathcal L$--twisted principal
	$G$--Higgs bundle of semiharmonic type. For any representation
	$\varphi:G\to \mathrm{GL}(W)$, the induced Higgs field on the associated
	bundle $E_G\times^G W$ is obtained by composing $\theta_{\mathcal L}$ with the differential $d\varphi:\mathfrak g\to\mathfrak{gl}(W)$. The integrability condition and vanishing of $\mathcal L$--Chern classes are preserved under this construction, hence $\rho_{E_G}(W)$ is an $\mathcal L$--twisted Higgs bundle of semiharmonic type. This defines a strict exact faithful tensor functor
	\[
	\rho_{E_G}:\texttt{Rep}(G)\longrightarrow
	\texttt{Higg}^{\mathrm{sh}}_{\mathcal L}(X).
	\]
	
	Conversely, let $\rho:\texttt{Rep}(G)\to
	\texttt{Higg}^{\mathrm{sh}}_{\mathcal L}(X)$ be a strict exact faithful
	tensor functor such that for every closed point $x\in X$ the functor
	$V\mapsto \rho(V)_x$ is a fiber functor. Since
	$\texttt{Higg}^{\mathrm{sh}}_{\mathcal L}(X)$ is a neutral Tannakian
	category, Nori’s Tannakian reconstruction yields a principal $G$--bundle
	$E_G$ on $X$. By \cite[\S\,6, remark after Lemma~6.13]{CT. Simpson-1992}, the Higgs fields
	on the objects $\rho(V)$ assemble functorially to an
	$\mathcal L$--twisted Higgs field $\theta_{\mathcal L}$ on the principal $G$--bundle $E_G$. By construction, $\theta_{\mathcal L}$ is of semiharmonic type.
	
\end{proof}



%%%%%%%%%%%%%%%%%%%%%%%%%%%%%%%%%%%%%%%%%%%%%%%%%%%%%%%%%%%%%%%%%%%%%%%%%
%%%%%%%%%%%%%%%%%%%%%%%%%%%%%%%%%%%%%%%%%%%%%%%%%%%%%%%%%%%%%%%%%%%%%%%%

\section{Moduli spaces of $\mathcal{L}$-twisted principal objects}

Before proceeding to the construction, we recall an important fact from
\cite[\S\,Proof of Thm.~3.8]{CT.Simpson-1994}, which will be used later.

\begin{lem}\label{lem:fact-from-simpson-thm-3-8}
	Let $S$ be a $\mathbb{C}$--scheme of finite type, and let
	$E$ and $F$ be coherent $\mathcal{O}_{X'}$--modules on
	$X' := X \times_{\mathbb{C}} S$.
	Assume that $E$ is flat over $S$ (no flatness assumption is imposed on $F$),
	and let $\varphi \colon F \to E$ be a morphism of coherent sheaves.
	Then there exists a closed subscheme $T \subset S$ such that for any morphism
	$f \colon S' \to S$, the pullback
	$f^*(\varphi)=0$ as a morphism of sheaves on
	$X'' := X' \times_S S'$ if and only if $f$ factors through $T$.
\end{lem}

\begin{proof}
	Since $X$ is projective and $F$ is a coherent $\mathcal{O}_{X'}$--module,
	by Serre's theorem \cite[\S\,Ch.~II, Thm.~5.17]{R. Hartshorne-1977},
	the sheaf $F(m)$ is generated by finitely many global sections for $m \gg 0$.
	Hence there exists a surjection
	\[
	\mathcal{O}_{X'}(-m)^{\oplus k} \twoheadrightarrow F \to 0 .
	\]
	This surjection remains surjective after any base change
	$X'' = X' \times_S S'$.
	
	Replacing $F$ by $\mathcal{O}_{X'}(-m)^{\oplus k}$, a morphism
	$\varphi \colon F \to E$ corresponds to a $k$--tuple of global sections of
	$E(m)$.
	For $m$ sufficiently large, the direct image $\pi_*E(m)$ is locally free on $S$
	and compatible with arbitrary base change.
	The desired closed subscheme $T \subset S$ is then defined as the intersection
	of the zero loci of the corresponding sections of $\pi_*E(m)$.
\end{proof}





%%%%%%%%%%%%%%%%%%%%%%%%%%%%%%%%%%%%%%%%%%%%%%%%%%%%%%%%%%%%%%%%%%%%%%%%%%
\subsection{The moduli spaces of integrable $\mathcal{L}$-connection on principal $G$-bundles}\label{susec: moduli space of L connection on PB}

We now define the moduli functor for principal $G$--bundles with integrable
$\mathcal L$--connections.

Fix a irreducible reductive algebraic group $G$ over $\mathbb{C}$ and a smooth projective variety
$X$ endowed with a transitive Lie algebroid $\mathcal L$. Let $\mathcal{P}_{0}$ denote the Hilbert polynomial of $\mathcal{O}_{X}$.
Let $\mathcal{M}^{\text{DR}^{\natural}}_{\mathcal L}(X,G)$ be the contravariant functor
\[
\mathcal{M}^{\text{DR}^{\natural}}_{\mathcal L}(X,G) :
\texttt{Sch}/\mathbb{C} \longrightarrow \texttt{Sets}
\]
defined as follows.

For a $\mathbb{C}$--scheme $S$, $\mathcal{M}^{\text{DR}^{\natural}}_{\mathcal L}(X,G)(S)$
is the set of isomorphism classes of pairs $(E_G,\nabla_{\mathcal L})$ on
$X_S := X \times_{\mathbb{C}} S$ such that:
\begin{itemize}\label{moduli of principal G- bundles with integrable connection}
	\item $E_G$ is a principal $G$--bundle on $X_S$;
	\item $\nabla_{\mathcal L}$ is an integrable $\mathcal L$--connection on
	$E_G$ relative to $S$;
	\item for every finite--dimensional representation
	$\rho : G \to \mathrm{GL}(V)$ of dimension $n$, the associated vector bundle
	\[
	E_G(V) := E_G \times^{G} V
	\]
	on $X_S$ is flat over $S$, and for every closed point
	$\bar{s} \to S$ the restriction $E_G(V)_{\bar{s}}$ has Hilbert polynomial
	$n \mathcal{P}_0$.
	\item for every closed point $s \in S$, the fiber
	$(E_{G,s},\nabla_{\mathcal L,s})$ is a principal
	$G$--bundle with integrable $\mathcal L$--connection on $X$.
\end{itemize}

Two such families $(E_G,\nabla_{\mathcal L})$ and $(E'_G,\nabla'_{\mathcal L})$
are isomorphic if there exists a $G$--bundle isomorphism
$\varphi : E_G \to E'_G$ commuting with the $\mathcal L$--connections.

For a morphism of schemes $\psi : T \to S$, the functor assigns the pullback
\[
(E_G,\nabla_{\mathcal L}) \longmapsto
(\psi^*E_G,\psi^*\nabla_{\mathcal L})
\]
on $X_T := X \times_{\mathbb{C}} T$.




The main result is that there exists a quasi-projective variety $\mathcal{M}^{\text{DR}}_{\mathcal{L}}(X,G)$
over $\mathbb{C}$ which universally corepresents the functor
$\mathcal{M}^{\text{DR}^{\natural}}_{\mathcal{L}}(X,G)$. 

We begin this construction by establishing a sequence of preparatory lemmas. The following lemma is well known in the case of vector bundles with integrable $\mathcal{T}_X$--connections
(cf.~\cite[\S\,9, Lemma~9.9]{CT.Simpson-1995}).
We extend it to the Lie algebroid setting by the same method.

\begin{lem}\label{lem:projective-scheme-N(E,k)}
	Let $\mathcal{L}$ be a transitive Lie algebroid on $X$, and let $(E,\nabla_{\mathcal L})$
	be a vector bundle equipped with an integrable $\mathcal L$--connection.
	Fix an integer $k$.
	Then there exists a projective scheme $\mathcal{N}(E,k)$ over $\mathbb{C}$
	representing the functor which assigns to each $\mathbb{C}$--scheme $S$
	the set of quotients
	\[
	p_X^*E \twoheadrightarrow F \to 0
	\quad \text{on } X_S := X\times S,
	\]
	which are compatible with the $\mathcal L$--connection.
	Moreover, the natural morphism
	\[
	\mathcal{N}(E,k) \longrightarrow
	\mathfrak{Grass}\bigl(E|_{\mathfrak p\times S},k\bigr)
	\]
	is a closed embedding.
\end{lem}

\begin{proof}
	Let $\Lambda^{\text{DR}}_{\mathcal L}$ denote the de Rham $\mathcal{L}$--twisted $D$--algebra associated to the Lie algebroid
	$\mathcal L$.
	By \cite[\S\, Ch.4. Prop. 38]{P.Tortella-2011T},
	the vector bundle $E$ equipped with integrable $\mathcal{L}$--connection may be regarded as a $\Lambda_{\mathcal L}$--module.
	
	Let $\mathcal P_0$ be the Hilbert polynomial of $\mathcal O_X$.
	The Quot scheme $\mathfrak{Quot}(E,k\mathcal P_0)$
	parametrizes quotient sheaves $E\twoheadrightarrow F\to 0$
	with $F$ flat over the base and Hilbert polynomial $k\mathcal P_0$.
	Let
	\[
	E^{\mathrm{univ}} \twoheadrightarrow F^{\mathrm{univ}} \to 0
	\]
	be the universal quotient over
	$X\times \mathfrak{Quot}(E,k\mathcal P_0)$,
	and let $K^{\mathrm{univ}}$ denote its kernel.
	
	Since $E^{\mathrm{univ}}$ is a $\Lambda^{\text{DR}^{\text{univ}}}_{\mathcal L}$--module,
	we obtain a natural morphism
	\[
	\psi^{\mathrm{univ}} :
	\Lambda^{\text{DR}^{\text{univ}}}_{\mathcal L}
	\otimes_{\mathcal O_X} K^{\mathrm{univ}}
	\longrightarrow F^{\mathrm{univ}}.
	\]
	By Lemma~\eqref{lem:fact-from-simpson-thm-3-8},
	the condition that $\psi^{\mathrm{univ}}$ vanishes defines a closed
	subscheme
	\[
	\mathcal{N}(E,k) \subset \mathfrak{Quot}(E,k\mathcal P_0).
	\]
	
	For any $S$--valued point
	$e:S\to \mathcal{N}(E,k)$,
	the induced morphism
	\[
	\psi_e :
	\Lambda^{\text{DR}^{S}}_{\mathcal L} \otimes K \to F
	\]
	vanishes, and hence the $\Lambda^{\text{DR}}_{\mathcal L}$--module structure on $p_X^*E$ descends to $F$.
	Thus $F$ acquires a $\Lambda^{\text{DR}}_{\mathcal L}$--module structure, i.e.,
	an integrable $\mathcal L$--connection.
	
	Since $F$ is flat over $S$ and all its geometric fibers are locally free,
	it follows that $F$ is locally free
	(cf.~\cite[\S\,1, Lemma~1.27]{CT.Simpson-1994}), and the Hilbert polynomial condition implies that $\operatorname{rk}(F)=k$.
	Therefore $\mathcal{N}(E,k)$ represents the desired functor.
	
	Finally, by
	\cite[\S\,4, Lemma~4.9]{CT.Simpson-1994},
	the natural morphism to the relative Grassmannian is a closed embedding.
\end{proof}


We now extend the notion of the monodromy group of principal $G$--objects
(cf.~\cite[\S\,9, p.~50]{CT.Simpson-1995}) to the setting of
$\mathcal{L}$--twisted principal $G$--objects.
Let $\mathfrak{p}\in X$ be a closed point, and let $G\subset H$ be an
algebraic subgroup.
Let $(E_H,\nabla_{\mathcal L})$ be a principal $H$--bundle on $X$ equipped
with an integrable $\mathcal{L}$--connection, and fix a point
$\mathfrak{b}\in E_H|_{\mathfrak{p}}\cong H$.

We say that the \emph{monodromy of $(E_H,\mathfrak{b})$ is contained in $G$}
if the following condition holds:
for every finite--dimensional linear representation $V$ of $H$ and every
subspace $W\subset V$ preserved by $G$, there exists a strict subbundle
\[
F \subset E_H\times^{H} V
\]
which is preserved by the induced $\mathcal{L}$--connection and satisfies
\[
F|_{\mathfrak{p}}=\{\mathfrak{b}\}\times W
\subset (E_H\times^{H} V)|_{\mathfrak{p}}.
\]

The \emph{monodromy group} $\mathbf{Mono}(E_H,\mathfrak{b})$ is defined to
be the intersection of all algebraic subgroups $G\subset H$ such that the
monodromy of $(E_H,\mathfrak{b})$ is contained in $G$.



The following two lemmas are well known in the case of principal Higgs
bundles of semiharmonic type on a smooth projective variety
(cf.~\cite[\S\,9, Lemmas~9.4 and~9.5]{CT.Simpson-1995}).
Here we derive analogous statements for principal $G$--bundles equipped
with integrable $\mathcal{L}$--connections.

\begin{lem}\label{monodromy related thm-for-PB with-L-connection}
	Let $G \subset H$ be algebraic groups. Let $X$ be a smooth projective variety
	and let $\mathcal{L}$ be a transitive Lie algebroid on $X$.
	Suppose $E_{G}$ is a principal $G$--bundle on $X$ equipped with an integrable
	$\mathcal{L}$--connection. Then the associated bundle
	\[
	E_{H} := E_{G} \times^{G} H
	\]
	admits a natural structure of a principal $H$--bundle with integrable
	$\mathcal{L}$--connection.
	
	Moreover, fixing a closed point $\mathfrak{p} \in X$, this construction induces
	a bijection between the following sets:
	\begin{enumerate}
		\item isomorphism classes of pairs $(E_{G},\mathfrak{b}')$, where $E_{G}$ is a
		principal $G$--bundle with integrable $\mathcal{L}$--connection and
		$\mathfrak{b}' \in (E_{G})_{\mathfrak{p}}$ is a point;
		
		\item isomorphism classes of pairs $(E_{H},\mathfrak{b})$, where $E_{H}$ is a
		principal $H$--bundle with integrable $\mathcal{L}$--connection and
		$\mathfrak{b} \in (E_{H})_{\mathfrak{p}}$ is a point, such that the monodromy of
		$(E_{H},\mathfrak{b})$ is contained in $G$.
	\end{enumerate}
\end{lem}

\begin{proof}
	Let $E_{G}$ be a principal $G$--bundle on $X$ equipped with an integrable
	$\mathcal{L}$--connection. Choose a faithful finite--dimensional representation
	$V$ of $H$. Restricting $V$ to $G$ yields a faithful representation of $G$, and
	we have a natural identification
	\[
	E_{H} \times^{H} V \;=\; E_{G} \times^{G} V .
	\]
	By Lemma~\eqref{lem:PB-with-L-conection-Tannakian equivalence},
	$E_{H} \times^{H} V$ is a vector bundle equipped with an integrable
	$\mathcal{L}$--connection. Applying the same lemma once again, this vector
	bundle determines a unique integrable $\mathcal{L}$--connection on the
	principal $H$--bundle $E_{H}$. This defines a functor from the category of
	objects in $(1)$ to the category of objects in $(2)$.
	
	\medskip
	For the converse, let $\texttt{Rep}(G,H)$ denote the category whose objects are
	pairs $(V,W)$, where $W$ is a finite--dimensional representation of $H$ and
	$V \subset W$ is a $G$--invariant subspace. Morphisms in this category are
	$G$--equivariant linear maps between the subspaces $V$. The tensor product is
	defined by
	\[
	(V_{1},W_{1}) \otimes (V_{2},W_{2})
	\;:=\;
	(V_{1} \otimes V_{2},\, W_{1} \otimes W_{2}).
	\]
	Forgetting the $H$--representation $W$ defines an equivalence of tensor
	categories
	\[
	\texttt{Rep}(G,H) \;\cong\; \texttt{Rep}(G).
	\]
	
	Now suppose $(E_{H},\mathfrak{b})$ is a principal $H$--bundle with integrable
	$\mathcal{L}$--connection and a point $\mathfrak{b} \in (E_{H})_{\mathfrak{p}}$
	such that the monodromy of $(E_{H},\mathfrak{b})$ is contained in $G$. By
	definition, for every object $(V,W)$ of $\texttt{Rep}(G,H)$ there exists a unique
	subbundle
	\[
	F(V,W) \subset E_{H} \times^{H} W
	\]
	preserved by the $\mathcal{L}$--connection and satisfying
	\[
	F(V,W)|_{\mathfrak{p}} = \{\mathfrak{b}\} \times V .
	\]
	
	Given objects $(V,W)$ and $(V',W')$ and a $G$--equivariant morphism
	$f \colon V \to V'$, consider the $G$--invariant subspace
	$L \subset W \oplus W'$ given by the graph of $f$. The assumption that the
	monodromy is contained in $G$ implies the existence of a subbundle
	\[
	L(f) \subset F(V,W) \oplus F(V',W')
	\]
	preserved by the $\mathcal{L}$--connection and satisfying
	$L(f)|_{\mathfrak{p}} = L$. This subbundle defines the graph of a morphism
	\[
	f^{\sharp} \colon F(V,W) \longrightarrow F(V',W'),
	\]
	which satisfies $f^{\sharp}|_{\mathfrak{p}} = f$. The uniqueness of
	$f^{\sharp}$ follows from \cite[\S\,4, Lemma~4.9]{CT.Simpson-1994}.
	
	In this way, we obtain a strict, exact, faithful tensor functor from
	$\texttt{Rep}(G,H)$ to the category of vector bundles on $X$ equipped with
	integrable $\mathcal{L}$--connections. Composing with the inverse of the above
	equivalence of tensor categories yields a strict, exact, faithful tensor
	functor from $\texttt{Rep}(G)$ to the category of vector bundles with integrable
	$\mathcal{L}$--connections on $X$. By
	Lemma~\eqref{lem:PB-with-L-conection-Tannakian equivalence}, this functor
	corresponds to a principal $G$--bundle equipped with an integrable
	$\mathcal{L}$--connection, as required.
\end{proof}




\begin{lem}\label{lem:Fr-bundle-induce-L-connection}
	Let $E$ be a vector bundle of rank $n$ on $X$ equipped with an integrable
	$\mathcal{L}$--connection, and let
	\[
	\beta \colon E|_{\mathfrak{p}} \xrightarrow{\;\sim\;} \mathbb{C}^{n}
	\]
	be a choice of framing at a closed point $\mathfrak{p} \in X$.
	Then the frame bundle $\mathbf{Fr}(E)$ carries a natural structure of a
	principal $GL(n,\mathbb{C})$--bundle with integrable $\mathcal{L}$--connection.
	Moreover, the vector bundle $E$ is recovered as the associated bundle
	\[
	E \;=\; \mathbf{Fr}(E) \times^{GL(n,\mathbb{C})} \mathbb{C}^{n}.
	\]
	This construction induces a bijection between the sets of isomorphism classes
	of pairs $(E,\beta)$ and $(\mathbf{Fr}(E),\mathfrak{b})$, where
	$\mathfrak{b} \in \mathbf{Fr}(E)|_{\mathfrak{p}}$ corresponds to the
	framing $\beta$.
\end{lem}

\begin{proof}
	Let $\texttt{Rep}(GL(n,\mathbb{C}),\mathrm{std})$ denote the category whose
	objects are pairs $(V,T^{a,b}(\mathbb{C}^{n}))$, where
	\[
	T^{a,b}(\mathbb{C}^{n})
	:= (\mathbb{C}^{n})^{\otimes a} \otimes ((\mathbb{C}^{n})^{*})^{\otimes b},
	\]
	and $V \subset T^{a,b}(\mathbb{C}^{n})$ is a $GL(n,\mathbb{C})$--invariant
	subspace. Morphisms are $GL(n,\mathbb{C})$--equivariant linear maps between such
	subspaces. This category is tensor--equivalent to
	$\texttt{Rep}(GL(n,\mathbb{C}))$.
	
	Let $V \subset T^{a,b}(\mathbb{C}^{n})$ be a $GL(n,\mathbb{C})$--invariant
	subspace. For any $n$--dimensional vector space $U$, the functoriality of tensor
	constructions yields a canonically defined subspace
	\[
	V \subset T^{a,b}(U),
	\]
	independent of the choice of basis of $U$.
	Applying this construction fiberwise to the vector bundle $E$, we obtain a
	subbundle
	\[
	F \subset T^{a,b}(E)
	\]
	such that $\beta(F|_{\mathfrak{p}}) = V$.
	
	Since $E$ is equipped with an integrable $\mathcal{L}$--connection and tensor
	operations preserve integrable $\mathcal{L}$--connections, the tensor bundle
	$T^{a,b}(E)$ carries a natural integrable $\mathcal{L}$--connection. Moreover,
	the construction of $F$ is compatible with infinitesimal automorphisms; hence
	$F$ is preserved by the $\mathcal{L}$--connection. Choosing a
	$GL(n,\mathbb{C})$--invariant complement $V^{\perp}$ of $V$ yields a
	corresponding complementary subbundle $F^{\perp}$, which is also preserved by
	the $\mathcal{L}$--connection. Consequently, $F$ is a vector bundle equipped
	with an integrable $\mathcal{L}$--connection.
	
	Morphisms in $\texttt{Rep}(GL(n,\mathbb{C}),\mathrm{std})$ induce morphisms
	between the corresponding vector bundles with integrable
	$\mathcal{L}$--connections, and these constructions are compatible with tensor
	products. Hence, we obtain a strict, exact, faithful tensor functor
	\[
	\texttt{Rep}(GL(n,\mathbb{C}),\mathrm{std})
	\longrightarrow
	\{\text{vector bundles with integrable $\mathcal{L}$--connections on } X\}.
	\]
	Composing with the inverse of the above tensor equivalence yields a strict,
	exact, faithful tensor functor from
	$\texttt{Rep}(GL(n,\mathbb{C}))$ to the category of vector bundles with
	integrable $\mathcal{L}$--connections on $X$. By
	Lemma~\eqref{lem:PB-with-L-conection-Tannakian equivalence}, this functor
	corresponds to a principal $GL(n,\mathbb{C})$--bundle with integrable
	$\mathcal{L}$--connection, which is precisely the frame bundle $\mathbf{Fr}(E)$.
\end{proof}



Fix a faithful representation $\rho \colon G \hookrightarrow GL(V)$ and  $\mathcal{P}_{0}$ be a Hilbert polynomial of $\mathcal{O}_{X}$.
Let $\mathfrak{p} \in X$ be a fixed closed point.
We define a contravariant functor
\[
\mathcal{M}^{\text{DR}^{\natural}}_{\mathcal{L}}(X,\mathfrak{p},G)
\colon \texttt{Sch}/\mathbb{C} \longrightarrow \texttt{Sets}
\]
as follows.

\begin{itemize}
	\item
	To a scheme $S$, the functor associates the set of isomorphism classes of
	triples $(\mathcal{F},\nabla_{\mathcal{L}},b)$, where:
	\begin{itemize}
		\item
		$\mathcal{F}$ is a principal $G$--bundle on $X_S := X \times S$,
		\item
		$\nabla_{\mathcal{L}}$ is an integrable $\mathcal{L}$--connection on
		$\mathcal{F}$,
		\item
		$b \colon S \to \mathcal{F}|_{\mathfrak{p}\times S}$ is a framing,
	\end{itemize}
	such that the associated vector bundle
	\[
	\mathcal{F}(V) := \mathcal{F} \times^{G} V
	\]
	is flat over $S$ and, for every geometric point $s \in S$,
	the fiber $\mathcal{F}_s(V)$ has Hilbert polynomial $\dim(V)\mathcal{P}_{0}$.
	
	\item for every closed point $s\in S$, the fiber $(E_{G,s}, \nabla_{\mathcal{L},s},b_{s})$ is a principal $G$--bundle with integrable $\mathcal{L}$--connection and $\mathfrak{b}=b_{s}$.
	
	\item
	Two such triples $(\mathcal{F},\nabla_{\mathcal{L}},b)$ and
	$(\mathcal{F}',\nabla'_{\mathcal{L}},b')$ over $S$ are said to be
	\emph{isomorphic} if there exists an isomorphism of principal $G$--bundles
	with $\mathcal{L}$--connections
	\[
	\varphi \colon (\mathcal{F},\nabla_{\mathcal{L}})
	\longrightarrow
	(\mathcal{F}',\nabla'_{\mathcal{L}})
	\]
	such that $\varphi \circ b = b'$.
	
	\item
	For a morphism of schemes $\psi \colon T \to S$, the functor assigns the pullback
	\[
	(\mathcal{F},\nabla_{\mathcal{L}},b)
	\longmapsto
	(\psi^{*}\mathcal{F}, \psi^{*}\nabla_{\mathcal{L}}, \psi^{*}b),
	\]
	which defines a $T$--family on $X \times T$.
\end{itemize}




The following theorem is well known for principal $G$--bundles with integrable $\mathcal{T}_{X}$--connections (cf.~\cite[\S\,9, Th.~9.10]{CT.Simpson-1995}). We extend it here to the
Lie algebroid setting.

\begin{thm}\label{thm: L-twisted - reprsentaion - space for L-connection}
	Let $\mathfrak{p}$ be a fixed base point of $X$. Then there exists a scheme $\mathcal{R}^{\text{DR}}_{\mathcal{L}}(X,\mathfrak{p},G)$
	representing the functor
	$\mathcal{M}^{\text{DR}^{\natural}}_{\mathcal{L}}(X,\mathfrak{p},G)$.
	Moreover, if $f \colon G \hookrightarrow H$ is a closed embedding of algebraic groups, then $f$ induces a closed embedding
	\[
	\mathcal{R}^{\text{DR}}_{\mathcal{L}}(X,\mathfrak{p},G)
	\hookrightarrow
	\mathcal{R}^{\text{DR}}_{\mathcal{L}}(X,\mathfrak{p},H).
	\]
\end{thm}

\begin{proof}
	We first treat the case $G = GL(n,\mathbb{C})$.
	By Lemma~\eqref{lem:Fr-bundle-induce-L-connection}, framed vector bundles with integrable $\mathcal{L}$--connections are equivalent to  framed principal $GL(n,\mathbb{C})$--bundles with integrable $\mathcal{L}$--connections. Hence, by Theorem~\eqref{eqn: moduli-of-L-connection on VB}, the functor
	$\mathcal{M}^{\text{DR}^{\natural}}_{\mathcal{L}}(X,\mathfrak{p},GL(n,\mathbb{C}))$ is represented by the scheme
	\[
	\mathcal{R}^{\text{DR}}_{\mathcal{L}}(X,\mathfrak{p},GL(n,\mathbb{C}))
	:= \mathcal{R}^{\text{DR}}_{\mathcal{L}}(X,\mathfrak{p},n).
	\]
	
	Now let $G \subset H$ be a closed algebraic subgroup, and assume that the scheme $\mathcal{R}^{\text{DR}}_{\mathcal{L}}(X,\mathfrak{p},H)$ representing
	$\mathcal{M}^{\text{DR}^{\natural}}_{\mathcal{L}}(X,\mathfrak{p},H)$ is already constructed.
	Let $
	(\mathcal{F}^{\mathrm{univ}}, \nabla^{\text{DR}^{\mathrm{univ}}}_{\mathcal{L}},b^{\mathrm{univ}})
	$
	denote the universal framed principal $H$--bundle with integrable
	$\mathcal{L}$--connection on
	$X \times \mathcal{R}^{\text{DR}}_{\mathcal{L}}(X,\mathfrak{p},H)$.
	
	Let $V$ be a finite--dimensional representation of $H$, and let
	$W \subset V$ be a subspace preserved by $G$.
	Set
	\[
	E^{\mathrm{univ}} := \mathcal{F}^{\mathrm{univ}} \times^{H} V,
	\]
	which is a vector bundle with integrable $\mathcal{L}$--connection.
	The framing $b^{\mathrm{univ}}$ induces an isomorphism
	\[
	\beta^{\mathrm{univ}} :
	E^{\mathrm{univ}}|_{\mathfrak{p} \times \mathcal{R}^{\text{DR}}_{\mathcal{L}}(H)}
	\;\xrightarrow{\;\sim\;}\;
	\mathcal{V} := V \otimes_{\mathbb{C}}
	\mathcal{O}_{\mathcal{R}^{\text{DR}}_{\mathcal{L}}(H)}.
	\]
	
	Let $k := \dim(V) - \dim(W)$, and denote by
	$\mathcal{W} := W \otimes_{\mathbb{C}}
	\mathcal{O}_{\mathcal{R}_{\mathcal{L}}(H)}$
	the corresponding trivial subbundle of $\mathcal{V}$.
	The quotient $\mathcal{V} \to \mathcal{V}/\mathcal{W}$ defines a section
	\[
	\sigma_{\mathcal{V}/\mathcal{W}} :
	\mathcal{R}^{\text{DR}}_{\mathcal{L}}(H)
	\longrightarrow
	\mathfrak{Grass}_{\mathcal{R}^{\text{DR}}_{\mathcal{L}}(H)}(\mathcal{V},k).
	\]
	
	By Lemma~\eqref{lem:projective-scheme-N(E,k)}, there exists a
	closed subscheme
	\[
	\mathcal{N}(E^{\mathrm{univ}},k)
	\subset
	\mathfrak{Grass}_{\mathcal{R}^{\text{DR}}_{\mathcal{L}}(H)}
	\bigl(E^{\mathrm{univ}}|_{\mathfrak{p}\times
		\mathcal{R}^{\text{DR}}_{\mathcal{L}}(H)}, k\bigr),
	\]
	parametrizing quotients compatible with the integrable
	$\mathcal{L}$--connection.
	Using the framing $\beta^{\mathrm{univ}}$, we view
	$\mathcal{N}(E^{\mathrm{univ}},k)$ as a closed subscheme of
	$\mathfrak{Grass}_{\mathcal{R}^{\text{DR}}_{\mathcal{L}}(H)}(\mathcal{V},k)$.
	Define
	\[
	\mathbf{C}(V,W)
	:=
	\sigma_{\mathcal{V}/\mathcal{W}}^{-1}
	\bigl(\mathcal{N}(E^{\mathrm{univ}},k)\bigr)
	\;\subset\;
	\mathcal{R}^{\text{DR}}_{\mathcal{L}}(H).
	\]
	
	By construction, a morphism $g \colon S \to \mathcal{R}^{\text{DR}}_{\mathcal{L}}(H)$ factors
	through $\mathbf{C}(V,W)$ if and only if the pullback
	$g^{*}(E^{\mathrm{univ}})$ admits a strict subbundle preserved by the
	$\mathcal{L}$--connection whose fiber at $\mathfrak{p}$ corresponds to $W$.
	Equivalently, this condition expresses that the monodromy of the framed object
	$g^{*}(\mathcal{F}^{\mathrm{univ}}, b^{\mathrm{univ}})$ preserves $W$.
	
	Finally, define
	\[
	\mathcal{R}^{\text{DR}}_{\mathcal{L}}(X,\mathfrak{p},G)
	:=
	\bigcap_{(V,W)} \mathbf{C}(V,W),
	\]
	where the intersection is taken over all representations $V$ of $H$ and all $G$--invariant subspaces $W \subset V$.
	This is a closed subscheme of
	$\mathcal{R}^{\text{DR}}_{\mathcal{L}}(X,\mathfrak{p},H)$.
	By construction, it represents the functor assigning to a scheme $S$ the set of framed principal $H$--bundles with integrable $\mathcal{L}$--connection whose monodromy is contained in $G$.
	By Lemma~\eqref{monodromy related thm-for-PB with-L-connection}, this functor identifies with $\mathcal{M}^{\text{DR}^{\natural}}_{\mathcal{L}}(X,\mathfrak{p},G)$.
	
	The final assertion concerning closed embeddings follows immediately from the
	construction.
\end{proof}

\medskip

\noindent\textbf{GIT quotient of the framed representation space.}
By Theorem~\ref{thm: L-twisted - reprsentaion - space for L-connection}, for any
closed embedding $G \subset GL(n,\mathbb{C})$ we obtain a closed immersion
\[
\mathcal{R}^{\text{DR}}_{\mathcal{L}}(X,\mathfrak{p},G)
\hookrightarrow
\mathcal{R}^{\text{DR}}_{\mathcal{L}}(X,\mathfrak{p},GL(n,\mathbb{C}))
=
\mathcal{R}^{\text{DR}}_{\mathcal{L}}(X,\mathfrak{p},n).
\]
Assume henceforth that $G$ is a reductive algebraic group and that
$G \hookrightarrow GL(n,\mathbb{C})$ is a faithful representation.

The group $GL(n,\mathbb{C})$ acts naturally on
$\mathcal{R}^{\text{DR}}_{\mathcal{L}}(X,\mathfrak{p},n)$ by change of framing at the base
point $\mathfrak{p}$, and this action restricts to an action of $G$ on the closed
subscheme $\mathcal{R}^{\text{DR}}_{\mathcal{L}}(X,\mathfrak{p},G)$.

Choose a $GL(n,\mathbb{C})$--linearized ample line bundle
$\mathscr{L}$ on $\mathcal{R}^{\text{DR}}_{\mathcal{L}}(X,\mathfrak{p},n)$ such that every
point is semistable for the action of $GL(n,\mathbb{C})$
(cf.~\cite[\S\,4, Th.~4.10]{CT.Simpson-1994}).
By Mumford’s numerical criterion for semistability
(cf.~\cite[\S\,Ch.~2, Th.~2.1]{D.Mumford-1965},
\cite[\S\,4, Th.~4.2.11]{Huybrechts-Lehn-2010}),
every point is also semistable for the induced action of the reductive subgroup $G$.
Therefore, all points of
$\mathcal{R}^{\text{DR}}_{\mathcal{L}}(X,\mathfrak{p},G)$ are $G$--semistable with respect to the restricted linearization
$\mathscr{L}|_{\mathcal{R}_{\mathcal{L}}(X,\mathfrak{p},G)}$.

Since $G$ is reductive, GIT yields a good quotient
\[
\mathcal{M}^{\text{DR}}_{\mathcal{L}}(X,G)
:=
\mathcal{R}^{\text{DR}}_{\mathcal{L}}(X,\mathfrak{p},G)\sslash G,
\]
which is a quasi-projective variety
(cf.~\cite[\S\,Ch.~1, Th.~1.10]{D.Mumford-1965}).




\begin{proposition}\label{prop:local-quotient-description}
	The moduli functor $\mathcal{M}^{\text{DR}^{\natural}}_{\mathcal{L}}(X,G)$ is locally
	isomorphic, in the étale topology, to the quotient functor associated with the natural action of $G$ on the framed representation space
	$\mathcal{R}^{\text{DR}}_{\mathcal{L}}(X,\mathfrak{p},G)$.
\end{proposition}

\begin{proof}
	Let $S$ be a scheme and let
	\[
	(\mathcal{E}_{G},\nabla_{\mathcal{L}})
	\in
	\mathcal{M}^{\text{DR}^{\natural}}_{\mathcal{L}}(X,G)(S)
	\]
	be an $S$--family of principal $G$--bundles with integrable
	$\mathcal{L}$--connection.
	
	By definition of the moduli functor, after passing to an étale cover
	$\{S_i \to S\}$, the restriction of $\mathcal{E}_{G}$ to
	$\mathfrak{p}\times S_i$ admits a trivialization. Hence, over each $S_i$,
	the family lifts to a framed object
	\[
	(\mathcal{E}_{G},\nabla_{\mathcal{L}},\beta_i)
	\in
	\mathcal{M}^{\text{DR}^{\natural}}_{\mathcal{L}}(X,\mathfrak{p},G)(S_i).
	\]
	
	Since the framed moduli functor
	$\mathcal{M}^{\text{DR}^{\natural}}_{\mathcal{L}}(X,\mathfrak{p},G)$ is represented by $\mathcal{R}^{\text{DR}}_{\mathcal{L}}(X,\mathfrak{p},G)$, each such framed family corresponds uniquely to a morphism
	\[
	S_i \longrightarrow \mathcal{R}^{\text{DR}}_{\mathcal{L}}(X,\mathfrak{p},G).
	\]
	
	On overlaps $S_{ij}=S_i\times_S S_j$, two choices of framing differ by the
	action of an element of $G$. Consequently, the induced morphisms to
	$\mathcal{R}^{\text{DR}}_{\mathcal{L}}(X,\mathfrak{p},G)$ agree up to the natural $G$--action. This shows that $\mathcal{M}^{\text{DR}^{\natural}}_{\mathcal{L}}(X,G)$
	is locally isomorphic to the quotient functor
	$\mathcal{R}^{\text{DR}}_{\mathcal{L}}(X,\mathfrak{p},G)/G$ in the étale topology.
\end{proof}

Hence, we have the following main theorem,
\begin{corollary}\label{cor:moduli of int-L-connection PB}
	The good GIT quotient $
	\mathcal{M}^{\text{DR}}_{\mathcal{L}}(X,G)
	=
	\mathcal{R}^{\text{DR}}_{\mathcal{L}}(X,\mathfrak{p},G)\sslash G
	$
	universally corepresents the functor
	$\mathcal{M}^{\text{DR}^{\natural}}_{\mathcal{L}}(X,G)$. In particular, $\mathcal{M}^{\text{DR}}_{\mathcal{L}}(X,G)$ is a quasi-projective variety over $\mathbb{C}$.
\end{corollary}




%\begin{thm}\label{thm:closed-orbits}
%	Suppose $G$ is a reductive algebraic group and 
	%$f \colon G \hookrightarrow GL(n,\mathbb{C})$ is a faithful representation.
	%Let $(E_G,b) \in \mathcal{R}_{\mathcal{L}}(X,\mathfrak{p},G)$ and let
	%$(E,\beta) \in \mathcal{R}_{\mathcal{L}}(X,\mathfrak{p},n)$ be the induced point
	%via extension of structure group by $f$.
	%Then $(E_G,b)$ lies in a closed $G$--orbit in 
	%$\mathcal{R}_{\mathcal{L}}(X,\mathfrak{p},G)$ if and only if
	%$(E,\beta)$ lies in a closed $GL(n,\mathbb{C})$--orbit in
	%$\mathcal{R}_{\mathcal{L}}(X,\mathfrak{p},n)$.
%\end{thm}


\subsection{The moduli spaces of $\mathcal{L}$-twisted principal $G$-Higgs bundles}

First, we define the moduli functor associated to our moduli problem.
Let $\mathcal{P}_{0}$ be a Hilbert polynomial of $\mathcal{O}_{X}$.
We consider the contravariant functor
\[
\mathcal{M}^{\mathrm{Dol}^{\natural}}_{\mathcal{L}}(X,G)
\colon (\texttt{Sch}/\mathbb{C}) \longrightarrow \texttt{Sets},
\]
defined as follows.

\begin{itemize}\label{moduli-of-L-twisted princpal G--Higgs bundles}
	\item
	For a $\mathbb{C}$--scheme $S$, the set
	$\mathcal{M}^{\mathrm{Dol}^{\natural}}_{\mathcal{L}}(X,G)(S)$
	consists of isomorphism classes of $\mathcal{L}$--twisted principal $G$--Higgs bundles  $(\mathcal{F},\theta_{\mathcal{L}})$ of harmonic type on $X'=X\times S$, such that the following conditions hold:
    \item[(i)]
 	For every finite--dimensional representation
 	$\rho \colon G \to GL(V)$,
 	the associated $\mathcal{L}$--twisted Higgs bundle
 	\[
 	(\mathcal{F}(V),\theta_{\mathcal{L}}(V))
 	\]
 	on $X_S$ is $S$--flat and for every closed point $\bar{s}\rightarrow S$ the restriction $E_{G}(V)_{\bar{s}}$ has Hilbert polynomial $\dim(V)\mathcal{P}_{0}$.
 	
 	\item[(ii)]
 	For every closed point $s \in S$,, the fiber $(E_{G,s}, \theta_{\mathcal L,s})$ is a $\mathcal{L}$--twisted principal $G$--Higgs bundle of harmonic type on $X$.
 	
 \item Two such family $(\mathcal{F},\theta_{\mathcal{L}})$ and $(\mathcal{F}',\theta'_{\mathcal L})$ over $S$ are said to be \textit{isomorphic} if there exist an isomorphism of $\mathcal{L}$--twisted principal $G$--Higgs bundle, \[\varphi:(\mathcal{F},\theta_{\mathcal L}\rightarrow (\mathcal{F}',\theta'_{\mathcal L}))\]
 
 \item For a morphism of schemes $\psi:T\rightarrow S$, the functor assigns the pullback \[(\mathcal{F}, \theta_{\mathcal L})\mapsto (\psi^{*}(\mathcal{F}),\psi^{*}(\theta_{\mathcal L}))\]which defines a $T$--family on $X\times T$.
 \end{itemize}
 

We now construct a $\mathcal{L}$--twisted Dolbeault moduli space $\mathcal{M}^{\mathrm{Dol}}_{\mathcal{L}}(X,G)$ as a quasi-projective variety over $\mathbb{C}$, by an argument entirely parallel to the construction of $\mathcal{M}^{\text{DR}}_{\mathcal{L}}(X,G)$ in the de Rham case.This variety universally corepresents the functor $\mathcal{M}^{\mathrm{Dol},\natural}_{\mathcal{L}}(X,G)$.



The following lemma is well known in the case of Higgs bundles of
semiharmonic type on $X$ (see \cite[\S\,9, Lemma~9.3]{CT.Simpson-1995}).
We extend it here to the setting of Lie algebroids by the same method.

\begin{lem}\label{projective scheme N(E,k) for L-twisted-Higgs bundles}
	Let $\mathcal{L}$ be a Lie algebroid on $X$, and let
	$E$ be an $\mathcal{L}$--twisted Higgs bundle of semiharmonic type.
	Fix an integer $k$.
	Then there exists a projective scheme
	$\mathcal{N}(E,k)$ over $\mathbb{C}$ representing the functor which
	associates to each $\mathbb{C}$--scheme $S$ the set of quotients
	\[
	p_X^{*}E \longrightarrow F \longrightarrow 0
	\]
	of $\mathcal{L}$--twisted Higgs bundles of semiharmonic type on
	$X_S := X \times S$, where $F$ has rank $k$.
	Moreover, the natural morphism
	\[
	\mathcal{N}(E,k)(S) \hookrightarrow
	\mathfrak{Grass}_{S}(E|_{\mathfrak{p}\times S},k)
	\]
	is a closed embedding.
\end{lem}

\begin{proof}
	Let $\Lambda^{\text{Dol}}_{\mathcal{L}}$ denote the $\mathcal{L}$--twisted Dolbeault $D$--algebara  associated to the Lie algebroid $\mathcal{L}$.
	By lemma~\eqref{lem: Lambda module structure--equivalent to L twisted Higgs bundle}
	the $\mathcal{L}$--twisted Higgs bundle $E$ may be viewed as a
	$\Lambda^{\mathrm{Dol}}_{\mathcal{L}}$--module.
	
	Let $\mathcal{P}_0$ denote the Hilbert polynomial of $\mathcal{O}_X$.
	The Quot scheme $\mathfrak{Quot}(E,k\mathcal{P}_0)$
	parame-trizes quotient sheaves
	$E \twoheadrightarrow F \to 0$
	that are flat over $\mathbb{C}$ and have Hilbert polynomial
	$k\mathcal{P}_0$.
	Let $e \colon S \to \mathfrak{Quot}(E,k\mathcal{P}_0)$ be an $S$--valued
	point corresponding to a quotient
	$p_X^{*}E \twoheadrightarrow F \to 0$
	on $X_S$, and let $K$ denote its kernel.
	
	Let
	\[
	E^{\mathrm{univ}} \twoheadrightarrow F^{\mathrm{univ}} \to 0
	\]
	be the universal quotient on
	$X^{\mathrm{univ}} := X \times \mathfrak{Quot}(E,k\mathcal{P}_0)$,
	and let $K^{\mathrm{univ}} \subset E^{\mathrm{univ}}$ be its kernel.
	Then $K = e^{*}(K^{\mathrm{univ}})$.
	The Higgs field on $E$ induces a universal morphism
	\[
	\theta_{\mathcal{L}}^{\mathrm{univ}} \colon
	K^{\mathrm{univ}} \longrightarrow
	F^{\mathrm{univ}} \otimes (\Omega^1_{\mathcal{L}})^{\mathrm{univ}},
	\]
	whose pullback via $e$ is the induced morphism
	\[
	\theta_{\mathcal{L}} \colon K \longrightarrow F \otimes \Omega^1_{\mathcal{L}}.
	\]
	
	By Lemma~\eqref{lem:fact-from-simpson-thm-3-8}, there exists a closed
	subscheme
	\[
	\mathcal{N}(E,k) \subset \mathfrak{Quot}(E,k\mathcal{P}_0)
	\]
	representing the condition that for any $S$--valued point
	$e \colon S \to \mathcal{N}(E,k)$,
	the induced morphism $\theta$ vanishes.
	In this case, the $\mathcal{L}$--twisted Higgs field on $p_X^{*}E$ descends to a morphism
	\[
	\theta_{\mathcal{L}} \colon
	F \longrightarrow F \otimes \Omega^1_{\mathcal{L}},
	\]
	making $F$ into an $\mathcal{L}$--twisted Higgs sheaf.
	
	Thus $\mathcal{N}(E,k)(S)$ parametrizes quotients of
	$\mathcal{L}$--twisted Higgs sheaves
	$p_X^{*}E \twoheadrightarrow F \to 0$
	on $X_S$, where $F$ is flat over $S$ with Hilbert polynomial
	$k\mathcal{P}_0$.
	For any $s \in S$, the fiber $F_s$ is a quotient $\mathcal{L}$--twisted Higgs sheaf of $E_s$
	with the same normalized Hilbert polynomial.
	Let $K_s$ denote the kernel of $E_s \to F_s$.
	Then $K_s$ is a $\mathcal{L}$--twisted Higgs subsheaf of $E_s$ with the same normalized Hilbert
	polynomial.
	By proposition \eqref{prop:locally-free-L-Higgs},
	$K_s$ is a strict subbundle with vanishing $\mathcal{L}$--Chern classes, and hence
	$F_s$ is locally free with vanishing $\mathcal{L}$--Chern classes.
	By \cite[\S\,1, Lemma~1.27]{CT.Simpson-1994}, it follows that $F$ is
	locally free.
	
	Therefore, $\mathcal{N}(E,k)(S)$ parametrizes quotients of
	$\mathcal{L}$--twisted Higgs bundles of semiharmonic type
	$p_X^{*}E \twoheadrightarrow F \to 0$
	on $X_S$ of rank $k$.
	
	Finally, $\mathcal{N}(E,k)$ is projective over $\mathbb{C}$, and the
	natural morphism to
	$\mathfrak{Grass}_{S}(E|_{\mathfrak{p}\times S},k)$
	is injective on $S$--valued points by
	\cite[\S\,4, Lemma~4.9]{CT.Simpson-1994}. Hence this morphism is a closed embedding.
\end{proof}


We now extend the notion of the monodromy group for principal $G$--objects
(cf.~\cite[\S\,9, p.~50]{CT.Simpson-1995}) to the setting of
$\mathcal{L}$--twisted principal Higgs bundles.

Let $\mathfrak{p} \in X$ be a fixed closed point, and let $G \subset H$
be an algebraic subgroup.
Let $E_H$ be a $\mathcal{L}$--twisted principal Higgs bundle for the group
$H$ on $X$, and fix a point
$\mathfrak{b} \in E_H|_{\mathfrak{p}} \cong H$.

We say that the \emph{monodromy of the pair $(E_H,\mathfrak{b})$ is
contained in $G$} if the following condition holds:
for every finite-dimensional linear representation $V$ of $H$, and every
subspace $W \subset V$ preserved by $G$, there exists a strict
$\mathcal{L}$--twisted Higgs subbundle of semiharmonic type
\[
F \subset E_H \times^{H} V
\]
such that
\[
F|_{\mathfrak{p}}
=
\{\mathfrak{b}\} \times W
\subset
(E_H \times^{H} V)|_{\mathfrak{p}}.
\]

The \emph{monodromy group} $\mathbf{Mono}(E_H,\mathfrak{b})$ is defined as
the intersection of all algebraic subgroups $G \subset H$ such that the
monodromy of $(E_H,\mathfrak{b})$ is contained in $G$.


The following two lemmas are well known in the case of principal Higgs
bundles of semiharmonic type (see
\cite[\S\,9, Lemmas~9.4 and~9.5]{CT.Simpson-1995}).
Their proofs carry over verbatim to the $\mathcal{L}$--twisted setting,
using the arguments of
Lemma~\eqref{monodromy related thm-for-PB with-L-connection} and
Lemma~\eqref{lem:Fr-bundle-induce-L-connection}.

\begin{lem}\label{monodromy related thm- with-L-twisted-Higgs-bundle-for-PB}
	Let $G \subset H$ be an algebraic subgroup, and let
	$E_G$ be a $\mathcal{L}$--twisted principal Higgs bundle on $X$ for the
	group $G$.
	Then the induced bundle
	\[
	E_H := E_G \times^{G} H
	\]
	has a natural structure of a $\mathcal{L}$--twisted principal Higgs
	bundle for the group $H$.
	
	Moreover, this construction induces a bijection between:
	\begin{enumerate}
		\item
		isomorphism classes of pairs $(E_G,\mathfrak{b}')$, where $E_G$ is a
		$\mathcal{L}$--twisted principal Higgs bundle of semiharmonic type for
		$G$ and $\mathfrak{b}' \in E_G|_{\mathfrak{p}}$;
		
		\item
		isomorphism classes of pairs $(E_H,\mathfrak{b})$, where $E_H$ is a
		$\mathcal{L}$--twisted principal Higgs bundle of semiharmonic type for
		$H$ and $\mathfrak{b} \in E_H|_{\mathfrak{p}}$, such that the monodromy of
		$(E_H,\mathfrak{b})$ is contained in $G$.
	\end{enumerate}
\end{lem}

\begin{proof}
	The proof is identical to that of
	Lemma~\eqref{monodromy related thm-for-PB with-L-connection}.
\end{proof}

\begin{lem}\label{lem:Fr-bundle as - L-twisted-principal-GLn - bundle}
	Let $E$ be a $\mathcal{L}$--twisted Higgs bundle of semiharmonic type of
	rank $n$ on $X$, and let
	$\beta \colon E|_{\mathfrak{p}} \xrightarrow{\sim} \mathbb{C}^{n}$ be a
	framing.
	Then the frame bundle $\mathbf{Fr}(E)$ carries a natural structure of a
	$\mathcal{L}$--twisted principal Higgs bundle for the group
	$GL(n,\mathbb{C})$.
	
	Moreover, the associated vector bundle recovers $E$:
	\[
	E \cong \mathbf{Fr}(E) \times^{GL(n,\mathbb{C})} \mathbb{C}^{n}.
	\]
	This construction induces a bijection between isomorphism classes of
	pairs $(E,\beta)$ and $(\mathbf{Fr}(E),\mathfrak{b})$, where
	$\mathfrak{b} \in \mathbf{Fr}(E)|_{\mathfrak{p}}$ corresponds to the
	chosen frame $\beta$.
\end{lem}

\begin{proof}
	The proof is identical to that of
	Lemma~\eqref{lem:Fr-bundle-induce-L-connection}.
\end{proof}


Fix a irreducible complex projective variety $X$ with transitive a Lie algebroid $\mathcal{L}$ on $X$, a closed point $\mathfrak{p} \in X$, and a
reductive algebraic group $G$.
Fix once and for all a faithful representation
$\rho \colon G \hookrightarrow GL(V)$.

We define a contravariant functor
\[
\mathcal{M}^{\mathrm{Dol}^{\natural}}_{\mathcal{L}}(X,\mathfrak{p},G)
\colon \mathbf{Sch}/\mathbb{C} \longrightarrow \mathbf{Sets}
\]
as follows.


For a scheme $S$, the set
$\mathcal{M}^{\mathrm{Dol}^{\natural}}_{\mathcal{L}}(X,\mathfrak{p},G)(S)$
consists of isomorphism classes of triples
\[
(\mathcal{F},\theta_{\mathcal{L}}, b),
\]
where:
\begin{itemize}
	\item
	$(\mathcal{F}, \theta_{\mathcal{L}})$ is a $\mathcal{L}$--twisted principal $G$--Higgs bundle of semi-harmonic type on $X_S := X \times S$;
	\item
	$b \colon S \to \mathcal{F}|_{\mathfrak{p}\times S}$ is a framing at the base point;
	
	\item
	for every closed point $s \in S$, the induced Higgs bundle
	\[
	(E_s,\theta_s,b_{s})
	:=
	\bigl(
	\mathcal{F}_s \times^{G} V,\,
	\theta_{\mathcal{L},s,}\mathfrak{b}
	\bigr)
	\]
	is a semiharmonic $\mathcal{L}$--twisted Higgs bundle on $X$ with fixed Hilbert polynomial $\mathcal{P}_{0}$.
\end{itemize}
Two triples $(\mathcal{F},\theta_{\mathcal{L}},b)$ and
$(\mathcal{F}',\theta_{\mathcal{L}}',b')$ over $S$ are isomorphic if there
exists an isomorphism of principal $G$--bundles
\[
\varphi \colon \mathcal{F} \xrightarrow{\sim} \mathcal{F}'
\]
such that:
\begin{itemize}
	\item
	$\varphi$ is compatible with the $\mathcal{L}$--twisted Higgs fields,
	i.e. $\varphi^*(\theta_{\mathcal{L}}') = \theta_{\mathcal{L}}$;
	
	\item
	$\varphi$ preserves the framing:
	$\varphi \circ b = b'$.
\end{itemize}

The following theorem is well known for principal $G$--Higgs bundles of
semiharmonic type (cf.~\cite[\S\,9, Th.~9.6]{CT.Simpson-1995}). We extend it
to the setting of $\mathcal{L}$--twisted principal Higgs bundles, following the same strategy as in
Theorem~\eqref{thm: L-twisted - reprsentaion - space for L-connection}.

\begin{thm}\label{thm:Dol-L-representation-space}
	Suppose $\mathfrak{p}$ is a fixed base point in $X$. Then there exists a scheme $
	\mathcal{R}^{\mathrm{Dol}}_{\mathcal{L}}(X,\mathfrak{p},G)$
	representing the functor
	$\mathcal{M}^{\mathrm{Dol}^{\natural}}_{\mathcal{L}}(X,\mathfrak{p},G)$.
	Moreover, if $f \colon G \hookrightarrow H$ is a closed embedding of
	algebraic groups, then $f$ induces a closed embedding
	\[
	\mathcal{R}^{\mathrm{Dol}}_{\mathcal{L}}(X,\mathfrak{p},G)
	\hookrightarrow
	\mathcal{R}^{\mathrm{Dol}}_{\mathcal{L}}(X,\mathfrak{p},H).
	\]
\end{thm}

\begin{proof}
	We first treat the case $G = GL(n,\mathbb{C})$. By
	Lemma~\eqref{lem:Fr-bundle-induce-L-connection} and
	Theorem~\eqref{eqn: moduli-of-L- twisted Higgs bundles}, the functor
	$\mathcal{M}^{\mathrm{Dol}^{\natural}}_{\mathcal{L}}(X,\mathfrak{p},GL(n,\mathbb{C}))$
	is represented by the Dolbeault representation space for framed
	$\mathcal{L}$--twisted Higgs bundles of rank $n$. We therefore set
	\[
	\mathcal{R}^{\mathrm{Dol}}_{\mathcal{L}}(X,\mathfrak{p},GL(n,\mathbb{C}))
	:= \mathcal{R}^{\mathrm{Dol}}_{\mathcal{L}}(X,\mathfrak{p},n),
	\]
	which represents the desired functor in this case.
	
	Now let $G$ be an arbitrary reductive group. Choose a faithful
	representation $G \hookrightarrow GL(n,\mathbb{C})$. By
	Lemma~\eqref{projective scheme N(E,k) for L-twisted-Higgs bundles}, for every representation $V$ of $GL(n,\mathbb{C})$ and every subspace $W \subset V$ preserved by $G$, the condition that a framed
	$\mathcal{L}$--twisted Higgs bundle admits a strict
	$\mathcal{L}$--twisted Higgs subbundle of semiharmonic type with prescribed fiber at $\mathfrak{p}$ is represented by a closed subscheme of $\mathcal{R}^{\mathrm{Dol}}_{\mathcal{L}}(X,\mathfrak{p},n)$.
	
	Intersecting these closed subschemes over all such pairs $(V,W)$, exactly as in the proof of
	Theorem~\eqref{thm: L-twisted - reprsentaion - space for L-connection}, we obtain a closed subscheme
	\[
	\mathcal{R}^{\mathrm{Dol}}_{\mathcal{L}}(X,\mathfrak{p},G)
	\subset
	\mathcal{R}^{\mathrm{Dol}}_{\mathcal{L}}(X,\mathfrak{p},n),
	\]
	which represents the functor
	$\mathcal{M}^{\mathrm{Dol}^{\natural}}_{\mathcal{L}}(X,\mathfrak{p},G)$.
	
	The functoriality with respect to closed embeddings
	$f \colon G \hookrightarrow H$ is immediate from the construction, and
	induces a closed embedding
	\[
	\mathcal{R}^{\mathrm{Dol}}_{\mathcal{L}}(X,\mathfrak{p},G)
	\hookrightarrow
	\mathcal{R}^{\mathrm{Dol}}_{\mathcal{L}}(X,\mathfrak{p},H).
	\]
	This completes the proof.
\end{proof}

\medskip
\noindent \textbf{GIT quotient of the framed $\mathcal{L}$--twisted Dolbeault moduli space}.
We now apply Mumford's Geometric Invariant Theory to construct the
moduli space of $\mathcal{L}$--twisted principal $G$--Higgs bundles.

By Theorem~\ref{thm:Dol-L-representation-space}, for any closed embedding
$G \subset GL(n,\mathbb{C})$ we have a closed subscheme
\[
\mathcal{R}^{\mathrm{Dol}}_{\mathcal{L}}(X,\mathfrak{p},G)
\subset
\mathcal{R}^{\mathrm{Dol}}_{\mathcal{L}}(X,\mathfrak{p},n),
\]
where
$\mathcal{R}^{\mathrm{Dol}}_{\mathcal{L}}(X,\mathfrak{p},GL(n,\mathbb{C}))
= \mathcal{R}^{\mathrm{Dol}}_{\mathcal{L}}(X,\mathfrak{p},n)$.

Assume henceforth that $G$ is reductive and that
$G \hookrightarrow GL(n,\mathbb{C})$ is a faithful representation.
Then $G$ acts on
$\mathcal{R}^{\mathrm{Dol}}_{\mathcal{L}}(X,\mathfrak{p},n)$
via its inclusion in $GL(n,\mathbb{C})$, and this action preserves the
closed subscheme
$\mathcal{R}^{\mathrm{Dol}}_{\mathcal{L}}(X,\mathfrak{p},G)$.


By \cite[\S\,4, Th.~4.10]{CT.Simpson-1994}, there exists a
$GL(n,\mathbb{C})$--linearized ample line bundle
$\mathscr{L}$ on
$\mathcal{R}^{\mathrm{Dol}}_{\mathcal{L}}(X,\mathfrak{p},n)$
such that every point is semistable for the action of $GL(n,\mathbb{C})$.

By the Hilbert--Mumford numerical criterion
(cf.~\cite[\S\,Ch.~2, Th.~2.1]{D.Mumford-1965},
\cite[\S\,4, Th.~4.2.11]{Huybrechts-Lehn-2010}),
every point is then also semistable for the action of the subgroup $G$.
Consequently, every point of
$\mathcal{R}^{\mathrm{Dol}}_{\mathcal{L}}(X,\mathfrak{p},G)$
is $G$--semistable with respect to the induced linearization
$\mathscr{L}|_{\mathcal{R}^{\mathrm{Dol}}_{\mathcal{L}}(X,\mathfrak{p},G)}$.

\begin{thm}\label{thm:moduli-L-twisted-principal-G-Higgs}
	The good GIT quotient $
	\mathcal{M}^{\mathrm{Dol}}_{\mathcal{L}}(X,G)
	=
	\mathcal{R}^{\mathrm{Dol}}_{\mathcal{L}}(X,\mathfrak{p},G)
	\sslash G$ universally corepresents the functor $\mathcal{M}^{\mathrm{Dol}^{\natural}}_{\mathcal{L}}(X,G)$. In particulur, $\mathcal{M}^{\mathrm{Dol}}_{\mathcal{L}}(X,G)$ 
	
\end{thm}

\begin{proof}
	The proof is identical to the $\mathcal{L}$--connection case.
	Étale locally on the base, any family of $\mathcal{L}$--twisted principal
	$G$--Higgs bundles admits a framing at $\mathfrak{p}$, hence is locally
	isomorphic to a $G$--equivariant family over
	$\mathcal{R}^{\mathrm{Dol}}_{\mathcal{L}}(X,\mathfrak{p},G)$.
	Two such local liftings differ by the action of $G$, so the functor
	$\mathcal{M}^{\mathrm{Dol},\natural}_{\mathcal{L}}(X,G)$ is locally
	isomorphic to the quotient functor
	$[\mathcal{R}^{\mathrm{Dol}}_{\mathcal{L}}(X,\mathfrak{p},G)/G]$.
\end{proof}



%\begin{thm}
	%Suppose $G$ is a reductive group and $f:G\rightarrow GL(n,\mathbb{C})$ is a faithful representation. Suppose $(E_{G},b)\in \mathcal{R}^{\text{Dol}}_{\mathcal{L}}(X,\mathfrak{p},G)$ maps to $(E,\beta)\in \mathcal{R}^{\text{Dol}}_{\mathcal{L}}(X,\mathfrak{p},n)$. Then $(E_{G},b)$ is in a closed $G$-orbit in $\mathcal{R}^{\text{Dol}}_{\mathcal{L}}(X,\mathfrak{p},G)$ if and only if $(E,\beta)$ is in a closed $GL(n,\mathbb{C})$-orbit in $\mathcal{R}^{\text{Dol}}_{\mathcal{L}}(X,\mathfrak{p},n)$.
	
%\end{thm}


\subsection{The $\mathcal{L}$-Hodge moduli spaces for principal $G$-bundles}
\label{The-L-Hodge-moduli-spaces-for-principal-bundles}

Recall from Subsection~\eqref{subsec: L-Hodge modul space} that to the Lie algebroid $\mathcal{L}$ one can associate a $\mathcal{L}$--twisted de Rham $D$--algebra $\Lambda^{\text{DR}}_{\mathcal{L}}$. Using the Rees construction, we obtain a family of $D$--algebra 
$\Lambda_{\mathcal{L}}^{\mathrm{red}}$ on $X \times \mathbf{A}^{1}$ over
$\mathbf{A}^{1}$ whose fibre over $1$ is $\Lambda^{\text{DR}}_{\mathcal{L}}$ and whose fibre
over $0$ is isomorphic to its associated graded algebra
\[
\Lambda_{\mathcal{L}}^{\mathrm{Dol}}
:= \operatorname{Gr}_{\bullet}(\Lambda_{\mathcal{L}})
\cong \operatorname{Sym}^{\bullet}(V).
\]


Let $\mathcal{L}^{\mathrm{red}}$ denote the Lie algebroid corresponding to
$\Lambda^{\mathrm{red}}_{\mathcal{L}}$ on $X \times \mathbf{A}^{1}$.
We now consider the moduli space
$\mathcal{M}_{\mathcal{L}}^{\mathrm{Hod}}(X,G)$ of principal $G$-bundles equipped
with integrable $\mathcal{L}^{\mathrm{red}}$-connections on
$X \times \mathbf{A}^{1}$.
By Subsection~\eqref{susec: moduli space of L connection on PB}, this moduli space
is a quasi-projective variety endowed with a morphism
\[
\pi:\mathcal{M}_{\mathcal{L}}^{\mathrm{Hod}}(X,G) \longrightarrow \mathbf{A}^{1}.
\]

It follows that for all $\lambda \in \mathbf{A}^{1} \setminus \{0\}$,
\[
\pi^{-1}(\lambda)
= \mathcal{M}_{\mathcal{L}_{\lambda}}(X,G)
\cong \mathcal{M}_{\mathcal{L}_{1}}(X,G)
= \pi^{-1}(1),
\]
while
\[
\pi^{-1}(0)=\mathcal{M}_{\mathcal{L}}^{\mathrm{Dol}}(X,G).
\]
Thus, the space $\mathcal{M}_{\mathcal{L}}^{\mathrm{Hod}}(X,G)$ interpolates
between the moduli space of $\mathcal{L}_{\lambda}$-connections on principal
$G$-bundles and the moduli space of $\mathcal{L}$-twisted principal $G$-Higgs
bundles. The moduli space
$\mathcal{M}_{\mathcal{L}}^{\mathrm{Hod}}(X,G)$ is called the
$\mathcal{L}$-\emph{Hodge moduli space for principal $G$-bundles}.

\begin{thm}
	The $\mathcal{L}$-Hodge moduli space for principal $G$-bundles,
	$\mathcal{M}_{\mathcal{L}}^{\mathrm{Hod}}(X,G)$, is a quasi-projective variety
	equipped with a morphism
	\[
	\pi:\mathcal{M}_{\mathcal{L}}^{\mathrm{Hod}}(X,G)\longrightarrow \mathbf{A}^{1}.
	\]
\end{thm}



%%%%%%%%%%%%%%%%%%%%%%%%%%%%%%%%%%%%%%%%%%%%%%%%%%%%%%%%%%%%%%%%%%%%%%%%%
%%%%%%%%%%%%%%%%%%%%%%%%%%%%%%%%%%%%%%%%%%%%%%%%%%%%%%%%%%%%%%%%%%%%%%%%%








\section{Semiprojectivity of the moduli spaces of $\mathcal{L}$-twisted principal objects}
Moduli spaces of $\mathcal{L}$-twisted Higgs bundles, and $\mathcal{L}$-twisted principal $G$--Higgs bundles of harmonic types are quasi-projective but, in general, non-proper. 
The notion of \emph{semiprojectivity} provides a natural framework to control this non-properness via an algebraic $\mathbf{G}_{m}:=\mathbb{C}^{\ast}$--action with projective fixed point locus and well-defined limits. 
This plays a central role in the study of topology, Hodge theory, and degeneration phenomena for moduli spaces. 
In particular, for smooth semiprojective varieties, the Grothendieck motivic class admits an explicit expression in terms of the projective fixed point locus via the Białynicki--Birula decomposition.

Let $Y$ be a quasi-projective variety endowed with an algebraic $\mathbf{G}_{m}$-action
\[
\mathbf{G}_{m} \times Y \longrightarrow Y, \qquad (t,y) \longmapsto t \cdot y .
\]

\begin{defn}\cite[\S\,2.1,]{Hausel-Hitchin-2022}
	A quasi-projective variety $Y$ with an algebraic $\mathbf{G}_{m}$-action is called \emph{semiprojective} if the following conditions are satisfied:
	\begin{itemize}
		\item For every $y \in Y$, there is a $p\in Y^{\mathbf{G}_{m}}$ such that $
		\lim_{t \to 0} t \cdot y = p$. (i.e, we mean that there exists a $\mathbf{G}_{m}$-equivariant morphism $g : \mathbf{A}^{1} \rightarrow Y$ with $g(1) = y$ and $g(0) = p$, where
		$\mathbf{G}_{m}$ acts on $\mathbf{A}^{1}$ in the standard way.)
		
		\item The fixed point locus $
		Y^{\mathbf{G}_{m}} := \{\, y \in Y \mid t \cdot y = y \text{ for all } t \in \mathbf{G}_{m} \,\}$ is proper (equivalently, projective).
	\end{itemize}
\end{defn}
Examples of semi-projective varieties include cotangent bundles of smooth
projective varieties, moduli spaces of Higgs bundles and Nakajima quiver
varieties.

Let $J$ denote the standard Hermitian metric on $\mathbb{C}^{n}$.  
Define $\mathcal{R}^{\text{DR}^{J}}_{\mathcal{L}}(X,\mathfrak{p},n)$ to be the space whose points are pairs $(E,\beta)$, where
\begin{itemize}
	\item $E$ is a vector bundle on $X$ equipped with an integrable $\mathcal{L}$-connection, and
	\item $\beta \colon E_{\mathfrak{p}} \xrightarrow{\sim} \mathbb{C}^{n}$ is a framing at $\mathfrak{p}$,
\end{itemize}
such that there exists a harmonic metric $K$ on $E$ satisfying
\[
\beta(K_{\mathfrak{p}}) = J .
\]

Similarly, define $\mathcal{R}^{\mathrm{Dol}^{J}}_{\mathcal{L}}(X,\mathfrak{p},n)$ to be the space of pairs $(E,\beta)$, where
\begin{itemize}
	\item $E$ is an $\mathcal{L}$-twisted Higgs bundle on $X$, and
	\item $\beta \colon E_{\mathfrak{p}} \xrightarrow{\sim} \mathbb{C}^{n}$ is a framing,
\end{itemize}
such that there exists a harmonic metric $K$ on $E$ with $\beta(K_{\mathfrak{p}})=J$.

Now suppose $G$ is a reductive affine algebraic group.  
Fix a maximal compact subgroup $V \subset G$ and choose an embedding $
G \hookrightarrow GL_{n}(\mathbb{C})$ such that $V = G \cap U(n)$.

Define
\[
\mathcal{R}^{\text{DR}^{J}}_{\mathcal{L}}(X,\mathfrak{p},G)
:= \mathcal{R}^{\text{DR}}_{\mathcal{L}}(X,\mathfrak{p},G)
\cap
\mathcal{R}^{\text{DR}^{J}}_{\mathcal{L}}(X,\mathfrak{p},n),
\]
and endow it with the subspace topology induced from the analytic topology on
$\mathcal{R}^{\text{DR}^{\mathrm{an}}}_{\mathcal{L}}(X,\mathfrak{p},n)$.

Similarly, define
\[
\mathcal{R}^{\mathrm{Dol}^{J}}_{\mathcal{L}}(X,\mathfrak{p},G)
:= \mathcal{R}^{\mathrm{Dol}}_{\mathcal{L}}(X,\mathfrak{p},G)
\cap
\mathcal{R}^{\mathrm{Dol}^{J}}_{\mathcal{L}}(X,\mathfrak{p},n),
\]
with the induced topology from
$\mathcal{R}^{\mathrm{Dol}^{\mathrm{an}}}_{\mathcal{L}}(X,\mathfrak{p},n)$.

The arguments of \cite[\S\,9, Lemmas~9.13 and~9.14]{CT.Simpson-1995} extend verbatim to the Lie algebroid setting, by replacing vector bundles with integrable connections (respectively Higgs bundles) by their $\mathcal{L}$-twisted analogues throughout. Consequently, we obtain the following result, which generalizes \cite[\S\,9, Lemma~9.15]{CT.Simpson-1995}.

\begin{thm}\label{proper map between moduli of principal objects}
	Let $G$ and $H$ be reductive affine algebraic groups and let $
	G \hookrightarrow H$ 
	be an injective homomorphism. Then the induced morphisms 
	\[
	\mathcal{M}^{\text{DR}}_{\mathcal{L}}(X,G) \longrightarrow \mathcal{M}^{\text{DR}}_{\mathcal{L}}(X,H),
	\qquad
	\mathcal{M}^{\mathrm{Dol}}_{\mathcal{L}}(X,G)
	\longrightarrow
	\mathcal{M}^{\mathrm{Dol}}_{\mathcal{L}}(X,H)
	\]
	are proper.In particular, the natural morphisms
	\[
	\mathcal{M}^{\text{DR}}_{\mathcal{L}}(X,G) \hookrightarrow \mathcal{M}^{\text{DR}}_{\mathcal{L}}(X,n),
	\qquad
	\mathcal{M}^{\mathrm{Dol}}_{\mathcal{L}}(X,G)
	\hookrightarrow
	\mathcal{M}^{\mathrm{Dol}}_{\mathcal{L}}(X,n)
	\]
	are proper. The same conclusion holds for the $\mathcal{L}$--Hodge moduli spaces
	\[
	\mathcal{M}^{\mathrm{Hod}}_{\mathcal{L}}(X,G)
	\longrightarrow
	\mathcal{M}^{\mathrm{Hod}}_{\mathcal{L}}(X,n).
	\]
\end{thm}

\begin{proof}
	Since $G \hookrightarrow H$ is injective, the induced map $
	\mathcal{R}^{\text{DR}^{J}}_{\mathcal{L}}(X,\mathfrak{p},G)
	\longrightarrow
	\mathcal{R}^{\text{DR}^{J}}_{\mathcal{L}}(X,\mathfrak{p},H)
	$
	identifies $\mathcal{R}^{\text{DR}^{J}}_{\mathcal{L}}(X,\mathfrak{p},G)$ with a closed subset of
	$\mathcal{R}^{\text{DR}^{J}}_{\mathcal{L}}(X,\mathfrak{p},H)$. Hence the composite map
	\[
	\mathcal{R}^{\text{DR}^{J}}_{\mathcal{L}}(X,\mathfrak{p},G)
	\longrightarrow
	\mathcal{M}^{\text{DR}}_{\mathcal{L}}(X,H)
	\]
	is proper.
	
	This map factors through the quotient
	$\mathcal{M}_{\mathcal{L}}(X,G)$, and the projection
	\[
	\mathcal{R}^{\text{DR}^{J}}_{\mathcal{L}}(X,\mathfrak{p},G)
	\longrightarrow
	\mathcal{M}^{\text{DR}}_{\mathcal{L}}(X,G)
	\]
	is surjective. Therefore the induced morphism
	\[
	\mathcal{M}^{\text{DR}}_{\mathcal{L}}(X,G)
	\longrightarrow
	\mathcal{M}^{\text{DR}}_{\mathcal{L}}(X,H)
	\]
	is proper. The same argument applies to the $\mathcal{L}$--twisted Dolbeault moduli spaces.
	
	For the $\mathcal{L}$--Hodge moduli spaces, properness follows from the
	topological trivialization
	\[
	\mathcal{M}^{\mathrm{Hod}}_{\mathcal{L}}(X,G)
	\cong
	\mathcal{M}^{\text{DR}}_{\mathcal{L}}(X,G) \times \mathbf{A}^{1},
	\]
	which is functorial in $G$.
\end{proof}



One defines the Hitchin map for the $\mathcal{L}$--twisted Dolbeault moduli space
(cf. \cite[\S\,2, p.~9]{Alfaya-Oliveire-2024})
\[
\mathcal{H}:
\mathcal{M}^{\mathrm{Dol}}_{\mathcal{L}}(X,n)
\longrightarrow
\mathcal{W}, \ \text{where} \ \mathcal{W} \ \text{is the Htchin base.}
\]

\begin{lem}\cite[\S\,2, Lem.~2.4]{Alfaya-Oliveire-2024}.\label{lem:proper Hitchin map for ML(X,n)}
	The Hitchin map $
	\mathcal{H}:
	\mathcal{M}^{\mathrm{Dol}}_{\mathcal{L}}(X,n)\longrightarrow \mathcal{W}$ is proper.
\end{lem}




\subsection{Semiprojectivity of moduli space $\mathcal{M}_{\mathcal{L}}^{\text{Dol}}(X,G)$.}
There is a natural $\mathbf{G}_{m}$--action on the category of
$\mathcal{L}$--twisted Higgs bundles defined by
\[
t\cdot (E,\theta_{\mathcal{L}}):=(E,t\cdot\theta_{\mathcal{L}}).
\]
Since $G$ is reductive, this induces a $\mathbf{G}_{m}$--action on
$\mathcal{M}^{\mathrm{Dol}}_{\mathcal{L}}(X,n)$.

\begin{proposition}\cite[\S\,4, Prop. 4.4]{Alfaya-Oliveire-2024}\label{pro:semiprojectivty of moduli of L twisted Higgs bundle}
	The moduli space $\mathcal{M}^{\mathrm{Dol}}_{\mathcal{L}}(X,n)$
	is a semiprojective variety.
\end{proposition}

Similarly, there is a natural $\mathbf{G}_{m}$--action on the category of
$\mathcal{L}$--twisted principal $G$--Higgs bundles defined by
\[
t\cdot (E_{G},\theta_{\mathcal{L}}):=(E_{G},t\theta_{\mathcal{L}}).
\]
If $G$ is reductive, this induces a $\mathbf{G}_{m}$--action on
$\mathcal{M}^{\mathrm{Dol}}_{\mathcal{L}}(X,G)$.
This action is compatible with the functoriality induced by morphisms of
reductive groups and coincides with the standard action when
$G=GL_{n}(\mathbb{C})$.

\begin{lem}\label{limit exist for MLDOL of PB}
	For any $y\in \mathcal{M}^{\mathrm{Dol}}_{\mathcal{L}}(X,G)$,
	the limit $
	\lim_{t\to 0} t\cdot y$ 
	exists and is a fixed point of the $\mathbf{G}_{m}$--action.
\end{lem}

\begin{proof}
	Proposition~\eqref{pro:semiprojectivty of moduli of L twisted Higgs bundle}
	proves the statement for $G=GL_{n}(\mathbb{C})$.
	Choose a faithful representation
	$G\hookrightarrow GL_{n}(\mathbb{C})$, which induces a closed immersion
	\[
	i:\mathcal{M}^{\mathrm{Dol}}_{\mathcal{L}}(X,G)
	\hookrightarrow
	\mathcal{M}^{\mathrm{Dol}}_{\mathcal{L}}(X,n).
	\]
	The orbit map
	$\mathbf{G}_{m}\to \mathcal{M}^{\mathrm{Dol}}_{\mathcal{L}}(X,n)$
	extends to a morphism
	$\mathbf{A}^{1}\to \mathcal{M}^{\mathrm{Dol}}_{\mathcal{L}}(X,n)$.
	By the properness of the maps in
	Theorem~\ref{proper map between moduli of principal objects},
	the induced orbit
	$\mathbf{G}_{m}\to \mathcal{M}^{\mathrm{Dol}}_{\mathcal{L}}(X,G)$
	extends to a morphism
	$\mathbf{A}^{1}\to \mathcal{M}^{\mathrm{Dol}}_{\mathcal{L}}(X,G)$.
	The image of $0\in\mathbf{A}^{1}$ is the required fixed point.
\end{proof}

\begin{lem}\label{fixed point for MLDOL of PB}
	The fixed point locus of the $\mathbf{G}_{m}$--action on
	$\mathcal{M}^{\mathrm{Dol}}_{\mathcal{L}}(X,G)$
	is a proper scheme contained in
	$(\mathcal{H}\circ i)^{-1}(0)$,
	where
	\[
	i:\mathcal{M}^{\mathrm{Dol}}_{\mathcal{L}}(X,G)
	\hookrightarrow
	\mathcal{M}^{\mathrm{Dol}}_{\mathcal{L}}(X,n)
	\]
	is the closed immersion.
\end{lem}

\begin{proof}
	The Hitchin map $\mathcal{H}$ is proper and $\mathbf{G}_{m}$--equivariant;
	hence the composition $\mathcal{H}\circ i$ is also proper and
	$\mathbf{G}_{m}$--equivariant.
	Therefore the fixed point locus satisfies
	\[
	\mathcal{M}^{\mathrm{Dol}}_{\mathcal{L}}(X,G)^{\mathbf{G}_{m}}
	\subseteq
	(\mathcal{H}\circ i)^{-1}(\mathcal{W}^{\mathbf{G}_{m}})
	=
	(\mathcal{H}\circ i)^{-1}(0).
	\]
	Since the fiber over $0$ of a proper morphism is proper, the fixed point
	locus is proper.
\end{proof}

\begin{proposition}\label{semiprojectivity of MLDOL(X,G)}
	The moduli space $\mathcal{M}^{\mathrm{Dol}}_{\mathcal{L}}(X,G)$
	is a semiprojective variety.
\end{proposition}

\begin{proof}
	The variety $\mathcal{M}^{\mathrm{Dol}}_{\mathcal{L}}(X,G)$ is
	quasi-projective.
	By Lemma~\eqref{limit exist for MLDOL of PB}, limits of $\mathbf{G}_{m}$--orbits
	exist, and by Lemma~\eqref{fixed point for MLDOL of PB} the fixed point locus
	is proper.
	Hence $\mathcal{M}^{\mathrm{Dol}}_{\mathcal{L}}(X,G)$ is semiprojective.
\end{proof}


\subsection{Semiprojectivity of moduli space $\mathcal{M}_{\mathcal{L}}^{\text{Hod}}(X,G)$.}
The semiprojectivity of the Hodge moduli space
$\mathcal{M}^{\mathrm{Hod}}_{\mathcal{L}}(X,n)$
is well known.

\begin{thm}\label{thm: semiprojectivity of MLHOD(X,n)}
	\cite[\S\,4, Th.~4.12]{Alfaya-Oliveire-2024}
	The $\mathcal{L}$--Hodge moduli space for vector bundles
	$\mathcal{M}^{\mathrm{Hod}}_{\mathcal{L}}(X,n)$
	is a semiprojective variety.
\end{thm}

There is a natural $\mathbf{G}_{m}$--action on the category of principal
$G$--bundles with integrable $\mathcal{L}^{\mathrm{red}}$--connections,
defined by
\[
t\cdot (E_{G},\nabla_{\mathcal{L}^{\mathrm{red}}})
:=(E_{G},t\nabla_{\mathcal{L}^{\mathrm{red}}}).
\]
If $G$ is reductive, this induces a $\mathbf{G}_{m}$--action on
$\mathcal{M}^{\mathrm{Hod}}_{\mathcal{L}}(X,G)$.
This action is compatible with functoriality with respect to morphisms of
reductive groups and coincides with the standard action when
$G=GL_{n}(\mathbb{C})$.

\begin{lem}\label{limit exist for MLHOD of PB}
	For any $y\in \mathcal{M}^{\mathrm{Hod}}_{\mathcal{L}}(X,G)$,
	the limit $
	\lim_{t\to 0} t\cdot y$
	exists and is a fixed point of the $\mathbf{G}_{m}$--action on
	$\mathcal{M}^{\mathrm{Hod}}_{\mathcal{L}}(X,G)$.
\end{lem}

\begin{proof}
	Theorem~\eqref{thm: semiprojectivity of MLHOD(X,n)} proves the statement
	for $G=GL_{n}(\mathbb{C})$.
	Choose a faithful representation
	$G\hookrightarrow GL_{n}(\mathbb{C})$, inducing a closed immersion
	\[
	i:\mathcal{M}^{\mathrm{Hod}}_{\mathcal{L}}(X,G)
	\hookrightarrow
	\mathcal{M}^{\mathrm{Hod}}_{\mathcal{L}}(X,n).
	\]
	The orbit map
	$\mathbf{G}_{m}\to \mathcal{M}^{\mathrm{Hod}}_{\mathcal{L}}(X,n)$
	extends to a morphism
	$\mathbf{A}^{1}\to \mathcal{M}^{\mathrm{Hod}}_{\mathcal{L}}(X,n)$.
	By the properness of the morphisms in
	Theorem~\ref{proper map between moduli of principal objects},
	the induced orbit
	$\mathbf{G}_{m}\to \mathcal{M}^{\mathrm{Hod}}_{\mathcal{L}}(X,G)$
	extends to a morphism
	$\mathbf{A}^{1}\to \mathcal{M}^{\mathrm{Hod}}_{\mathcal{L}}(X,G)$.
	The image of $0\in \mathbf{A}^{1}$ is the desired fixed point.
\end{proof}

 
\begin{lem}\label{fixed point for MLHOD of PB}
	The fixed point locus of the $\mathbf{G}_{m}$--action on
	$\mathcal{M}^{\mathrm{Hod}}_{\mathcal{L}}(X,G)$
	is proper.
\end{lem}

\begin{proof}
	The argument follows \cite[\S\,10, Lemma~10.4]{CT.Simpson-1997}.
	The fixed point locus of the $\mathbf{G}_{m}$--action lies over the origin
	$0\in \mathbf{A}^{1}$ under the Hodge projection
	\[
	\pi:\mathcal{M}^{\mathrm{Hod}}_{\mathcal{L}}(X,G)\to \mathbf{A}^{1}.
	\]
	Therefore it coincides with the fixed point locus of the
	$\mathbf{G}_{m}$--action on the Dolbeault moduli space
	$\mathcal{M}^{\mathrm{Dol}}_{\mathcal{L}}(X,G)$.
	By Lemma~\eqref{fixed point for MLDOL of PB}, this fixed point locus is
	proper.
\end{proof}

 
\begin{proposition}\label{semiprojectivity of MLHOD(X,G)}
	The moduli space
	$\mathcal{M}^{\mathrm{Hod}}_{\mathcal{L}}(X,G)$
	is a semiprojective variety.
\end{proposition}

\begin{proof}
	The variety $\mathcal{M}^{\mathrm{Hod}}_{\mathcal{L}}(X,G)$ is
	quasi-projective.
	By Lemma~\eqref{limit exist for MLHOD of PB}, limits of $\mathbf{G}_{m}$--orbits
	exist, and by Lemma~\eqref{fixed point for MLHOD of PB}, the fixed point locus
	is proper.
	Hence $\mathcal{M}^{\mathrm{Hod}}_{\mathcal{L}}(X,G)$ is semiprojective.
\end{proof}

%%%%%%%%%%%%%%%%%%%%%%%%%%%%%%%%%%%%%%%%%%%%%%%%%%%%%%%%%%%%%%%%%%%%%%%%%%%%%%%%%%%%%%%%%%%%%%%%%%%%%%%%%%%%%%%%%%%%%%%%%%%%%%%%%%%%%%%%%%%%%%%%%%%%

\section{Motives of moduli of $\mathcal{L}$--twisted principal $G$--objects}

Another goal of this paper is to compare the class of the moduli space $\mathcal{M}^{\text{DR}}_{\mathcal{L}}(X,G)$ and $\mathcal{M}^{\text{Dol}}_{\mathcal{L}}(X,G)$ in the (completion of the)
Grothendieck ring of varieties. We recall such ring and its basic properties (cf. \cite[\S\,5]{Alfaya-Oliveire-2024}).
Denote by $Var_{\mathbb{C}}$ the category of quasi-projective varieties over $\mathbb{C}$. For each $Y \in Var_{\mathbb{C}}$, let $[Y]$ denote the corresponding isomorphism class. Consider the group obtained by the free abelian group on isomorphism classes $[Y]$, modulo the relation
\[[Y] = [Y'] + [Y \backslash Y'],\]
where $Y' \subset Y$ is a Zariski-closed subset. In particular, in such group,
\[[Y] + [Z] = [Y \sqcup Z],\]
where $\sqcup$ denotes disjoint union. If we define the product
\[[Y] \cdot [Z] = [Y \times Z],\]
in this quotient, then we obtain a commutative ring, known as the \textit{Grothendieck ring of varieties} and denoted by $\mathcal{K}(Var_{\mathbb{C}})$.
Then $0 = [\emptyset]$ and $1 = [Spec(\mathbb{C})]$ are the additive and multiplicative units of this ring. 

The class of the affine line, sometimes called the \textit{Lefschetz object}, is denoted by $\mathbb{L} := [\mathbf{A}^{1}] = [\mathbb{C}]$.
Of course, $\mathbb{L}^{n} = [\mathbf{A}^{n}] = [\mathbb{C}^{n}]$. We will consider the localization $\mathcal{K}(Var_{\mathbb{C}})[\mathbb{L}^{-1}]$, and then the dimensional completion,

\[\hat{\mathcal{K}}(VarC) := \left\{\sum_{r\geqslant 0}[Y_{r}]\mathbb{L}^{-r}: [Y_{r}]\in \mathcal{K}(Var_{\mathbb{C}}) \ \text{with} \ \dim Y_{r}-r\longrightarrow -\infty\right\}
\]

Notice that we have a map $\mathcal{K}(Var_{\mathbb{C}}) \rightarrow \hat{\mathcal{K}}(Var_{\mathbb{C}})$. Observe also that $\mathbb{L}^{n}-1$ is invertible in $\hat{\mathcal{K}}(Var_{\mathbb{C}})$, for every $n$, with
inverse equal to $\mathbb{L}^{-n}\sum_{k=0}^{\infty}\mathbb{L}^{-nk}$. This is the reason why we had to introduce the completion $\hat{\mathcal{K}}(Var\mathbb{C})$: there will be
computations in which we will need to invert elements of the form $\mathbb{L}^{n}$ or $\mathbb{L}^{n} - 1$.
In this paper, by motive we mean the following.

\begin{defn}
	Let $Y$ be a quasi-projective variety. The class $[Y]$ in $\mathcal{K}(Var_{\mathbb{C}})$ or in $\hat{\mathcal{K}}(Var_{\mathbb{C}})$ is called the \textit{motive, or motivic class}, of $Y$.
\end{defn}

Recall that, if $\pi: Y \rightarrow B$ is an algebraic fibre bundle (thus Zariski locally trivial), with fibre $F$, then $[Y] = [F] \cdot [B]$ .

\medskip

\noindent
\textbf{Mixed Hodge Structures and \(E\)--polynomials}.
For the convenience of the reader, we briefly recall the notion of mixed Hodge structures and \(E\)--polynomials following \cite[\S\,2]{Hausel-Rodriguez-2008}. Let \(Y\) be a complex algebraic variety. Deligne showed that the cohomology group \(H^{k}(Y,\mathbb{Q})\) carries a canonical mixed Hodge structure consisting of an increasing weight filtration
\[
W_{0}\subseteq W_{1}\subseteq \cdots \subseteq W_{2k}=H^{k}(Y,\mathbb{Q}),
\]
and a decreasing Hodge filtration
\[
H^{k}(Y,\mathbb{C})=F^{0}\supseteq F^{1}\supseteq \cdots \supseteq F^{k}\supseteq 0,
\]such that the filtration induced by $F$ on the complexification of the graded pieces $\text{Gr}^{W}_{\ell}:= W_{\ell}/W_{\ell - 1}$ of the
weight filtration endows every graded piece with a pure Hodge structure of weight $\ell$.

The \textit{mixed Hodge numbers} are defined by
\[
h^{p,q;k}_{c}(Y)
=
\dim_{\mathbb{C}}
\text{Gr}^{p}_{F}\text{Gr}^{W}_{p+q}H^{k}_{c}(Y,\mathbb{C}),
\]
where \(H^{k}_{c}(Y,\mathbb{C})\) denotes the compactly supported cohomology.

The corresponding \textit{mixed Hodge polynomial} is
\[
H_{c}(Y;x,y,t)
=
\sum_{p,q,k}
h^{p,q;k}_{c}(Y)x^{p}y^{q}t^{k},
\]
and the \textit{\(E\)--polynomial (or Hodge--Deligne polynomial)} is obtained by specialization:
\[
E(Y;x,y)
=
H_{c}(Y;x,y,-1).
\]


A fundamental property of the \(E\)--polynomial is its additivity: if \(Z\subseteq Y\) is a closed subvariety, then
\[
E(Y;x,y)=E(Z;x,y)+E(Y\setminus Z;x,y).
\]
Moreover, if \(X\to Y\) is a Zariski locally trivial fibration with fibre \(F\), then
\[
E(X;x,y)=E(Y;x,y)\,E(F;x,y).
\]

These properties make \(E\)--polynomials particularly useful in the study of moduli spaces

\begin{rem}
Actually, the $E$ polynomial can be seen as a ring map,
\[E:\hat{\mathcal{K}}(\text{Var}_{\mathbb{C}})\rightarrow \mathbb{Z}[x,y] \left \llbracket \frac{1}{xy} \right\rrbracket\]with values in the Laurent series in $xy$, which takes values in $\mathbb{Z}[x,y]$ when restricted to $\mathcal{K}(\text{Var}_{\mathbb{C}})$.
\end{rem}

\begin{example}
	\begin{enumerate}
		\item For the affine line $\mathbb{A}^{1}$, $
		E(\mathbb{A}^{1};x,y)=xy.$
		 More generally, $$
		E(\mathbb{A}^{n};x,y)=(xy)^{n}.$$
		\item For the multiplicative group $\mathbb{C}^{*}$, $
		E(\mathbb{C}^{*};x,y)=xy-1$.
		Indeed, \(\mathbb{C}^{*}=\mathbb{A}^{1}\setminus \{0\}\), and the \(E\)--polynomial is additive with respect to stratifications.
		
		\item Let \(X\) be a smooth projective curve of genus \(g\). Since \(X\) is smooth and projective, its mixed Hodge structure is pure. The nonzero Hodge numbers are
		\[
		h^{0,0}=1,\qquad h^{1,0}=h^{0,1}=g,\qquad h^{1,1}=1.
		\]
		Hence $
		E(X;x,y)
		=
		\sum_{p,q}(-1)^{p+q}h^{p,q}x^{p}y^{q}$,
		which gives
		\[
		E(X;x,y)
		=
		1-gx-gy+xy.
		\]
		For example, if \(X=\mathbb{P}^{1}\), then \(g=0\), and therefore $
		E(\mathbb{P}^{1};x,y)=1+xy$.
		
		If $X$ is an elliptic curve, then \(g=1\), so
		\[
		E(X;x,y)=1-x-y+xy=(1-x)(1-y).
		\]
\end{enumerate}
\end{example}

Let $Y$ be a semiprojective variety and $\alpha \in Y^{\mathbf{G}_{m}}$. The \textit{upward flow from $\alpha$} is,
\[ F_{\alpha}^{+}:= \{x \in Y : \lim_{t\rightarrow 0}t\cdot x = \alpha\} \subseteq Y.\]
The upward flows define the \textit{Białynicki-Birula decomposition},
\[Y = \bigsqcup_{\alpha \in Y^{\mathbf{G}_{m}}}F_{\alpha}^{+}.\]
Similarly, \textit{the downward flow from $\alpha$} is
\[ F_{\alpha}^{-}:= \{x \in Y : \lim_{t\rightarrow \infty}t\cdot x = \alpha\} \subseteq Y.\]
and the \textit{core of $Y$} is defined to be, 
\[\mathcal{C}:=\bigsqcup_{\alpha \in X^{\mathbf{G}_{m}}}F_{\alpha}^{-}\subseteq Y\]
Given a smooth fixed point $\alpha \in Y^{\mathbf{G}_{m}}$, its upward and downward flows can be described by studying the tangent space $T_{\alpha}Y$. Since $\alpha$ is fixed, $T_{\alpha}Y$ has an induced $\mathbf{G}_{m}$--action which provides a weight space decomposition $T_{\alpha}Y =\bigoplus_{k\in \mathbb{Z}}(T_{\alpha}Y)_{k}$. Let $T^{+}_{\alpha}Y := \bigoplus_{k(>0)\in \mathbb{Z}}(T_{\alpha}Y)_{k}$, be the subspace of positive weights and $T^{-}_{\alpha}X := \bigoplus_{k(<0)\in \mathbb{Z}}(T_{\alpha}Y)_{k}$, the subspace of negative weights \cite[\S\,2, Proposition 2.1]{Hausel-Hitchin-2022}. Let $Y$ be a semiprojective variety and $\alpha \in  Y^{G_{m}}$ a smooth fixed point. The upward flow $F_{\alpha}^{+}$ (resp. the downward flow $F_{\alpha}^{-}$) is a locally closed $\mathbf{G}_{m}$-invariant subvariety of $Y$ isomorphic to $T^{+}_{\alpha}Y (\text{resp.} T^{-}_{\alpha}X)$ as $\mathbf{G}_{m}$--varieties. Define $N^{+}_{\alpha}:=\dim T^{+}_{\alpha}Y, \ N^{-}_{\alpha}:=\dim T^{-}_{\alpha}Y$.

We now recall the definition \cite[Definition 2.12]{Hausel-Hitchin-2022} of a very stable point in a semiprojective variety.

\begin{defn}
	A smooth fixed point $\alpha \in Y^{\mathbf{G}_{m}}$ is \textit{very stable} if $F^{+}_{\alpha} \bigcap \mathcal{C}= \{\alpha\}$.
\end{defn}


\begin{proposition}\cite[\S\,5, Theorem 5.6]{Alfaya-Oliveire-2024}\label{prop: motivic equality of smooth complex variety}
Let $Y$ be a smooth semiprojective complex variety endowed with a $\mathbb{G}_{m}$--equivariant surjective submersion $\pi:Y\rightarrow \mathbb{A}^{1}$ covering the standard scalling action on $\mathbb{C}$. Then the following motivic equalites holds in $\hat{\mathcal{K}}(\text{Var}_{\mathbb{C}})$:
\[[\pi^{-1}(0)]=[\pi^{-1}(1)] \ \ \text{and}  \ \ [X]=\mathbb{L}[\pi^{-1}(0)]\]where $\mathbb{L}$ is the Lefschetz object.
\end{proposition}

Let \(Y\) be a semiprojective variety over \(\mathbb{C}\). We denote by \(\widehat{Y}\) the smooth locus of \(Y\). Note that smooth locus of semiprojective variety is again a semiprojective variety.
\begin{thm}\label{thm:motivic equality of dr-dol-Hod principal moduli space}
	Let $\mathcal{L} = (V,[·,·],\delta)$ be a transitive Lie algebroid on $X$. Then the following equalities hold $\hat{\mathcal{K}}(Var_{\mathbb{C}})$,
	
	\begin{enumerate}
		\item $[\widehat{\mathcal{M}^{\text{Dol}}_{\mathcal{L}}(X,G)}]=\sum_{\alpha\in X^{\mathbf{G}_{m}}}\mathbb{L}^{N_{\alpha}^{+}}[F_{\alpha}]$
		
		\item $[\widehat{\mathcal{M}^{\text{DR}}_{\mathcal{L}}(X,G)}] = [\widehat{\mathcal{M}_{\mathcal{L}}^{\text{Dol}}(X,G)}, \ \  [\widehat{\mathcal{M}_{\mathcal{L}}^{\text{Hod}}(X,G)]} = \mathbb{L}[\widehat{\mathcal{M}_{\mathcal{L}}^{\text{Dol}}(X,G)}]$,
		
		\item $E(\widehat{\mathcal{M}_{\mathcal{L}}^{\text{DR}}(X,G))} = E(\widehat{\mathcal{M}_{\mathcal{L}}^{\text{Dol}}(X,G))}$, \ \  $E(\widehat{\mathcal{M}_{\mathcal{L}}^{\text{Hod}}(X,G))} = xyE(\widehat{\mathcal{M}_{\mathcal{L}}^{\text{Dol}}(X,G))}$,
		
		\item we have an isomorphism of Hodge structures,
		\[H^{\bullet}(\widehat{\mathcal{M}_{\mathcal{L}}^{\text{DR}}(X,G)}) \cong H^{\bullet}(\widehat{\mathcal{M}_{\mathcal{L}}^{\text{Dol}}(X,G)})\]and both $\widehat{\mathcal{M}_{\mathcal{L}}^{\text{DR}}(X,G)}$ and $\widehat{\mathcal{M}_{\mathcal{L}}^{\text{Hod}}(X,G)}$ have pure mixed Hodge structure.
		
		
		\item A smooth fixed point 
		\[
		\alpha \in \mathcal{M}^{\mathrm{DR}}_{\mathcal{L}}(X,G)
		\quad 
		\left(\text{respectively} \ 
		\mathcal{M}^{\mathrm{Dol}}_{\mathcal{L}}(X,G),\ 
		\mathcal{M}^{\mathrm{Hod}}_{\mathcal{L}}(X,G)
		\right)
		\]
		is very stable if and only if the corresponding upward flow \(F^{+}_{\alpha}\) is closed.
	\end{enumerate}
	
\end{thm}
\begin{proof}
	 Since, $\widehat{\mathcal{M}_{\mathcal{L}}^{\text{Dol}}(X,G)}$ is smooth semiprojective variety, then  $(1)$ is follows from \cite[lemma 5.5]{Alfaya-Oliveire-2024}. The moduli space $Z=\widehat{\mathcal{M}_{\mathcal{L}}^{\text{Hod}}(X,G)}$ is a smooth semiprojective variety for the $\mathbf{G}_{m}$-action . Moreover, the map $\pi|_{Z}$ from \eqref{The-L-Hodge-moduli-spaces-for-principal-bundles} is a surjective $\mathbf{G}_{m}$-equivariant submersion covering the standard $\mathbf{G}_{m}$--action on $\mathbb{C}$. Then proposition \eqref{prop: motivic equality of smooth complex variety} gives the desired motivic equalities,
	\begin{gather*}
		[\widehat{\mathcal{M}_{\mathcal{L}}^{\text{Dol}}(X,G)}]=[\pi|_{Z}^{-1}(0)] = [\pi|_{Z}^{-1}(1)]=[\widehat{\mathcal{M}_{\mathcal{L}}^{\text{DR}}(X,G)}], \ \text{and} \\ [\widehat{\mathcal{M}_{\mathcal{L}}^{\text{Hod}}(X,G)}]=\mathbb{L}[\pi|_{Z}^{-1}(0)]=\mathbb{L}[\widehat{\mathcal{M}_{\mathcal{L}}^{\text{Dol}}(X,G)}]
	\end{gather*}
	which yield the corresponding equalities of $E$--polynomials.
	\[E(\widehat{\mathcal{M}_{\mathcal{L}}^{\text{DR}}(X,G)};x,y) = E(\widehat{\mathcal{M}_{\mathcal{L}}^{\text{Dol}}(X,G)};x,y), \ \  E(\widehat{\mathcal{M}_{\mathcal{L}}^{\text{Hod}}(X,G)};x,y) = xyE(\widehat{\mathcal{M}_{\mathcal{L}}^{\text{Dol}}(X,G)};x,y).\]
	
	Moreover, by \cite[Corollary 1.3.3]{T.Hausel-F. Villegas}, the fibres, $$\widehat{\mathcal{M}_{\mathcal{L}}^{\text{Dol}}(X,G)}=\pi|_{Z}^{-1}(0) \ \text{and} \ \widehat{\mathcal{M}_{\mathcal{L}}^{\text{DR}}(X,G)} = \pi|_{Z}^{-1}(1)$$ have isomorphic cohomology supporting pure mixed Hodge structures. As $\widehat{\mathcal{M}_{\mathcal{L}}^{\text{Hod}}(\mathcal{L},G)}$ is also smooth and semiprojective, its cohomology is also pure by \cite[Corollary 1.3.2]{T.Hausel-F. Villegas}. Finally, by \cite[Proposition 2.14]{Hausel-Hitchin-2022}, a smooth fixed point 
	\[
	\alpha \in \mathcal{M}^{\mathrm{DR}}_{\mathcal{L}}(X,G)
	\quad 
	\left(\text{respectively} \ 
	\mathcal{M}^{\mathrm{Dol}}_{\mathcal{L}}(X,G),\ 
	\mathcal{M}^{\mathrm{Hod}}_{\mathcal{L}}(X,G)
	\right)
	\]
	is very stable if and only if the corresponding upward flow \(F^{+}_{\alpha}\) is closed.
\end{proof}
\begin{rem}
If one can construct moduli spaces of \(\mathcal{L}\)--twisted principal objects for an arbitrary Lie algebroid \(\mathcal{L}\), then the above theorem remains valid for any Lie algebroid.
\end{rem}


 
 

 
 
 
 
 
 %%%%%%%%%%%%%%%%%%%%%%%%%%%%%%%%%%%%%%%%%%%%%%%%%%%%%%%%%%%%%%%%%%%%%%%%%%%%%%%%%%%%%%%%%%%%%%%%%%%%%%%%%%%%%%%%%%%%%%%%%%%%%%%%%%%%%%%%%%%%%%%%%%
\section*{Acknowledgment}
The first named author is supported by the \textit{National Board of Higher Mathematics (NBHM)} 
through the Doctoral Research Fellowship Program. 
The second named author is partially supported by the DST INSPIRE Faculty Fellowship (Research Grant No.: DST/INSPIRE/04/2020/000649, IFA-20-MA-144), the Ministry of Science \& Technology, 
Government of India. 


%\bibliographystyle{halpha.bst}
%\bibliography{references}

\begin{thebibliography}{BDPS17}
	\expandafter\ifx\csname url\endcsname\relax
	\def\url#1{\texttt{#1}}\fi
	\expandafter\ifx\csname doi\endcsname\relax
	\def\doi#1{\burlalt{doi:#1}{http://doi.org/#1}}\fi
	\expandafter\ifx\csname urlprefix\endcsname\relax\def\urlprefix{URL }\fi
	\expandafter\ifx\csname href\endcsname\relax
	\def\href#1#2{#2}\fi
	\expandafter\ifx\csname burlalt\endcsname\relax
	\def\burlalt#1#2{\href{#2}{#1}}\fi
	
	\bibitem[Bal06]{Balaji-2006}
	V. Balaji.
	\newblock Lectures on principal bundles.
	\newblock  \href{https://cel.hal.science/cel-00392133/document}{3rd cycle. Guanajuato (Mexique), 2006, pp.17.
	cel-00392133}.
	
	\bibitem[Sor99]{C.Sorger-1999}
	Christoph Sorger.
	\newblock Lectures on moduli of principal
	$G$-bundles over algebraic curves.
	\newblock \href{https://www.math.sciences.univ-nantes.fr/~sorger/assets/pdf/trieste.pdf}{Lecture given at the School on Algebraic Geometry
	Trieste, 26 July - 13 August 1999}.
	
	\bibitem[GP25]{Ghosh-Paul-2025}
	 Samit Ghosh and Arjun Paul.
	\newblock Lie algebroid connections on principal bundles.
	\newblock 	\href{https://doi.org/10.48550/arXiv.2505.18821}{arXiv:2505.18821}.
	
	
	\bibitem[Sim92]{CT. Simpson-1992}
	 Carlos T. Simpson.
	\newblock Higgs bundles and local systems.
	\newblock \href{https://www.scopus.com/pages/publications/51649132477}{Inst. Hautes $\acute{E}$tudes Sci. Publ. Math. 75 (1992), 5–95.}
	
	
	\bibitem[Sim94]{CT.Simpson-1994}
	\newblock Carlos T. Simpson.
	\newblock Moduli of representations of the fundamental group of a smooth projective variety I.
	\newblock \href{https://www.scopus.com/pages/publications/51249163404}{Publi.Math. I.H.$\acute{E}$.S., 79:47–129, 1994.}
	
	\bibitem[Sim95]{CT.Simpson-1995}
	\newblock Carlos T. Simpson.
	\newblock Moduli of representations of the fundamental group of a smooth projective variety II.
	\newblock \href{https://www.numdam.org/item/PMIHES_1994__80__5_0/}{Publi.Math. I.H.$\acute{E}$.S., 80:5–79, 1995.}
	
	\bibitem[Sim97]{CT.Simpson-1997}
	\newblock Carlos T. Simpson.
	\newblock The Hodge filtration on nonabelian cohomology.
	\newblock \href{https://arxiv.org/pdf/alg-geom/9604005}{Proceedings of Symposia in Pure Mathematics, 62.2:217–284, 1997.}
	
	\bibitem[Tor11]{P.Tortella-2011}
	\newblock Pietro Tortella.
	\newblock $\Lambda$-modules and holomorphic Lie algebroid connections.
	\newblock  \href{https://link.springer.com/article/10.2478/s11533-012-0065-z}{Cent. Eur. J. Math. 10 (4) (2012) 1422–1441.}
	
	\bibitem[To11T]{P.Tortella-2011T}
	\newblock Pietro Tortella.
	\newblock $\Lambda$-modules and holomorphic Lie algebroids.
	\newblock  \href{https://iris.sissa.it/retrieve/dd8a4bf7-0696-20a0-e053-d805fe0a8cb0/1963_6243_PhD_Tortella_Pietro.pdf}{PhD thesis, Scuola Internazionale Superiore di Studi Avanzati, 2011.}
	
	\bibitem[Mil25]{J.S.Milne-2025}
	\newblock J.S. Milne.
	\newblock Tannakian Category.
	\newblock \href{https://jmilne.org/math/Books/tcdraft.pdf}{Penultimate version, February 16, 2025.}
	
	\bibitem[DMOS]{Deligne-Milne-Ogus-Shih-1982}
	\newblock P. Deligne, J.S. Milne, A.Ogus and K.Shih.
	\newblock Hodge Cycles, Motives, and Shimura Varieties.
	\newblock \href{https://link.springer.com/book/10.1007/978-3-540-38955-2}{Lect. Notes in Math, 900 (1982), Heidelberg, Springer.}
	
	\bibitem[No76]{M.V.Nori-1976}
	\newblock M.V.Nori.
	\newblock On the representation of the fundamental group.
	\newblock \href{https://eudml.org/doc/89294}{Compositio Math, 33 (1976), 29-41.}
	
	\bibitem[Be83]{J.Bernstein-1983}
	\newblock J. Bernstein.
	\newblock Course on $\mathscr{D}$-modules.
	\newblock \href{https://www.math.columbia.edu/~khovanov/resources/Bernstein-dmod.pdf}{Harvard, 1983-1984.}
	
	\bibitem[Ha77]{R. Hartshorne-1977}
	\newblock Robin Hartshorne.
	\newblock Algebraic Geometry.
	\newblock \href{https://link.springer.com/book/10.1007/978-1-4757-3849-0}{Graduate Texts in Mathematics, vol. 52, Springer-Verlag, New York, 1977.}
	
	\bibitem[Mu65]{D.Mumford-1965}
	\newblock D. Mumford, John Fogarty and Frances Kirwan.
	\newblock Geometric Invariant theory.
	\newblock \href{https://link.springer.com/book/9783642634000}{Springer Verlag, New York (1965).}
	
	\bibitem[HV15]{T.Hausel-F. Villegas}
	\newblock Tam$\acute{a}$s Hausel, F. Rodriguez-Villegas.
	\newblock Cohomology of Large semiprojective hyperkahler varieties.
	\newblock \href{https://hausel.pages.ist.ac.at/wp-content/uploads/sites/229/2024/08/large.pdf}{Ast$\acute{e}$risque 2015 (370) (2015) 113 - 156.}
	
	\bibitem[CM04]{Cortes-Martinez-2004}
	\newblock Jorge Cortes, Eduardo Martinez. 
	\newblock Mechanical Control Systems on Lie algebroids.
	\newblock \href{https://academic.oup.com/imamci/article-abstract/21/4/457/723154?redirectedFrom=fulltext}{IMA J. Math.Control Inf. 21(4) (2004) 457-492.}
	
	\bibitem[LM12]{Lazzarini-Masson-2012}
	\newblock Serge Lazzarini, Thierry Masson.
	\newblock Connections on Lie algebroids and on derivation-based noncommutative geometry.
	\newblock  \href{https://www.sciencedirect.com/science/article/pii/S0393044011002506}{J. Geom. Phys. 62 (2) (2012)
	387–402.}
	
	\bibitem[Hi87]{N.J.Hitchin-1987}
	\newblock N.J. Hitchin.
	\newblock The self-duality equations on a Riemann surface.
	\newblock \href{https://londmathsoc.onlinelibrary.wiley.com/doi/abs/10.1112/plms/s3-55.1.59}{Proc. Lond. Math. Soc. (3) 55 (1) (1987) 59–126.}
	
	\bibitem[Ni93]{N.Nitsure-1993}
	\newblock N. Nitsure.
	\newblock  Moduli of semistable logarithmic connections.
	\newblock \href{https://www.scopus.com/pages/publications/84968466784}{J. Am. Math. Soc. 6 (3) (1993) 597–609.}
	
	
	\bibitem[Ni91]{N.Nitsure-1991}
	\newblock N. Nitsure.
	\newblock Moduli space of semistable pairs on a curve.
	\newblock  \href{https://londmathsoc.onlinelibrary.wiley.com/doi/abs/10.1112/plms/s3-62.2.275}{Proc. Lond. Math. Soc. 62 (3) (1991) 275–300.}
	
	
	\bibitem[De70]{P.Deligne-1970}
	\newblock  P. Deligne.
	\newblock $\acute{E}$quations diff$\acute{e}$rentielles $\grave{a}$ points singuliers r$\acute{e}$guliers.
	\newblock \href{http://dx.doi.org/10.1007/BFb0061194}{Lecture Notes in Mathematics, vol. 163, Springer-Verlag, Berlin-New York, 1970.}
	
	
	
	\bibitem[Bo02]{P.Boalch-2002}
	\newblock Philip Boalch.
	\newblock Quasi-Hamiltonian geometry of meromorphic connections.
	\newblock \href{https://doi.org/10.1215/S0012-7094-07-13924-3}{Duke Math. J. 139 (2002) 369–405.}
	
	\bibitem[BS12]{Bremer-Sage-2012}
	\newblock Christopher L. Bremer, Daniel S. Sage.
	\newblock Moduli spaces of irregular singular connections.
	\newblock \href{https://doi.org/10.1093/imrn/rns103}{Int. Math. Res. Not. 2013 (8) (2012) 1800–1872.}

    \bibitem[Li1]{K.Libor-2009}
	\newblock Krizka Libor.
	\newblock Moduli spaces of Lie algebroid connections.
	\newblock \href{https://eudml.org/doc/250509}{\em Arch. Math. 44 (5) (2009) 403–418.}
	
	\bibitem[Li2]{K.Libor-2010}
	\newblock Krizka Libor.
	\newblock Moduli spaces of flat Lie algebroid connections.
	\newblock \href{http://eudml.org/doc/250509}{arXiv:1012.3180, 2010.}
	
	\bibitem[RS24]{Roy-Singh-2024}
	\newblock Sumit Roy, Anoop Singh.
	\newblock Semiprojectivity of the moduli of principal $G$--bundles with $\lambda$--connections.
	\newblock \href{https://doi.org/10.48550/arXiv.2310.15652}{arxiv.org/abs/2310.15652v2}
	
	\bibitem[Bor87]{A.Borel-1987}
	\newblock A. Borel (et al.).
	\newblock Algebraic D-modules (Perspectives in Mathematics; Vol. 2).
	\newblock \href{https://openlibrary.org/works/OL18822165W?edition=algebraicdmodule00bore}{\em Academic Press, 1987}
	
	\bibitem[GLP18]{GLP-2018}
	\newblock Marco Gualtieri, Songhao Li, and Brent Pym.
	\newblock The Stokes groupoids.
	\newblock \href{https://ora.ox.ac.uk/objects/uuid:e0cc04c7-838d-48dc-a954-68482a8d582f}{\em Journal für die reine und angewandte Mathematik (Crelle's Journal), 2018.}
	%\bibitem[AAB02]{HN-reduction-Anchouche-Azad-Biswas-2002}
	%Boudjemaa Anchouche, Hassan Azad, and Indranil Biswas.
	%\newblock Harder-{N}arasimhan reduction for principal bundles over a compact
	%{K}\"ahler manifold.
	%\newblock {\em Math. Ann.}, 323(4):693--712, 2002.
	%\newblock \doi{10.1007/s002080200322}.
	
	%\bibitem[AB01]{Anchouche-Biswas-2001}
	%Boudjemaa Anchouche and Indranil Biswas.
	%\newblock Einstein-{H}ermitian connections on polystable principal bundles over
	%a compact {K}\"ahler manifold.
	%\newblock {\em Amer. J. Math.}, 123(2):207--228, 2001.
	%\newblock
	%\urlprefix\url{http://muse.jhu.edu/journals/american_journal_of_mathematics/v123/123.2anchouche.pdf}.
	
	\bibitem[AB02]{Azad-Biswas-2002}
	Hassan Azad and Indranil Biswas.
	\newblock On holomorphic principal bundles over a compact {R}iemann surface
	admitting a flat connection.
	\newblock {\em Math. Ann.}, 322(2):333--346, 2002.
	\newblock \doi{10.1007/s002080100273}.
	
	\bibitem[AO24]{Alfaya-Oliveire-2024}
	\newblock David Alfaya and Andr\'{e} Oliveira.
	\newblock Lie algebroid connections, twisted {H}iggs bundles and motives of
	moduli spaces.
	\newblock {\em J. Geom. Phys.}, 201:Paper No. 105195, 55, 2024.
	\newblock \doi{10.1016/j.geomphys.2024.105195}.
	
	\bibitem[Ati57]{Atiyah-1957}
	\newblock M.~F. Atiyah.
	\newblock Complex analytic connections in fibre bundles.
	\newblock {\em Trans. Amer. Math. Soc.}, 85:181--207, 1957.
	\newblock \doi{10.2307/1992969}.
	
	\bibitem[ABKA25]{ABKA-2025}
	\newblock David Alfaya,Indranil Biswas, Pradip Kumar and 
	Anoop Singh.
	\newblock A criterion for holomorphic Lie algebroid
	connections.
	\newblock \href{https://www.sciencedirect.com/science/article/pii/S0021869325003230}{\em Journal of Algebra, Volume 681, 2025, Pages 343-366,ISSN 0021-8693}.
	
	\bibitem[Bi26]{Biswas-2026}
	\newblock Indranil Biswas.
	\newblock Holomorphic Lie Algebroid Connections on
	Holomorphic Principal Bundles on Compact Riemann Surfaces.
	\newblock \href{https://www.ms.u-tokyo.ac.jp/journal/33-1-4.pdf}{\em J. Math. Sci. Univ. Tokyo, 33 (2026), 97–125}.
	
	%\bibitem[PO17]{PO-2017}
	%\newblock Oscar Garcia-Prada and Andr$\acute{e}$ Oliveira.
	%\newblock Connectedness Of Higgs Bundle Moduli For Complex
	%Reductive Lie Groups.
	%\newblock \href{https://archive.intlpress.com/site/pub/files/_fulltext/journals/ajm/2017/0021/0005/AJM-2017-0021-0005-a001.pdf}{\em ASIAN J. MATH. Vol. 21, No. 5, pp. 791–810, October 2017}.
	
	%\bibitem[Pi25]{Piotr-2025}
	%\newblock Piotr Achinger.
	%\newblock Fundamental groups in algebraic geometry.
	%\newblock \href{https://achinger.impan.pl/pi1/notes.pdf}{\em Lecture notes}
	
	%\bibitem[Co98]{Cohen-1998}
	%\newblock Ralph L. Cohen.
	%\newblock The Topology of Fiber Bundles.
	%\newblock \href{https://math.stanford.edu/~ralph/fiber.pdf}{\em Lecture Notes}
	
	%\bibitem[No51]{Norman-1951}
	%\newblock Norman Steenrod.
	%\newblock The Topology of Fibre Bundles (PMS-14).
	%\newblock \href{http://www.jstor.org/stable/j.ctt1bpm9t5.}{\em  Princeton University Press, 1951.}
	
	%\bibitem[BPG08]{BPG-2008}
	%\newblock Steven B. Bradlow, Oscar García-Prada and Peter B. Gothen.
	%\newblock Homotopy groups of moduli spaces of representations.
	%\newblock \href{https://www.sciencedirect.com/science/article/pii/S0040938307000602}{\em Topology, Volume 47, Issue 4,2008,Pages 203-224,
		%ISSN 0040-9383,}
		
	%\bibitem[KKT23]{KKT-2023}
	%\newblock Masaki Kameko, Daisuke Kishimoto and Masahiro Takeda.
	%\newblock Homotopy types of gauge groups over Riemann surfaces.
	%\newblock \href{https://msp.org/agt/2023/23-5/agt-v23-n5-p12-p.pdf}{\em Algebraic \& Geometric Topology 23:5 (2023) 2309–2327}
	
	%\bibitem[Si22]{Simon-2022}
	%\newblock Simon Rea.
	%\newblock Homotopy types of gauge groups of
	%principal bundles with certain non-simply
	%connected structure groups.
	%\newblock \href{https://eprints.soton.ac.uk/467266/1/Simon_Rea_Thesis.pdf}{PhD Thesis.}
	
	%\bibitem[AB83]{AB-1983}
	%\newblock  M. F. Atiyah and R. Bott.
	%\newblock  The Yang–Mills equations over Riemann surfaces.
	%\newblock \href{https://www.jstor.org/stable/pdf/37156.pdf}{\em Philos. Trans. Roy. Soc. London Ser. A 308 (1983) 523–615.}
	
	%\bibitem[DU95]{DU-1995}
	%\newblock G. D. Daskalopoulos and K. K. Uhlenbeck.
	%\newblock An application of transversality to the
	%topology of the moduli space of stable bundles.
	%\newblock \href{https://www.sciencedirect.com/science/article/pii/0040938394E0014B}{\em Topology, Volume 34, Issue 1, 1995, Pages 203-215,
		%ISSN 0040-9383}
		
	%\bibitem[Th13]{Th-2013}
	%\newblock S. D. Theriault.
	%\newblock The homotopy types of gauge groups of nonorientable surfaces and applications to moduli spaces.
	%\newblock \href{https://projecteuclid.org/journals/illinois-journal-of-mathematics/volume-57/issue-1/The-homotopy-types-of-gauge-groups-of-nonorientable-surfaces-and/10.1215/ijm/1403534486.pdf}{\em  Illinois J. Math. 57 (2013), no. 1, 59–85.
	%}
	
	%\bibitem[Th11]{Th-2011}
	%\newblock S. D. Theriault.
	%\newblock Homotopy decompositions of gauge groups over riemann surfaces and applications to moduli spaces.
	%\newblock \href{https://doi.org/10.1142/S0129167X11007690}{\em  Int. J. Math. 22 (2011), no. 12, 1711–1719.}
	
	
	%\bibitem[BDP18]{Biswas-Dan-Paul-2018}
	%Indranil Biswas, Ananyo Dan, and Arjun Paul.
	%\newblock Criterion for logarithmic connections with prescribed residues.
	%\newblock {\em Manuscripta Math.}, 155(1-2):77--88, 2018.
	%\newblock \doi{10.1007/s00229-017-0935-6}.
	
	%\bibitem[BDPS17]{Biswas-Dan-Paul-Saha-2017}
	%Indranil Biswas, Ananyo Dan, Arjun Paul, and Arideep Saha.
	%\newblock Logarithmic connections on principal bundles over a {R}iemann
	%surface.
	%\newblock {\em Internat. J. Math.}, 28(12):1750088, 18, 2017.
	%\newblock \doi{10.1142/S0129167X17500884}.
	
	%\bibitem[BKS24]{Biswas-Kumar-Singh-2024}
	%Indranil Biswas, Pradip Kumar, and Anoop Singh.
	%\newblock A criterion for {L}ie algebroid connections on a compact {R}iemann
	%surface.
	%\newblock {\em Geom. Dedicata}, 218(4):Paper No. 87, 2024.
	%\newblock \doi{10.1007/s10711-024-00938-8}.
	
	%\bibitem[BP17]{Biswas-Paul-2017}
	%Indranil Biswas and Arjun Paul.
	%\newblock Equivariant bundles and connections.
	%\newblock {\em Ann. Global Anal. Geom.}, 51(4):347--358, 2017.
	%\newblock \doi{10.1007/s10455-016-9538-9}.
	
	%\bibitem[GP20]{Gurjar-Paul-2020}
	%Sudarshan Gurjar and Arjun Paul.
	%\newblock Criterion for existence of a logarithmic connection on a principal
	%bundle over a smooth complex projective variety.
	%\newblock {\em Ann. Global Anal. Geom.}, 58(3):241--251, 2020.
	%\newblock \doi{10.1007/s10455-020-09723-8}.
	
	\bibitem[HL10]{Huybrechts-Lehn-2010}
	Daniel Huybrechts and Manfred Lehn.
	\newblock {\em The geometry of moduli spaces of sheaves}.
	\newblock Cambridge Mathematical Library. Cambridge University Press,
	Cambridge, second edition, 2010.
	\newblock \doi{10.1017/CBO9780511711985}.
	
	%\bibitem[Ram75]{Ramanathan-1975}
	%A.~Ramanathan.
	%\newblock Stable principal bundles on a compact {R}iemann surface.
	%\newblock {\em Math. Ann.}, 213:129--152, 1975.
	%\newblock \doi{10.1007/BF01343949}.
	
	%\bibitem[Ram96]{Ramanathan-I}
	%A.~Ramanathan.
	%\newblock Moduli for principal bundles over algebraic curves. {I}.
	%\newblock {\em Proc. Indian Acad. Sci. Math. Sci.}, 106(3):301--328, 1996.
	%\newblock \doi{10.1007/BF02867438}.
	
	%\bibitem[Wei38]{Weil-1938}
	%Andre~A. Weil.
	%\newblock Généralisation des fonctions abéliennes.
	%\newblock {\em J. Math. Pures Appl.}, 17:47--87, 1938.
	
	%\bibitem[FGN19]{FGN-2019}
	%\newblock C. Florentino,P.B. Gothen and A. Nozad.
	%\newblock Homotopy type of moduli spaces of G-Higgs bundles and reducibility of the nilpotent cone.
	%\newblock \href{https://www.sciencedirect.com/science/article/pii/S0007449718300927}{\em Bulletin des Sciences Mathématiques,Volume 150,2019,Pages 84-101,ISSN 0007-4497}
	
	%\bibitem[PM20]{PM-2020}
	%\newblock Zahra Pirbodaghi and Morteza Mirmohammad Rezaii.
	%\newblock Forms and Chern Classes on Hermitian Lie Algebroids.
	%\newblock \href{https://doi.org/10.1007/s41980-019-00238-y}{\em Bull. Iran. Math. Soc. 46, 19–36 (2020)}
	
	%\bibitem[Fa93]{Faltings-1993}
	%\newblock Gerd Faltings.
	%\newblock Stable--bundles and projective connections.
	%\newblock {\em Journal of Algebraic Geometry, Vol. 2, No. 3, 507–568, 1993.}
	
	%\bibitem[BH12]{BH-2012}
	%\newblock I.Biswas and N.Hoffmann.
	%\newblock Poincar$\acute{e}$ families of $G$--bundles on curve.
	%\newblock \href{https://arxiv.org/abs/1001.2123}{\em Math.Ann.352 (2012)}
	
	\bibitem[Sch08]{Sch-2008}
	\newblock Alexander H. W. Schmitt.
	\newblock Geometric Invariant Theory and Decorated Principal Bundles.
	\newblock \href{https://ems.press/books/zlam/224}{\em European Mathematical Society Publishing House, Zürich. ISBN 978-3-03719-065-4.}
	
	\bibitem[GL21]{GL-2021}
	\newblock Marco Gualtieri, Kevin Luk. 
	\newblock Log Picard Algebroids and Meromorphic Line Bundles.
	\newblock \href{https://doi.org/10.1093/imrn/rnz272}{International Mathematics Research Notices, Volume 2021, Issue 21, November 2021, Pages 16592–16635.}
	
	\bibitem[GSX09]{GSX-2009}
	\newblock Camille Laurent-Gengoux, Mathieu Stienon and Ping Xu.
	\newblock Integration of holomorphic Lie algebroids.
	\newblock \href{https://link.springer.com/article/10.1007/s00208-009-0388-7}{\em Math. Ann. 345, 895–923 (2009)}
	
	\bibitem[MM03]{Moerdijk-2003}
	\newblock Moerdijk I, Mrcun J.
	\newblock Introduction to Foliations and Lie Groupoids.
	\newblock \href{https://www.cambridge.org/core/books/introduction-to-foliations-and-lie-groupoids/75BBFED277FF39A56594731202789016}{\em Cambridge University Press; 2003.}
	
	\bibitem[BB93]{Beilinson-Bernstein-1993}
	\newblock A. Beilinson and J. Bernstein.
	\newblock A proof of Jantzen conjectures.
	\newblock \href{https://www.math.tau.ac.il/~bernstei/Publication_list/publication_texts/Ber-Beil_ProofJatzenConj.pdf}{Advances in Soviet Mathematics, Vol. 16, Part 1, pp. 1–50, 1993}.
	
	\bibitem[FT17]{FT-2017}
	\newblock E. Franco and P. Tortella.
	\newblock  Moduli spaces of \(\Lambda \)-modules on abelian varieties.
	\newblock \href{https://www.sciencedirect.com/science/article/pii/S0001870816305849}{\em Advances in Mathematics, 318, 459-496, 2017,.}
	
	%\bibitem[To17]{Tortella-2017}
	%\newblock Pietro Tortella.
	%\newblock Representations of Atiyah Algebroids and Logarithmic Connections.
	%\newblock \href{https://doi.org/10.1093/imrn/rnw010}{\em Volume 2017, Issue 1, January 2017, Pages 29–46.}
	
	\bibitem[Mor07]{A. Moroianu-2007}
	\newblock A. Moroianu.
	\newblock  Lectures on Kähler Geometry (London Mathematical Society Student Texts).
	\newblock \href{https://www.cambridge.org/core/books/abs/lectures-on-kahler-geometry/contents/823EB6F94327152106F30C3046D6DBF4}{\em Cambridge University Press, 2007:v-viii.}
	
	\bibitem[Fer02]{R.L.Fernandez-2002}
	\newblock Rui Loja Fernandes.
	\newblock Lie Algebroids, Holonomy and Characteristic Classes.
	\newblock \href{https://www.sciencedirect.com/science/article/pii/S0001870801920705}{\em Advances in Mathematics,
	Volume 170, Issue 1,
	2002,Pages 119-179, ISSN 0001-8708.}
		
	\bibitem[CF09]{Crainic-2009}
	\newblock M. Crainic, R.L. Fernandes.
	\newblock Exotic characteristic classes of Lie algebroids.
	\newblock \href{https://www.math.tecnico.ulisboa.pt/~rfern/Meus-papers/HTML/jets.pdf}{\em Quantum Field Theory and Noncommutative Geometry, Lect. Notes, 2009.}
	
	\bibitem[HH22]{Hausel-Hitchin-2022}
	\newblock T. Hausel, N. Hitchin.
	\newblock Very stable Higgs bundles, equivariant multiplicity and mirror symmetry.
	\newblock \href{https://doi.org/10.1007/s00222-021-01093-7}{\em Invent. math. 228, 893–989 (2022).}
	
	\bibitem[HR08]{Hausel-Rodriguez-2008}
	\newblock T. Hausel, F. Rodriguez-Villegas.
	\newblock Mixed Hodge polynomials of character varieties.
	\newblock \href{https://doi.org/10.1007/s00222-008-0142-x}{\em Invent. math. 174, 555–624 (2008)}.
	
	\bibitem[GPN24]{Gallego-Oscar-Narasimhan-2024}
	\newblock Guillermo Gallego, Oscar García-Prada, and M. S. Narasimhan.
	\newblock Higgs bundles twisted by a vector bundle.
	\newblock \href{https://doi.org/10.1142/S0129167X24410076}{\em International Journal of Mathematics, Vol. 35, No. 09, 2441007 (2024)}.
	
	\bibitem[HL21]{Hoskins-Lehalleur-2021}
	\newblock V. Hoskins, S.P. Lehalleur.
	\newblock On the Voevodsky motive of the moduli space of Higgs bundles on a curve. 
	\newblock \href{https://doi.org/10.1007/s00029-020-00610-5}{\em Sel. Math. New Ser. 27, 11 (2021).}
	
\end{thebibliography}




























\end{document}
